% ============================================================================
%  Supporting Information for
%  "Activity cliffs on human acetylcholinesterase: a pre-specified escalation
%   from 2D descriptors to molecular dynamics"
%
%  ACS requires the supporting information as its own file, so this is a
%  standalone document that shares refs.bib and ../results with the main text.
%
%  Build:  latexmk -pdf si.tex
% ============================================================================
\documentclass[letterpaper]{article}

\usepackage[T1]{fontenc}
\usepackage[utf8]{inputenc}
\usepackage{geometry}
\geometry{margin = 1in}
\usepackage{achemso}
\usepackage{graphicx}
\usepackage{array}
\usepackage{booktabs}
\usepackage{amsmath}
\usepackage{siunitx}
\usepackage{longtable}
\usepackage[colorlinks=true,allcolors=blue]{hyperref}

\graphicspath{{../figures/}}
\setcounter{secnumdepth}{-1}

\newcommand{\framingA}{framing~A}
\newcommand{\framingB}{framing~B}
\newcommand{\framingC}{framing~C}

\renewcommand{\thetable}{S\arabic{table}}
\renewcommand{\thefigure}{S\arabic{figure}}

\title{Supporting Information\\
  Activity cliffs on human acetylcholinesterase:
  a pre-specified descriptor escalation}
\author{Isman Kurniawan, Lela Lailatul Khumaisah, Muhamad Salman Fareza}
\date{}

\begin{document}

\maketitle

Table~\ref{tab:sistructures} lists every molecular structure the study used. The
journal's guidance on QSAR submissions asks for these to be reported here rather than
only deposited. Each ligand appears with both of the forms it was carried in, with its
measured potency, and with the split and fold it was frozen into before any descriptor
was computed.

\footnotesize
\begin{longtable}{l p{4.4cm} p{4.4cm} r r l l}
\caption{\textbf{Every ligand used in this study.} The standardised neutral
  structure carries the fingerprints and the cliff definition; the protonated
  structure enters the docking and the simulations. Potency is pChEMBL, the
  negative logarithm of the molar IC\textsubscript{50}. The split and fold were frozen
  before any descriptor was computed.}\label{tab:sistructures}\\
\toprule
ChEMBL ID & SMILES, neutral & SMILES, pH 7.4 & pChEMBL & $q$ & Split & Group \\
\midrule
\endfirsthead
\multicolumn{7}{l}{\itshape Table \thetable{} continued}\\
\toprule
ChEMBL ID & SMILES, neutral & SMILES, pH 7.4 & pChEMBL & $q$ & Split & Group \\
\midrule
\endhead
\midrule \multicolumn{7}{r}{\itshape continued on the next page}\\
\endfoot
\bottomrule
\endlastfoot
CHEMBL104387 & \texttt{c1ccc2c(\allowbreak{}c1)\allowbreak{}ccc[n+]\allowbreak{}2CCCCCCCCCCCC[n+]\allowbreak{}1cccc2ccccc21} & \texttt{c1ccc2c(\allowbreak{}c1)\allowbreak{}ccc[n+]\allowbreak{}2CCCCCCCCCCCC[n+]\allowbreak{}1cccc2ccccc21} & 7.70 & +2 & test & cliff \\
CHEMBL104553 & \texttt{c1ccc2c[n+]\allowbreak{}(\allowbreak{}CCCCCCCCC[n+]\allowbreak{}3ccc4ccccc4c3)\allowbreak{}ccc2c1} & \texttt{c1ccc2c[n+]\allowbreak{}(\allowbreak{}CCCCCCCCC[n+]\allowbreak{}3ccc4ccccc4c3)\allowbreak{}ccc2c1} & 7.40 & +2 & fold 1 & cliff \\
CHEMBL104613 & \texttt{c1ccc2c[n+]\allowbreak{}(\allowbreak{}CCCCCCCCCCC[n+]\allowbreak{}3ccc4ccccc4c3)\allowbreak{}ccc2c1} & \texttt{c1ccc2c[n+]\allowbreak{}(\allowbreak{}CCCCCCCCCCC[n+]\allowbreak{}3ccc4ccccc4c3)\allowbreak{}ccc2c1} & 7.00 & +2 & fold 1 & cliff \\
CHEMBL105017 & \texttt{c1ccc2c(\allowbreak{}c1)\allowbreak{}ccc[n+]\allowbreak{}2CCCCCCCCC[n+]\allowbreak{}1cccc2ccccc21} & \texttt{c1ccc2c(\allowbreak{}c1)\allowbreak{}ccc[n+]\allowbreak{}2CCCCCCCCC[n+]\allowbreak{}1cccc2ccccc21} & 7.10 & +2 & test & cliff \\
CHEMBL105257 & \texttt{c1ccc2c(\allowbreak{}c1)\allowbreak{}ccc[n+]\allowbreak{}2CCCCCCCCCCC[n+]\allowbreak{}1cccc2ccccc21} & \texttt{c1ccc2c(\allowbreak{}c1)\allowbreak{}ccc[n+]\allowbreak{}2CCCCCCCCCCC[n+]\allowbreak{}1cccc2ccccc21} & 8.00 & +2 & test & cliff \\
CHEMBL1083790 & \texttt{NCCCCCCCCCNc1c2c(\allowbreak{}nc3ccccc13)\allowbreak{}CCCC2} & \texttt{[NH3+]\allowbreak{}CCCCCCCCCNc1c2c(\allowbreak{}[nH+]\allowbreak{}c3ccccc13)\allowbreak{}CCCC2} & 6.88 & +2 & fold 4 & control \\
CHEMBL1083791 & \texttt{CN(\allowbreak{}C)\allowbreak{}CCCCCCCNc1c2c(\allowbreak{}nc3ccccc13)\allowbreak{}CCCC2} & \texttt{C[NH+]\allowbreak{}(\allowbreak{}C)\allowbreak{}CCCCCCCNc1c2c(\allowbreak{}[nH+]\allowbreak{}c3ccccc13)\allowbreak{}CCCC2} & 7.28 & +2 & fold 2 & control \\
CHEMBL1083972 & \texttt{COC(\allowbreak{}=\allowbreak{}O)\allowbreak{}c1cc2c(\allowbreak{}N)\allowbreak{}c3c(\allowbreak{}nc2nc1C)\allowbreak{}CCC3} & \texttt{COC(\allowbreak{}=\allowbreak{}O)\allowbreak{}c1cc2c(\allowbreak{}N)\allowbreak{}c3c(\allowbreak{}[nH+]\allowbreak{}c2[nH+]\allowbreak{}c1C)\allowbreak{}CCC3} & 5.72 & +2 & fold 1 & control \\
CHEMBL1083973 & \texttt{CCOC(\allowbreak{}=\allowbreak{}O)\allowbreak{}c1cc2c(\allowbreak{}N)\allowbreak{}c3c(\allowbreak{}nc2nc1C)\allowbreak{}CCCCC3} & \texttt{CCOC(\allowbreak{}=\allowbreak{}O)\allowbreak{}c1cc2c(\allowbreak{}N)\allowbreak{}c3c(\allowbreak{}[nH+]\allowbreak{}c2[nH+]\allowbreak{}c1C)\allowbreak{}CCCCC3} & 5.70 & +2 & test & control \\
CHEMBL1083983 & \texttt{O=\allowbreak{}C(\allowbreak{}CCc1c[nH]\allowbreak{}c2ccccc12)\allowbreak{}NCCCCCCNc1c2c(\allowbreak{}nc3ccccc13)\allowbreak{}CCCC2} & \texttt{O=\allowbreak{}C(\allowbreak{}CCc1c[nH]\allowbreak{}c2ccccc12)\allowbreak{}NCCCCCCNc1c2c(\allowbreak{}[nH+]\allowbreak{}c3ccccc13)\allowbreak{}CCCC2} & 10.70 & +1 & fold 1 & cliff \\
CHEMBL1083984 & \texttt{O=\allowbreak{}C(\allowbreak{}CCc1c[nH]\allowbreak{}c2ccccc12)\allowbreak{}NCCCCCCCNc1c2c(\allowbreak{}nc3ccccc13)\allowbreak{}CCCC2} & \texttt{O=\allowbreak{}C(\allowbreak{}CCc1c[nH]\allowbreak{}c2ccccc12)\allowbreak{}NCCCCCCCNc1c2c(\allowbreak{}[nH+]\allowbreak{}c3ccccc13)\allowbreak{}CCCC2} & 10.22 & +1 & fold 1 & cliff \\
CHEMBL1084092 & \texttt{NCCCCCCCCCCNc1c2c(\allowbreak{}nc3ccccc13)\allowbreak{}CCCC2} & \texttt{[NH3+]\allowbreak{}CCCCCCCCCCNc1c2c(\allowbreak{}[nH+]\allowbreak{}c3ccccc13)\allowbreak{}CCCC2} & 6.83 & +2 & fold 3 & control \\
CHEMBL1084093 & \texttt{NCCCCCCCCCCCCNc1c2c(\allowbreak{}nc3ccccc13)\allowbreak{}CCCC2} & \texttt{[NH3+]\allowbreak{}CCCCCCCCCCCCNc1c2c(\allowbreak{}[nH+]\allowbreak{}c3ccccc13)\allowbreak{}CCCC2} & 6.75 & +2 & fold 4 & control \\
CHEMBL1084095 & \texttt{CN(\allowbreak{}C)\allowbreak{}CCCCCCCCCNc1c2c(\allowbreak{}nc3ccccc13)\allowbreak{}CCCC2} & \texttt{C[NH+]\allowbreak{}(\allowbreak{}C)\allowbreak{}CCCCCCCCCNc1c2c(\allowbreak{}[nH+]\allowbreak{}c3ccccc13)\allowbreak{}CCCC2} & 7.15 & +2 & fold 2 & control \\
CHEMBL1084096 & \texttt{c1ccc2c(\allowbreak{}NCCCCCCCCNc3ccncc3)\allowbreak{}c3c(\allowbreak{}nc2c1)\allowbreak{}CCCC3} & \texttt{c1ccc2c(\allowbreak{}NCCCCCCCCNc3cc[nH+]\allowbreak{}cc3)\allowbreak{}c3c(\allowbreak{}[nH+]\allowbreak{}c2c1)\allowbreak{}CCCC3} & 7.87 & +2 & test & control \\
CHEMBL1084262 & \texttt{O=\allowbreak{}C(\allowbreak{}CCc1c[nH]\allowbreak{}c2ccccc12)\allowbreak{}NCCCCCCCCCNc1c2c(\allowbreak{}nc3cc(\allowbreak{}Cl)\allowbreak{}ccc13)\allowbreak{}CCCC2} & \texttt{O=\allowbreak{}C(\allowbreak{}CCc1c[nH]\allowbreak{}c2ccccc12)\allowbreak{}NCCCCCCCCCNc1c2c(\allowbreak{}[nH+]\allowbreak{}c3cc(\allowbreak{}Cl)\allowbreak{}ccc13)\allowbreak{}CCCC2} & 8.36 & +1 & fold 4 & cliff \\
CHEMBL1084263 & \texttt{O=\allowbreak{}C(\allowbreak{}CCc1c[nH]\allowbreak{}c2ccccc12)\allowbreak{}NCCCCCCCCCCNc1c2c(\allowbreak{}nc3cc(\allowbreak{}Cl)\allowbreak{}ccc13)\allowbreak{}CCCC2} & \texttt{O=\allowbreak{}C(\allowbreak{}CCc1c[nH]\allowbreak{}c2ccccc12)\allowbreak{}NCCCCCCCCCCNc1c2c(\allowbreak{}[nH+]\allowbreak{}c3cc(\allowbreak{}Cl)\allowbreak{}ccc13)\allowbreak{}CCCC2} & 7.66 & +1 & fold 4 & cliff \\
CHEMBL1084369 & \texttt{c1ccc2c(\allowbreak{}NCCCCCCCNc3ccncc3)\allowbreak{}c3c(\allowbreak{}nc2c1)\allowbreak{}CCCC3} & \texttt{c1ccc2c(\allowbreak{}NCCCCCCCNc3cc[nH+]\allowbreak{}cc3)\allowbreak{}c3c(\allowbreak{}[nH+]\allowbreak{}c2c1)\allowbreak{}CCCC3} & 7.89 & +2 & test & control \\
CHEMBL1084517 & \texttt{COC(\allowbreak{}=\allowbreak{}O)\allowbreak{}c1cc2c(\allowbreak{}N)\allowbreak{}c3c(\allowbreak{}nc2nc1C)\allowbreak{}CCCCCC3} & \texttt{COC(\allowbreak{}=\allowbreak{}O)\allowbreak{}c1cc2c(\allowbreak{}N)\allowbreak{}c3c(\allowbreak{}[nH+]\allowbreak{}c2[nH+]\allowbreak{}c1C)\allowbreak{}CCCCCC3} & 5.75 & +2 & fold 1 & control \\
CHEMBL1084519 & \texttt{CCOC(\allowbreak{}=\allowbreak{}O)\allowbreak{}c1cc2c(\allowbreak{}N)\allowbreak{}c3c(\allowbreak{}nc2nc1C)\allowbreak{}CCCCCC3} & \texttt{CCOC(\allowbreak{}=\allowbreak{}O)\allowbreak{}c1cc2c(\allowbreak{}N)\allowbreak{}c3c(\allowbreak{}[nH+]\allowbreak{}c2[nH+]\allowbreak{}c1C)\allowbreak{}CCCCCC3} & 5.75 & +2 & test & control \\
CHEMBL1085011 & \texttt{CCCCCCNC(\allowbreak{}=\allowbreak{}O)\allowbreak{}Oc1ccc2c(\allowbreak{}c1)\allowbreak{}C1CCN(\allowbreak{}CCC)\allowbreak{}C1C2} & \texttt{CCCCCCNC(\allowbreak{}=\allowbreak{}O)\allowbreak{}Oc1ccc2c(\allowbreak{}c1)\allowbreak{}C1CC[NH+]\allowbreak{}(\allowbreak{}CCC)\allowbreak{}C1C2} & 6.26 & +1 & fold 1 & control \\
CHEMBL1085630 & \texttt{CCCCCCCNC(\allowbreak{}=\allowbreak{}O)\allowbreak{}Oc1ccc2c(\allowbreak{}c1)\allowbreak{}C1CCN(\allowbreak{}CCC)\allowbreak{}C1C2} & \texttt{CCCCCCCNC(\allowbreak{}=\allowbreak{}O)\allowbreak{}Oc1ccc2c(\allowbreak{}c1)\allowbreak{}C1CC[NH+]\allowbreak{}(\allowbreak{}CCC)\allowbreak{}C1C2} & 6.20 & +1 & fold 1 & control \\
CHEMBL1085933 & \texttt{CNC(\allowbreak{}=\allowbreak{}O)\allowbreak{}Oc1cccc(\allowbreak{}CN(\allowbreak{}C)\allowbreak{}CCCCCCCOc2ccc3c(\allowbreak{}c2)\allowbreak{}O/\allowbreak{}C(\allowbreak{}=\allowbreak{}C\textbackslash{}c2ccc4ccccc4c2)\allowbreak{}C3=\allowbreak{}O)\allowbreak{}c1} & \texttt{CNC(\allowbreak{}=\allowbreak{}O)\allowbreak{}Oc1cccc(\allowbreak{}C[NH+]\allowbreak{}(\allowbreak{}C)\allowbreak{}CCCCCCCOc2ccc3c(\allowbreak{}c2)\allowbreak{}O/\allowbreak{}C(\allowbreak{}=\allowbreak{}C\textbackslash{}c2ccc4ccccc4c2)\allowbreak{}C3=\allowbreak{}O)\allowbreak{}c1} & 8.68 & +1 & fold 4 & cliff \\
CHEMBL118968 & \texttt{Nc1c2c(\allowbreak{}nc3ccccc13)\allowbreak{}CCCCCC2} & \texttt{Nc1c2c(\allowbreak{}[nH+]\allowbreak{}c3ccccc13)\allowbreak{}CCCCCC2} & 5.04 & +1 & fold 2 & cliff \\
CHEMBL1257412 & \texttt{CCCCCCCCCCCCCCNc1c2c(\allowbreak{}nc3ccc(\allowbreak{}OC)\allowbreak{}cc13)\allowbreak{}CCCC2} & \texttt{CCCCCCCCCCCCCCNc1c2c(\allowbreak{}[nH+]\allowbreak{}c3ccc(\allowbreak{}OC)\allowbreak{}cc13)\allowbreak{}CCCC2} & 4.34 & +1 & fold 4 & cliff \\
CHEMBL1257413 & \texttt{CCCCCCCCCCCCCNc1c2c(\allowbreak{}nc3ccc(\allowbreak{}OC)\allowbreak{}cc13)\allowbreak{}CCCC2} & \texttt{CCCCCCCCCCCCCNc1c2c(\allowbreak{}[nH+]\allowbreak{}c3ccc(\allowbreak{}OC)\allowbreak{}cc13)\allowbreak{}CCCC2} & 4.00 & +1 & fold 4 & cliff \\
CHEMBL1257525 & \texttt{CCCCCCCCCCCCNc1c2c(\allowbreak{}nc3ccc(\allowbreak{}OC)\allowbreak{}cc13)\allowbreak{}CCCC2} & \texttt{CCCCCCCCCCCCNc1c2c(\allowbreak{}[nH+]\allowbreak{}c3ccc(\allowbreak{}OC)\allowbreak{}cc13)\allowbreak{}CCCC2} & 6.00 & +1 & fold 4 & cliff \\
CHEMBL1257766 & \texttt{CCCCCCCCNc1c2c(\allowbreak{}nc3ccc(\allowbreak{}OC)\allowbreak{}cc13)\allowbreak{}CCCC2} & \texttt{CCCCCCCCNc1c2c(\allowbreak{}[nH+]\allowbreak{}c3ccc(\allowbreak{}OC)\allowbreak{}cc13)\allowbreak{}CCCC2} & 6.44 & +1 & fold 4 & cliff \\
CHEMBL1257767 & \texttt{CCCCCCCNc1c2c(\allowbreak{}nc3ccc(\allowbreak{}OC)\allowbreak{}cc13)\allowbreak{}CCCC2} & \texttt{CCCCCCCNc1c2c(\allowbreak{}[nH+]\allowbreak{}c3ccc(\allowbreak{}OC)\allowbreak{}cc13)\allowbreak{}CCCC2} & 6.57 & +1 & fold 4 & cliff \\
CHEMBL1257885 & \texttt{CCCCCCNc1c2c(\allowbreak{}nc3ccc(\allowbreak{}OC)\allowbreak{}cc13)\allowbreak{}CCCC2} & \texttt{CCCCCCNc1c2c(\allowbreak{}[nH+]\allowbreak{}c3ccc(\allowbreak{}OC)\allowbreak{}cc13)\allowbreak{}CCCC2} & 7.00 & +1 & fold 4 & cliff \\
CHEMBL1257887 & \texttt{CCCCNc1c2c(\allowbreak{}nc3ccc(\allowbreak{}OC)\allowbreak{}cc13)\allowbreak{}CCCC2} & \texttt{CCCCNc1c2c(\allowbreak{}[nH+]\allowbreak{}c3ccc(\allowbreak{}OC)\allowbreak{}cc13)\allowbreak{}CCCC2} & 4.68 & +1 & fold 4 & cliff \\
CHEMBL127465 & \texttt{CN(\allowbreak{}CCCOc1ccc2c(\allowbreak{}=\allowbreak{}O)\allowbreak{}c3ccccc3oc2c1)\allowbreak{}Cc1cccc(\allowbreak{}OC(\allowbreak{}=\allowbreak{}O)\allowbreak{}NCCCCCCCCN2CCOCC2)\allowbreak{}c1} & \texttt{C[NH+]\allowbreak{}(\allowbreak{}CCCOc1ccc2c(\allowbreak{}=\allowbreak{}O)\allowbreak{}c3ccccc3oc2c1)\allowbreak{}Cc1cccc(\allowbreak{}OC(\allowbreak{}=\allowbreak{}O)\allowbreak{}NCCCCCCCC[NH+]\allowbreak{}2CCOCC2)\allowbreak{}c1} & 7.56 & +2 & fold 0 & control \\
CHEMBL127857 & \texttt{CN(\allowbreak{}CCCOc1ccc2c(\allowbreak{}=\allowbreak{}O)\allowbreak{}c3ccccc3oc2c1)\allowbreak{}Cc1cccc(\allowbreak{}OC(\allowbreak{}=\allowbreak{}O)\allowbreak{}NCCCCCCCCCN2CCOCC2)\allowbreak{}c1} & \texttt{C[NH+]\allowbreak{}(\allowbreak{}CCCOc1ccc2c(\allowbreak{}=\allowbreak{}O)\allowbreak{}c3ccccc3oc2c1)\allowbreak{}Cc1cccc(\allowbreak{}OC(\allowbreak{}=\allowbreak{}O)\allowbreak{}NCCCCCCCCC[NH+]\allowbreak{}2CCOCC2)\allowbreak{}c1} & 7.79 & +2 & fold 0 & control \\
CHEMBL129577 & \texttt{CCCCCCCNC(\allowbreak{}=\allowbreak{}O)\allowbreak{}Oc1cccc(\allowbreak{}CN(\allowbreak{}CCCOc2ccc3c(\allowbreak{}=\allowbreak{}O)\allowbreak{}c4ccccc4oc3c2)\allowbreak{}C(\allowbreak{}C)\allowbreak{}C)\allowbreak{}c1} & \texttt{CCCCCCCNC(\allowbreak{}=\allowbreak{}O)\allowbreak{}Oc1cccc(\allowbreak{}C[NH+]\allowbreak{}(\allowbreak{}CCCOc2ccc3c(\allowbreak{}=\allowbreak{}O)\allowbreak{}c4ccccc4oc3c2)\allowbreak{}C(\allowbreak{}C)\allowbreak{}C)\allowbreak{}c1} & 6.77 & +1 & fold 3 & control \\
CHEMBL129753 & \texttt{CNC(\allowbreak{}=\allowbreak{}O)\allowbreak{}Oc1cccc(\allowbreak{}CN(\allowbreak{}C)\allowbreak{}CCCCOc2ccc3c(\allowbreak{}=\allowbreak{}O)\allowbreak{}c4ccccc4oc3c2)\allowbreak{}c1} & \texttt{CNC(\allowbreak{}=\allowbreak{}O)\allowbreak{}Oc1cccc(\allowbreak{}C[NH+]\allowbreak{}(\allowbreak{}C)\allowbreak{}CCCCOc2ccc3c(\allowbreak{}=\allowbreak{}O)\allowbreak{}c4ccccc4oc3c2)\allowbreak{}c1} & 7.43 & +1 & fold 3 & cliff \\
CHEMBL131411 & \texttt{CCCCNC(\allowbreak{}=\allowbreak{}O)\allowbreak{}Oc1cccc(\allowbreak{}CN(\allowbreak{}CCCOc2ccc3c(\allowbreak{}=\allowbreak{}O)\allowbreak{}c4ccccc4oc3c2)\allowbreak{}C(\allowbreak{}C)\allowbreak{}C)\allowbreak{}c1} & \texttt{CCCCNC(\allowbreak{}=\allowbreak{}O)\allowbreak{}Oc1cccc(\allowbreak{}C[NH+]\allowbreak{}(\allowbreak{}CCCOc2ccc3c(\allowbreak{}=\allowbreak{}O)\allowbreak{}c4ccccc4oc3c2)\allowbreak{}C(\allowbreak{}C)\allowbreak{}C)\allowbreak{}c1} & 7.00 & +1 & fold 3 & control \\
CHEMBL13166 & \texttt{CN(\allowbreak{}C)\allowbreak{}Cc1ccc(\allowbreak{}CSCCCCCSCc2ccc(\allowbreak{}CN(\allowbreak{}C)\allowbreak{}C)\allowbreak{}o2)\allowbreak{}o1} & \texttt{C[NH+]\allowbreak{}(\allowbreak{}C)\allowbreak{}Cc1ccc(\allowbreak{}CSCCCCCSCc2ccc(\allowbreak{}C[NH+]\allowbreak{}(\allowbreak{}C)\allowbreak{}C)\allowbreak{}o2)\allowbreak{}o1} & 6.05 & +2 & fold 0 & control \\
CHEMBL162734 & \texttt{c1ccc2c(\allowbreak{}c1)\allowbreak{}CN1CCCCCCC1=\allowbreak{}N2} & \texttt{c1ccc2c(\allowbreak{}c1)\allowbreak{}CN1CCCCCCC1=\allowbreak{}[NH+]\allowbreak{}2} & 5.04 & +1 & fold 0 & control \\
CHEMBL163362 & \texttt{c1ccc2c(\allowbreak{}c1)\allowbreak{}CN1CCCCC1=\allowbreak{}N2} & \texttt{c1ccc2c(\allowbreak{}c1)\allowbreak{}CN1CCCCC1=\allowbreak{}[NH+]\allowbreak{}2} & 5.44 & +1 & fold 3 & control \\
CHEMBL1642973 & \texttt{CC(\allowbreak{}C)\allowbreak{}(\allowbreak{}C)\allowbreak{}c1cc[n+]\allowbreak{}(\allowbreak{}CCCC[n+]\allowbreak{}2ccc(\allowbreak{}C(\allowbreak{}C)\allowbreak{}(\allowbreak{}C)\allowbreak{}C)\allowbreak{}cc2)\allowbreak{}cc1} & \texttt{CC(\allowbreak{}C)\allowbreak{}(\allowbreak{}C)\allowbreak{}c1cc[n+]\allowbreak{}(\allowbreak{}CCCC[n+]\allowbreak{}2ccc(\allowbreak{}C(\allowbreak{}C)\allowbreak{}(\allowbreak{}C)\allowbreak{}C)\allowbreak{}cc2)\allowbreak{}cc1} & 5.43 & +2 & fold 0 & cliff \\
CHEMBL1642974 & \texttt{CC(\allowbreak{}C)\allowbreak{}(\allowbreak{}C)\allowbreak{}c1cc[n+]\allowbreak{}(\allowbreak{}CCCCC[n+]\allowbreak{}2ccc(\allowbreak{}C(\allowbreak{}C)\allowbreak{}(\allowbreak{}C)\allowbreak{}C)\allowbreak{}cc2)\allowbreak{}cc1} & \texttt{CC(\allowbreak{}C)\allowbreak{}(\allowbreak{}C)\allowbreak{}c1cc[n+]\allowbreak{}(\allowbreak{}CCCCC[n+]\allowbreak{}2ccc(\allowbreak{}C(\allowbreak{}C)\allowbreak{}(\allowbreak{}C)\allowbreak{}C)\allowbreak{}cc2)\allowbreak{}cc1} & 4.62 & +2 & fold 0 & cliff \\
CHEMBL1642975 & \texttt{CC(\allowbreak{}C)\allowbreak{}(\allowbreak{}C)\allowbreak{}c1cc[n+]\allowbreak{}(\allowbreak{}CCCCCC[n+]\allowbreak{}2ccc(\allowbreak{}C(\allowbreak{}C)\allowbreak{}(\allowbreak{}C)\allowbreak{}C)\allowbreak{}cc2)\allowbreak{}cc1} & \texttt{CC(\allowbreak{}C)\allowbreak{}(\allowbreak{}C)\allowbreak{}c1cc[n+]\allowbreak{}(\allowbreak{}CCCCCC[n+]\allowbreak{}2ccc(\allowbreak{}C(\allowbreak{}C)\allowbreak{}(\allowbreak{}C)\allowbreak{}C)\allowbreak{}cc2)\allowbreak{}cc1} & 7.50 & +2 & fold 0 & cliff \\
CHEMBL1642976 & \texttt{CC(\allowbreak{}C)\allowbreak{}(\allowbreak{}C)\allowbreak{}c1cc[n+]\allowbreak{}(\allowbreak{}CCCCCCC[n+]\allowbreak{}2ccc(\allowbreak{}C(\allowbreak{}C)\allowbreak{}(\allowbreak{}C)\allowbreak{}C)\allowbreak{}cc2)\allowbreak{}cc1} & \texttt{CC(\allowbreak{}C)\allowbreak{}(\allowbreak{}C)\allowbreak{}c1cc[n+]\allowbreak{}(\allowbreak{}CCCCCCC[n+]\allowbreak{}2ccc(\allowbreak{}C(\allowbreak{}C)\allowbreak{}(\allowbreak{}C)\allowbreak{}C)\allowbreak{}cc2)\allowbreak{}cc1} & 7.80 & +2 & fold 0 & cliff \\
CHEMBL1642977 & \texttt{CC(\allowbreak{}C)\allowbreak{}(\allowbreak{}C)\allowbreak{}c1cc[n+]\allowbreak{}(\allowbreak{}CCCCCCCC[n+]\allowbreak{}2ccc(\allowbreak{}C(\allowbreak{}C)\allowbreak{}(\allowbreak{}C)\allowbreak{}C)\allowbreak{}cc2)\allowbreak{}cc1} & \texttt{CC(\allowbreak{}C)\allowbreak{}(\allowbreak{}C)\allowbreak{}c1cc[n+]\allowbreak{}(\allowbreak{}CCCCCCCC[n+]\allowbreak{}2ccc(\allowbreak{}C(\allowbreak{}C)\allowbreak{}(\allowbreak{}C)\allowbreak{}C)\allowbreak{}cc2)\allowbreak{}cc1} & 8.30 & +2 & fold 0 & cliff \\
CHEMBL1642978 & \texttt{CC(\allowbreak{}C)\allowbreak{}(\allowbreak{}C)\allowbreak{}c1cc[n+]\allowbreak{}(\allowbreak{}CCCCCCCCC[n+]\allowbreak{}2ccc(\allowbreak{}C(\allowbreak{}C)\allowbreak{}(\allowbreak{}C)\allowbreak{}C)\allowbreak{}cc2)\allowbreak{}cc1} & \texttt{CC(\allowbreak{}C)\allowbreak{}(\allowbreak{}C)\allowbreak{}c1cc[n+]\allowbreak{}(\allowbreak{}CCCCCCCCC[n+]\allowbreak{}2ccc(\allowbreak{}C(\allowbreak{}C)\allowbreak{}(\allowbreak{}C)\allowbreak{}C)\allowbreak{}cc2)\allowbreak{}cc1} & 7.43 & +2 & fold 0 & cliff \\
CHEMBL1642979 & \texttt{CC(\allowbreak{}C)\allowbreak{}(\allowbreak{}C)\allowbreak{}c1cc[n+]\allowbreak{}(\allowbreak{}CCCCCCCCCC[n+]\allowbreak{}2ccc(\allowbreak{}C(\allowbreak{}C)\allowbreak{}(\allowbreak{}C)\allowbreak{}C)\allowbreak{}cc2)\allowbreak{}cc1} & \texttt{CC(\allowbreak{}C)\allowbreak{}(\allowbreak{}C)\allowbreak{}c1cc[n+]\allowbreak{}(\allowbreak{}CCCCCCCCCC[n+]\allowbreak{}2ccc(\allowbreak{}C(\allowbreak{}C)\allowbreak{}(\allowbreak{}C)\allowbreak{}C)\allowbreak{}cc2)\allowbreak{}cc1} & 7.92 & +2 & fold 0 & cliff \\
CHEMBL1642980 & \texttt{CC(\allowbreak{}C)\allowbreak{}(\allowbreak{}C)\allowbreak{}c1cc[n+]\allowbreak{}(\allowbreak{}CCCCCCCCCCC[n+]\allowbreak{}2ccc(\allowbreak{}C(\allowbreak{}C)\allowbreak{}(\allowbreak{}C)\allowbreak{}C)\allowbreak{}cc2)\allowbreak{}cc1} & \texttt{CC(\allowbreak{}C)\allowbreak{}(\allowbreak{}C)\allowbreak{}c1cc[n+]\allowbreak{}(\allowbreak{}CCCCCCCCCCC[n+]\allowbreak{}2ccc(\allowbreak{}C(\allowbreak{}C)\allowbreak{}(\allowbreak{}C)\allowbreak{}C)\allowbreak{}cc2)\allowbreak{}cc1} & 7.58 & +2 & fold 0 & cliff \\
CHEMBL1643944 & \texttt{c1ccc2c[n+]\allowbreak{}(\allowbreak{}CCCCCCCCCC[n+]\allowbreak{}3ccc4ccccc4c3)\allowbreak{}ccc2c1} & \texttt{c1ccc2c[n+]\allowbreak{}(\allowbreak{}CCCCCCCCCC[n+]\allowbreak{}3ccc4ccccc4c3)\allowbreak{}ccc2c1} & 7.30 & +2 & fold 1 & cliff \\
CHEMBL165304 & \texttt{c1ccc2c(\allowbreak{}c1)\allowbreak{}CN1CCCCCC1=\allowbreak{}N2} & \texttt{c1ccc2c(\allowbreak{}c1)\allowbreak{}CN1CCCCCC1=\allowbreak{}[NH+]\allowbreak{}2} & 5.42 & +1 & fold 3 & control \\
CHEMBL166359 & \texttt{c1ccc2c(\allowbreak{}c1)\allowbreak{}CN1CCCCCCCC1=\allowbreak{}N2} & \texttt{c1ccc2c(\allowbreak{}c1)\allowbreak{}CN1CCCCCCCC1=\allowbreak{}[NH+]\allowbreak{}2} & 4.87 & +1 & fold 0 & control \\
CHEMBL1669482 & \texttt{c1ccc2c[n+]\allowbreak{}(\allowbreak{}CCCC[n+]\allowbreak{}3ccc4ccccc4c3)\allowbreak{}ccc2c1} & \texttt{c1ccc2c[n+]\allowbreak{}(\allowbreak{}CCCC[n+]\allowbreak{}3ccc4ccccc4c3)\allowbreak{}ccc2c1} & 4.44 & +2 & fold 1 & cliff \\
CHEMBL1669483 & \texttt{c1ccc2c[n+]\allowbreak{}(\allowbreak{}CCCCC[n+]\allowbreak{}3ccc4ccccc4c3)\allowbreak{}ccc2c1} & \texttt{c1ccc2c[n+]\allowbreak{}(\allowbreak{}CCCCC[n+]\allowbreak{}3ccc4ccccc4c3)\allowbreak{}ccc2c1} & 4.89 & +2 & fold 1 & cliff \\
CHEMBL1669484 & \texttt{c1ccc2c[n+]\allowbreak{}(\allowbreak{}CCCCCC[n+]\allowbreak{}3ccc4ccccc4c3)\allowbreak{}ccc2c1} & \texttt{c1ccc2c[n+]\allowbreak{}(\allowbreak{}CCCCCC[n+]\allowbreak{}3ccc4ccccc4c3)\allowbreak{}ccc2c1} & 6.52 & +2 & fold 1 & cliff \\
CHEMBL1669485 & \texttt{c1ccc2c[n+]\allowbreak{}(\allowbreak{}CCCCCCC[n+]\allowbreak{}3ccc4ccccc4c3)\allowbreak{}ccc2c1} & \texttt{c1ccc2c[n+]\allowbreak{}(\allowbreak{}CCCCCCC[n+]\allowbreak{}3ccc4ccccc4c3)\allowbreak{}ccc2c1} & 6.30 & +2 & fold 1 & cliff \\
CHEMBL1669486 & \texttt{c1ccc2c[n+]\allowbreak{}(\allowbreak{}CCCCCCCC[n+]\allowbreak{}3ccc4ccccc4c3)\allowbreak{}ccc2c1} & \texttt{c1ccc2c[n+]\allowbreak{}(\allowbreak{}CCCCCCCC[n+]\allowbreak{}3ccc4ccccc4c3)\allowbreak{}ccc2c1} & 8.30 & +2 & fold 1 & cliff \\
CHEMBL1761992 & \texttt{c1ccc2c(\allowbreak{}c1)\allowbreak{}ccc[n+]\allowbreak{}2CCCC[n+]\allowbreak{}1cccc2ccccc21} & \texttt{c1ccc2c(\allowbreak{}c1)\allowbreak{}ccc[n+]\allowbreak{}2CCCC[n+]\allowbreak{}1cccc2ccccc21} & 4.42 & +2 & test & cliff \\
CHEMBL1761993 & \texttt{c1ccc2c(\allowbreak{}c1)\allowbreak{}ccc[n+]\allowbreak{}2CCCCC[n+]\allowbreak{}1cccc2ccccc21} & \texttt{c1ccc2c(\allowbreak{}c1)\allowbreak{}ccc[n+]\allowbreak{}2CCCCC[n+]\allowbreak{}1cccc2ccccc21} & 5.40 & +2 & test & cliff \\
CHEMBL1761994 & \texttt{c1ccc2c(\allowbreak{}c1)\allowbreak{}ccc[n+]\allowbreak{}2CCCCCCC[n+]\allowbreak{}1cccc2ccccc21} & \texttt{c1ccc2c(\allowbreak{}c1)\allowbreak{}ccc[n+]\allowbreak{}2CCCCCCC[n+]\allowbreak{}1cccc2ccccc21} & 7.05 & +2 & test & cliff \\
CHEMBL1761995 & \texttt{c1ccc2c(\allowbreak{}c1)\allowbreak{}ccc[n+]\allowbreak{}2CCCCCCCC[n+]\allowbreak{}1cccc2ccccc21} & \texttt{c1ccc2c(\allowbreak{}c1)\allowbreak{}ccc[n+]\allowbreak{}2CCCCCCCC[n+]\allowbreak{}1cccc2ccccc21} & 7.10 & +2 & test & cliff \\
CHEMBL1819170 & \texttt{c1ccc2c(\allowbreak{}NCCCCCCCN3CCOCC3)\allowbreak{}c3c(\allowbreak{}nc2c1)\allowbreak{}CCCC3} & \texttt{c1ccc2c(\allowbreak{}NCCCCCCC[NH+]\allowbreak{}3CCOCC3)\allowbreak{}c3c(\allowbreak{}[nH+]\allowbreak{}c2c1)\allowbreak{}CCCC3} & 7.35 & +2 & fold 0 & control \\
CHEMBL1819172 & \texttt{c1ccc2c(\allowbreak{}NCCCCCCCCCN3CCOCC3)\allowbreak{}c3c(\allowbreak{}nc2c1)\allowbreak{}CCCC3} & \texttt{c1ccc2c(\allowbreak{}NCCCCCCCCC[NH+]\allowbreak{}3CCOCC3)\allowbreak{}c3c(\allowbreak{}[nH+]\allowbreak{}c2c1)\allowbreak{}CCCC3} & 7.47 & +2 & fold 0 & control \\
CHEMBL190528 & \texttt{CNC(\allowbreak{}=\allowbreak{}O)\allowbreak{}Oc1cccc(\allowbreak{}CN(\allowbreak{}C)\allowbreak{}CCCCCCCOc2ccc3c(\allowbreak{}c2)\allowbreak{}O/\allowbreak{}C(\allowbreak{}=\allowbreak{}C\textbackslash{}c2cc(\allowbreak{}OC)\allowbreak{}c(\allowbreak{}OC)\allowbreak{}c(\allowbreak{}OC)\allowbreak{}c2)\allowbreak{}C3=\allowbreak{}O)\allowbreak{}c1} & \texttt{CNC(\allowbreak{}=\allowbreak{}O)\allowbreak{}Oc1cccc(\allowbreak{}C[NH+]\allowbreak{}(\allowbreak{}C)\allowbreak{}CCCCCCCOc2ccc3c(\allowbreak{}c2)\allowbreak{}O/\allowbreak{}C(\allowbreak{}=\allowbreak{}C\textbackslash{}c2cc(\allowbreak{}OC)\allowbreak{}c(\allowbreak{}OC)\allowbreak{}c(\allowbreak{}OC)\allowbreak{}c2)\allowbreak{}C3=\allowbreak{}O)\allowbreak{}c1} & 6.28 & +1 & fold 4 & cliff \\
CHEMBL1912058 & \texttt{O=\allowbreak{}C(\allowbreak{}CCCC[C@H]\allowbreak{}1CCSS1)\allowbreak{}NCCCNc1c2c(\allowbreak{}nc3cc(\allowbreak{}Cl)\allowbreak{}ccc13)\allowbreak{}CCCC2} & \texttt{O=\allowbreak{}C(\allowbreak{}CCCC[C@H]\allowbreak{}1CCSS1)\allowbreak{}NCCCNc1c2c(\allowbreak{}[nH+]\allowbreak{}c3cc(\allowbreak{}Cl)\allowbreak{}ccc13)\allowbreak{}CCCC2} & 9.33 & +1 & fold 3 & cliff \\
CHEMBL1912059 & \texttt{O=\allowbreak{}C(\allowbreak{}CCCC[C@@H]\allowbreak{}1CCSS1)\allowbreak{}NCCCNc1c2c(\allowbreak{}nc3cc(\allowbreak{}Cl)\allowbreak{}ccc13)\allowbreak{}CCCC2} & \texttt{O=\allowbreak{}C(\allowbreak{}CCCC[C@@H]\allowbreak{}1CCSS1)\allowbreak{}NCCCNc1c2c(\allowbreak{}[nH+]\allowbreak{}c3cc(\allowbreak{}Cl)\allowbreak{}ccc13)\allowbreak{}CCCC2} & 9.64 & +1 & fold 3 & cliff \\
CHEMBL191386 & \texttt{CNC(\allowbreak{}=\allowbreak{}O)\allowbreak{}Oc1cccc(\allowbreak{}CN(\allowbreak{}C)\allowbreak{}CCCOc2ccc3c(\allowbreak{}c2)\allowbreak{}O/\allowbreak{}C(\allowbreak{}=\allowbreak{}C\textbackslash{}c2ccc4ccccc4c2)\allowbreak{}C3=\allowbreak{}O)\allowbreak{}c1} & \texttt{CNC(\allowbreak{}=\allowbreak{}O)\allowbreak{}Oc1cccc(\allowbreak{}C[NH+]\allowbreak{}(\allowbreak{}C)\allowbreak{}CCCOc2ccc3c(\allowbreak{}c2)\allowbreak{}O/\allowbreak{}C(\allowbreak{}=\allowbreak{}C\textbackslash{}c2ccc4ccccc4c2)\allowbreak{}C3=\allowbreak{}O)\allowbreak{}c1} & 5.47 & +1 & fold 4 & cliff \\
CHEMBL191416 & \texttt{CNC(\allowbreak{}=\allowbreak{}O)\allowbreak{}Oc1cccc(\allowbreak{}CN(\allowbreak{}C)\allowbreak{}CCCOc2ccc3c(\allowbreak{}c2)\allowbreak{}O/\allowbreak{}C(\allowbreak{}=\allowbreak{}C\textbackslash{}c2cc(\allowbreak{}OC)\allowbreak{}c(\allowbreak{}OC)\allowbreak{}c(\allowbreak{}OC)\allowbreak{}c2)\allowbreak{}C3=\allowbreak{}O)\allowbreak{}c1} & \texttt{CNC(\allowbreak{}=\allowbreak{}O)\allowbreak{}Oc1cccc(\allowbreak{}C[NH+]\allowbreak{}(\allowbreak{}C)\allowbreak{}CCCOc2ccc3c(\allowbreak{}c2)\allowbreak{}O/\allowbreak{}C(\allowbreak{}=\allowbreak{}C\textbackslash{}c2cc(\allowbreak{}OC)\allowbreak{}c(\allowbreak{}OC)\allowbreak{}c(\allowbreak{}OC)\allowbreak{}c2)\allowbreak{}C3=\allowbreak{}O)\allowbreak{}c1} & 5.52 & +1 & fold 4 & cliff \\
CHEMBL2047528 & \texttt{CN(\allowbreak{}CCCOc1ccc2ccc(\allowbreak{}=\allowbreak{}O)\allowbreak{}oc2c1)\allowbreak{}Cc1cccc(\allowbreak{}OC(\allowbreak{}=\allowbreak{}O)\allowbreak{}NCCCCCc2ccccc2)\allowbreak{}c1} & \texttt{C[NH+]\allowbreak{}(\allowbreak{}CCCOc1ccc2ccc(\allowbreak{}=\allowbreak{}O)\allowbreak{}oc2c1)\allowbreak{}Cc1cccc(\allowbreak{}OC(\allowbreak{}=\allowbreak{}O)\allowbreak{}NCCCCCc2ccccc2)\allowbreak{}c1} & 7.13 & +1 & fold 4 & control \\
CHEMBL2178782 & \texttt{COc1ccc2c3c1O[C@H]\allowbreak{}1C[C@@H]\allowbreak{}(\allowbreak{}O)\allowbreak{}C=\allowbreak{}C[C@@]\allowbreak{}31CCN(\allowbreak{}CCCCNC13CC4CC(\allowbreak{}C)\allowbreak{}(\allowbreak{}CC(\allowbreak{}C)\allowbreak{}(\allowbreak{}C4)\allowbreak{}C1)\allowbreak{}C3)\allowbreak{}C2} & \texttt{COc1ccc2c3c1O[C@H]\allowbreak{}1C[C@@H]\allowbreak{}(\allowbreak{}O)\allowbreak{}C=\allowbreak{}C[C@@]\allowbreak{}31CC[NH+]\allowbreak{}(\allowbreak{}CCCC[NH2+]\allowbreak{}C13CC4CC(\allowbreak{}C)\allowbreak{}(\allowbreak{}CC(\allowbreak{}C)\allowbreak{}(\allowbreak{}C4)\allowbreak{}C1)\allowbreak{}C3)\allowbreak{}C2} & 5.95 & +2 & test & cliff \\
CHEMBL2178784 & \texttt{COc1ccc2c3c1O[C@H]\allowbreak{}1C[C@@H]\allowbreak{}(\allowbreak{}O)\allowbreak{}C=\allowbreak{}C[C@@]\allowbreak{}31CCN(\allowbreak{}CCCCCCNC13CC4CC(\allowbreak{}C)\allowbreak{}(\allowbreak{}CC(\allowbreak{}C)\allowbreak{}(\allowbreak{}C4)\allowbreak{}C1)\allowbreak{}C3)\allowbreak{}C2} & \texttt{COc1ccc2c3c1O[C@H]\allowbreak{}1C[C@@H]\allowbreak{}(\allowbreak{}O)\allowbreak{}C=\allowbreak{}C[C@@]\allowbreak{}31CC[NH+]\allowbreak{}(\allowbreak{}CCCCCC[NH2+]\allowbreak{}C13CC4CC(\allowbreak{}C)\allowbreak{}(\allowbreak{}CC(\allowbreak{}C)\allowbreak{}(\allowbreak{}C4)\allowbreak{}C1)\allowbreak{}C3)\allowbreak{}C2} & 8.94 & +2 & test & cliff \\
CHEMBL219550 & \texttt{COc1cc(\allowbreak{}C(\allowbreak{}=\allowbreak{}O)\allowbreak{}NCCCCC(\allowbreak{}=\allowbreak{}O)\allowbreak{}NNc2c3c(\allowbreak{}nc4ccccc24)\allowbreak{}CCCC3)\allowbreak{}cc(\allowbreak{}OC)\allowbreak{}c1OC} & \texttt{COc1cc(\allowbreak{}C(\allowbreak{}=\allowbreak{}O)\allowbreak{}NCCCCC(\allowbreak{}=\allowbreak{}O)\allowbreak{}NNc2c3c(\allowbreak{}[nH+]\allowbreak{}c4ccccc24)\allowbreak{}CCCC3)\allowbreak{}cc(\allowbreak{}OC)\allowbreak{}c1OC} & 7.77 & +1 & test & control \\
CHEMBL220141 & \texttt{COc1cc(\allowbreak{}CCC(\allowbreak{}=\allowbreak{}O)\allowbreak{}NCCC(\allowbreak{}=\allowbreak{}O)\allowbreak{}NNc2c3c(\allowbreak{}nc4ccccc24)\allowbreak{}CCCC3)\allowbreak{}cc(\allowbreak{}OC)\allowbreak{}c1OC} & \texttt{COc1cc(\allowbreak{}CCC(\allowbreak{}=\allowbreak{}O)\allowbreak{}NCCC(\allowbreak{}=\allowbreak{}O)\allowbreak{}NNc2c3c(\allowbreak{}[nH+]\allowbreak{}c4ccccc24)\allowbreak{}CCCC3)\allowbreak{}cc(\allowbreak{}OC)\allowbreak{}c1OC} & 6.87 & +1 & fold 2 & control \\
CHEMBL223714 & \texttt{O=\allowbreak{}C(\allowbreak{}NCCCCCCCCNc1c2c(\allowbreak{}nc3cc(\allowbreak{}Cl)\allowbreak{}ccc13)\allowbreak{}CCCC2)\allowbreak{}c1c[nH]\allowbreak{}c2ccccc12} & \texttt{O=\allowbreak{}C(\allowbreak{}NCCCCCCCCNc1c2c(\allowbreak{}[nH+]\allowbreak{}c3cc(\allowbreak{}Cl)\allowbreak{}ccc13)\allowbreak{}CCCC2)\allowbreak{}c1c[nH]\allowbreak{}c2ccccc12} & 7.34 & +1 & fold 1 & control \\
CHEMBL2238001 & \texttt{Nc1c2c(\allowbreak{}nc3nc(\allowbreak{}-c4ccccc4)\allowbreak{}cc(\allowbreak{}C(\allowbreak{}F)\allowbreak{}(\allowbreak{}F)\allowbreak{}F)\allowbreak{}c13)\allowbreak{}CCCCC2} & \texttt{Nc1c2c(\allowbreak{}[nH+]\allowbreak{}c3[nH+]\allowbreak{}c(\allowbreak{}-c4ccccc4)\allowbreak{}cc(\allowbreak{}C(\allowbreak{}F)\allowbreak{}(\allowbreak{}F)\allowbreak{}F)\allowbreak{}c13)\allowbreak{}CCCCC2} & 4.99 & +2 & fold 1 & control \\
CHEMBL2238002 & \texttt{Nc1c2c(\allowbreak{}nc3nc(\allowbreak{}-c4ccccc4)\allowbreak{}cc(\allowbreak{}C(\allowbreak{}F)\allowbreak{}(\allowbreak{}F)\allowbreak{}F)\allowbreak{}c13)\allowbreak{}CCCCCC2} & \texttt{Nc1c2c(\allowbreak{}[nH+]\allowbreak{}c3[nH+]\allowbreak{}c(\allowbreak{}-c4ccccc4)\allowbreak{}cc(\allowbreak{}C(\allowbreak{}F)\allowbreak{}(\allowbreak{}F)\allowbreak{}F)\allowbreak{}c13)\allowbreak{}CCCCCC2} & 5.11 & +2 & fold 1 & control \\
CHEMBL224158 & \texttt{O=\allowbreak{}C(\allowbreak{}CCc1c[nH]\allowbreak{}c2ccccc12)\allowbreak{}NCCCCCNc1c2c(\allowbreak{}nc3ccccc13)\allowbreak{}CCCC2} & \texttt{O=\allowbreak{}C(\allowbreak{}CCc1c[nH]\allowbreak{}c2ccccc12)\allowbreak{}NCCCCCNc1c2c(\allowbreak{}[nH+]\allowbreak{}c3ccccc13)\allowbreak{}CCCC2} & 7.16 & +1 & fold 1 & cliff \\
CHEMBL224335 & \texttt{O=\allowbreak{}C(\allowbreak{}CCc1c[nH]\allowbreak{}c2ccccc12)\allowbreak{}NCCCCCNc1c2c(\allowbreak{}nc3cc(\allowbreak{}Cl)\allowbreak{}ccc13)\allowbreak{}CCCC2} & \texttt{O=\allowbreak{}C(\allowbreak{}CCc1c[nH]\allowbreak{}c2ccccc12)\allowbreak{}NCCCCCNc1c2c(\allowbreak{}[nH+]\allowbreak{}c3cc(\allowbreak{}Cl)\allowbreak{}ccc13)\allowbreak{}CCCC2} & 8.40 & +1 & fold 4 & cliff \\
CHEMBL224433 & \texttt{CN(\allowbreak{}CCCCCCCOc1ccc2c(\allowbreak{}=\allowbreak{}O)\allowbreak{}c3ccccc3oc2c1)\allowbreak{}Cc1ccccc1} & \texttt{C[NH+]\allowbreak{}(\allowbreak{}CCCCCCCOc1ccc2c(\allowbreak{}=\allowbreak{}O)\allowbreak{}c3ccccc3oc2c1)\allowbreak{}Cc1ccccc1} & 5.55 & +1 & fold 1 & control \\
CHEMBL224469 & \texttt{CN(\allowbreak{}CCCCCCCCOc1ccc2c(\allowbreak{}=\allowbreak{}O)\allowbreak{}c3ccccc3oc2c1)\allowbreak{}Cc1ccccc1} & \texttt{C[NH+]\allowbreak{}(\allowbreak{}CCCCCCCCOc1ccc2c(\allowbreak{}=\allowbreak{}O)\allowbreak{}c3ccccc3oc2c1)\allowbreak{}Cc1ccccc1} & 4.99 & +1 & fold 4 & control \\
CHEMBL224998 & \texttt{O=\allowbreak{}C(\allowbreak{}NCCCCCCNc1c2c(\allowbreak{}nc3cc(\allowbreak{}Cl)\allowbreak{}ccc13)\allowbreak{}CCCC2)\allowbreak{}c1c[nH]\allowbreak{}c2ccccc12} & \texttt{O=\allowbreak{}C(\allowbreak{}NCCCCCCNc1c2c(\allowbreak{}[nH+]\allowbreak{}c3cc(\allowbreak{}Cl)\allowbreak{}ccc13)\allowbreak{}CCCC2)\allowbreak{}c1c[nH]\allowbreak{}c2ccccc12} & 7.48 & +1 & fold 1 & control \\
CHEMBL225021 & \texttt{O=\allowbreak{}C(\allowbreak{}CCCc1c[nH]\allowbreak{}c2ccccc12)\allowbreak{}NCCCCCNc1c2c(\allowbreak{}nc3cc(\allowbreak{}Cl)\allowbreak{}ccc13)\allowbreak{}CCCC2} & \texttt{O=\allowbreak{}C(\allowbreak{}CCCc1c[nH]\allowbreak{}c2ccccc12)\allowbreak{}NCCCCCNc1c2c(\allowbreak{}[nH+]\allowbreak{}c3cc(\allowbreak{}Cl)\allowbreak{}ccc13)\allowbreak{}CCCC2} & 9.71 & +1 & fold 4 & cliff \\
CHEMBL225567 & \texttt{O=\allowbreak{}C(\allowbreak{}CCc1c[nH]\allowbreak{}c2ccccc12)\allowbreak{}NCCCCCCNc1c2c(\allowbreak{}nc3cc(\allowbreak{}Cl)\allowbreak{}ccc13)\allowbreak{}CCCC2} & \texttt{O=\allowbreak{}C(\allowbreak{}CCc1c[nH]\allowbreak{}c2ccccc12)\allowbreak{}NCCCCCCNc1c2c(\allowbreak{}[nH+]\allowbreak{}c3cc(\allowbreak{}Cl)\allowbreak{}ccc13)\allowbreak{}CCCC2} & 10.43 & +1 & fold 4 & cliff \\
CHEMBL225577 & \texttt{c1ccc2c(\allowbreak{}c1)\allowbreak{}ccc[n+]\allowbreak{}2CCCCCCCCCC[n+]\allowbreak{}1cccc2ccccc21} & \texttt{c1ccc2c(\allowbreak{}c1)\allowbreak{}ccc[n+]\allowbreak{}2CCCCCCCCCC[n+]\allowbreak{}1cccc2ccccc21} & 9.00 & +2 & test & cliff \\
CHEMBL225611 & \texttt{CN(\allowbreak{}CCCCOc1ccc2c(\allowbreak{}=\allowbreak{}O)\allowbreak{}c3ccccc3oc2c1)\allowbreak{}Cc1ccccc1} & \texttt{C[NH+]\allowbreak{}(\allowbreak{}CCCCOc1ccc2c(\allowbreak{}=\allowbreak{}O)\allowbreak{}c3ccccc3oc2c1)\allowbreak{}Cc1ccccc1} & 5.20 & +1 & fold 4 & control \\
CHEMBL225689 & \texttt{c1cc[n+]\allowbreak{}(\allowbreak{}CCCCCCCCCC[n+]\allowbreak{}2ccccc2)\allowbreak{}cc1} & \texttt{c1cc[n+]\allowbreak{}(\allowbreak{}CCCCCCCCCC[n+]\allowbreak{}2ccccc2)\allowbreak{}cc1} & 6.40 & +2 & fold 1 & cliff \\
CHEMBL2419683 & \texttt{N\#Cc1c(\allowbreak{}Cl)\allowbreak{}nc(\allowbreak{}NCCCC2CCN(\allowbreak{}Cc3ccccc3)\allowbreak{}CC2)\allowbreak{}c(\allowbreak{}C\#N)\allowbreak{}c1-c1ccccc1} & \texttt{N\#Cc1c(\allowbreak{}Cl)\allowbreak{}[nH+]\allowbreak{}c(\allowbreak{}NCCCC2CC[NH+]\allowbreak{}(\allowbreak{}Cc3ccccc3)\allowbreak{}CC2)\allowbreak{}c(\allowbreak{}C\#N)\allowbreak{}c1-c1ccccc1} & 5.41 & +2 & test & cliff \\
CHEMBL2419684 & \texttt{N\#Cc1c(\allowbreak{}Cl)\allowbreak{}nc(\allowbreak{}NCCCCC2CCN(\allowbreak{}Cc3ccccc3)\allowbreak{}CC2)\allowbreak{}c(\allowbreak{}C\#N)\allowbreak{}c1-c1ccccc1} & \texttt{N\#Cc1c(\allowbreak{}Cl)\allowbreak{}[nH+]\allowbreak{}c(\allowbreak{}NCCCCC2CC[NH+]\allowbreak{}(\allowbreak{}Cc3ccccc3)\allowbreak{}CC2)\allowbreak{}c(\allowbreak{}C\#N)\allowbreak{}c1-c1ccccc1} & 7.51 & +2 & test & cliff \\
CHEMBL2441668 & \texttt{COc1cccc(\allowbreak{}CCCCCCCc2cccc(\allowbreak{}OC)\allowbreak{}[n+]\allowbreak{}2C)\allowbreak{}[n+]\allowbreak{}1C} & \texttt{COc1cccc(\allowbreak{}CCCCCCCc2cccc(\allowbreak{}OC)\allowbreak{}[n+]\allowbreak{}2C)\allowbreak{}[n+]\allowbreak{}1C} & 6.92 & +2 & fold 2 & control \\
CHEMBL2441669 & \texttt{COc1cccc(\allowbreak{}CCCCCCCCc2cccc(\allowbreak{}OC)\allowbreak{}[n+]\allowbreak{}2C)\allowbreak{}[n+]\allowbreak{}1C} & \texttt{COc1cccc(\allowbreak{}CCCCCCCCc2cccc(\allowbreak{}OC)\allowbreak{}[n+]\allowbreak{}2C)\allowbreak{}[n+]\allowbreak{}1C} & 6.73 & +2 & fold 2 & control \\
CHEMBL258928 & \texttt{NCCCCCCCCNc1c2c(\allowbreak{}nc3ccccc13)\allowbreak{}CCCC2} & \texttt{[NH3+]\allowbreak{}CCCCCCCCNc1c2c(\allowbreak{}[nH+]\allowbreak{}c3ccccc13)\allowbreak{}CCCC2} & 7.05 & +2 & fold 3 & control \\
CHEMBL284369 & \texttt{CCCCCCNC(\allowbreak{}=\allowbreak{}O)\allowbreak{}Oc1ccc2c(\allowbreak{}c1)\allowbreak{}[C@@H]\allowbreak{}1CCN(\allowbreak{}CC)\allowbreak{}C1C2} & \texttt{CCCCCCNC(\allowbreak{}=\allowbreak{}O)\allowbreak{}Oc1ccc2c(\allowbreak{}c1)\allowbreak{}[C@@H]\allowbreak{}1CC[NH+]\allowbreak{}(\allowbreak{}CC)\allowbreak{}C1C2} & 6.68 & +1 & test & control \\
CHEMBL303574 & \texttt{CN1CC[C@@]\allowbreak{}2(\allowbreak{}C)\allowbreak{}c3cc(\allowbreak{}OC(\allowbreak{}=\allowbreak{}O)\allowbreak{}NCCCCCN4CCOCC4)\allowbreak{}ccc3N(\allowbreak{}C)\allowbreak{}[C@@H]\allowbreak{}12} & \texttt{CN1c2ccc(\allowbreak{}OC(\allowbreak{}=\allowbreak{}O)\allowbreak{}NCCCCC[NH+]\allowbreak{}3CCOCC3)\allowbreak{}cc2[C@]\allowbreak{}2(\allowbreak{}C)\allowbreak{}CC[NH+]\allowbreak{}(\allowbreak{}C)\allowbreak{}[C@H]\allowbreak{}12} & 5.12 & +2 & fold 2 & cliff \\
CHEMBL3093173 & \texttt{COc1ccc(\allowbreak{}/\allowbreak{}C=\allowbreak{}C2\textbackslash{}COc3cc(\allowbreak{}OCCCCCNc4c5c(\allowbreak{}nc6ccccc46)\allowbreak{}CCCC5)\allowbreak{}ccc3C2=\allowbreak{}O)\allowbreak{}cc1} & \texttt{COc1ccc(\allowbreak{}/\allowbreak{}C=\allowbreak{}C2\textbackslash{}COc3cc(\allowbreak{}OCCCCCNc4c5c(\allowbreak{}[nH+]\allowbreak{}c6ccccc46)\allowbreak{}CCCC5)\allowbreak{}ccc3C2=\allowbreak{}O)\allowbreak{}cc1} & 7.01 & +1 & fold 4 & control \\
CHEMBL3093175 & \texttt{COc1ccc(\allowbreak{}/\allowbreak{}C=\allowbreak{}C2\textbackslash{}COc3cc(\allowbreak{}OCCCCCCCCNc4c5c(\allowbreak{}nc6ccccc46)\allowbreak{}CCCC5)\allowbreak{}ccc3C2=\allowbreak{}O)\allowbreak{}cc1} & \texttt{COc1ccc(\allowbreak{}/\allowbreak{}C=\allowbreak{}C2\textbackslash{}COc3cc(\allowbreak{}OCCCCCCCCNc4c5c(\allowbreak{}[nH+]\allowbreak{}c6ccccc46)\allowbreak{}CCCC5)\allowbreak{}ccc3C2=\allowbreak{}O)\allowbreak{}cc1} & 6.92 & +1 & fold 4 & control \\
CHEMBL3093797 & \texttt{c1cc[n+]\allowbreak{}(\allowbreak{}CCCCCCCCCC[n+]\allowbreak{}2ccc3ccccc3c2)\allowbreak{}cc1} & \texttt{c1cc[n+]\allowbreak{}(\allowbreak{}CCCCCCCCCC[n+]\allowbreak{}2ccc3ccccc3c2)\allowbreak{}cc1} & 7.70 & +2 & fold 1 & control \\
CHEMBL3093798 & \texttt{c1cc[n+]\allowbreak{}(\allowbreak{}CCCCCCCCCCC[n+]\allowbreak{}2ccc3ccccc3c2)\allowbreak{}cc1} & \texttt{c1cc[n+]\allowbreak{}(\allowbreak{}CCCCCCCCCCC[n+]\allowbreak{}2ccc3ccccc3c2)\allowbreak{}cc1} & 7.52 & +2 & fold 1 & control \\
CHEMBL3093799 & \texttt{c1cc[n+]\allowbreak{}(\allowbreak{}CCCCCCC[n+]\allowbreak{}2cccc3ccccc32)\allowbreak{}cc1} & \texttt{c1cc[n+]\allowbreak{}(\allowbreak{}CCCCCCC[n+]\allowbreak{}2cccc3ccccc32)\allowbreak{}cc1} & 6.27 & +2 & fold 1 & cliff \\
CHEMBL3093802 & \texttt{c1cc[n+]\allowbreak{}(\allowbreak{}CCCCCCCCCC[n+]\allowbreak{}2cccc3ccccc32)\allowbreak{}cc1} & \texttt{c1cc[n+]\allowbreak{}(\allowbreak{}CCCCCCCCCC[n+]\allowbreak{}2cccc3ccccc32)\allowbreak{}cc1} & 8.30 & +2 & fold 1 & cliff \\
CHEMBL3093804 & \texttt{c1cc[n+]\allowbreak{}(\allowbreak{}CCCCCCCCCCCC[n+]\allowbreak{}2cccc3ccccc32)\allowbreak{}cc1} & \texttt{c1cc[n+]\allowbreak{}(\allowbreak{}CCCCCCCCCCCC[n+]\allowbreak{}2cccc3ccccc32)\allowbreak{}cc1} & 8.33 & +2 & fold 1 & cliff \\
CHEMBL3216554 & \texttt{CC1=\allowbreak{}CC2Cc3nc4cc(\allowbreak{}Cl)\allowbreak{}ccc4c(\allowbreak{}NCCCCCCCCCCNc4c5c(\allowbreak{}nc6ccccc46)\allowbreak{}CCCC5)\allowbreak{}c3C(\allowbreak{}C1)\allowbreak{}C2} & \texttt{CC1=\allowbreak{}CC2Cc3[nH+]\allowbreak{}c4cc(\allowbreak{}Cl)\allowbreak{}ccc4c(\allowbreak{}NCCCCCCCCCCNc4c5c(\allowbreak{}[nH+]\allowbreak{}c6ccccc46)\allowbreak{}CCCC5)\allowbreak{}c3C(\allowbreak{}C1)\allowbreak{}C2} & 8.38 & +2 & fold 1 & control \\
CHEMBL3216778 & \texttt{CC1=\allowbreak{}CC2Cc3nc4cc(\allowbreak{}Cl)\allowbreak{}ccc4c(\allowbreak{}NCCCCCCCCCNc4c5c(\allowbreak{}nc6cc(\allowbreak{}Cl)\allowbreak{}ccc46)\allowbreak{}CCCC5)\allowbreak{}c3C(\allowbreak{}C1)\allowbreak{}C2} & \texttt{CC1=\allowbreak{}CC2Cc3[nH+]\allowbreak{}c4cc(\allowbreak{}Cl)\allowbreak{}ccc4c(\allowbreak{}NCCCCCCCCCNc4c5c(\allowbreak{}[nH+]\allowbreak{}c6cc(\allowbreak{}Cl)\allowbreak{}ccc46)\allowbreak{}CCCC5)\allowbreak{}c3C(\allowbreak{}C1)\allowbreak{}C2} & 8.49 & +2 & fold 1 & control \\
CHEMBL3233826 & \texttt{CC1=\allowbreak{}CC2Cc3nc4cc(\allowbreak{}Cl)\allowbreak{}ccc4c(\allowbreak{}NCCCCCNC(\allowbreak{}=\allowbreak{}O)\allowbreak{}c4cc(\allowbreak{}O)\allowbreak{}c5c(\allowbreak{}c4)\allowbreak{}C(\allowbreak{}=\allowbreak{}O)\allowbreak{}c4cccc(\allowbreak{}O)\allowbreak{}c4C5=\allowbreak{}O)\allowbreak{}c3C(\allowbreak{}C1)\allowbreak{}C2} & \texttt{CC1=\allowbreak{}CC2Cc3[nH+]\allowbreak{}c4cc(\allowbreak{}Cl)\allowbreak{}ccc4c(\allowbreak{}NCCCCCNC(\allowbreak{}=\allowbreak{}O)\allowbreak{}c4cc(\allowbreak{}O)\allowbreak{}c5c(\allowbreak{}c4)\allowbreak{}C(\allowbreak{}=\allowbreak{}O)\allowbreak{}c4cccc(\allowbreak{}O)\allowbreak{}c4C5=\allowbreak{}O)\allowbreak{}c3C(\allowbreak{}C1)\allowbreak{}C2} & 8.97 & +1 & fold 3 & cliff \\
CHEMBL3233827 & \texttt{CC1=\allowbreak{}CC2Cc3nc4cc(\allowbreak{}Cl)\allowbreak{}ccc4c(\allowbreak{}NCCCCCCNC(\allowbreak{}=\allowbreak{}O)\allowbreak{}c4cc(\allowbreak{}O)\allowbreak{}c5c(\allowbreak{}c4)\allowbreak{}C(\allowbreak{}=\allowbreak{}O)\allowbreak{}c4cccc(\allowbreak{}O)\allowbreak{}c4C5=\allowbreak{}O)\allowbreak{}c3C(\allowbreak{}C1)\allowbreak{}C2} & \texttt{CC1=\allowbreak{}CC2Cc3[nH+]\allowbreak{}c4cc(\allowbreak{}Cl)\allowbreak{}ccc4c(\allowbreak{}NCCCCCCNC(\allowbreak{}=\allowbreak{}O)\allowbreak{}c4cc(\allowbreak{}O)\allowbreak{}c5c(\allowbreak{}c4)\allowbreak{}C(\allowbreak{}=\allowbreak{}O)\allowbreak{}c4cccc(\allowbreak{}O)\allowbreak{}c4C5=\allowbreak{}O)\allowbreak{}c3C(\allowbreak{}C1)\allowbreak{}C2} & 8.82 & +1 & fold 3 & cliff \\
CHEMBL3233828 & \texttt{CC1=\allowbreak{}CC2Cc3nc4cc(\allowbreak{}Cl)\allowbreak{}ccc4c(\allowbreak{}NCCCCCCCNC(\allowbreak{}=\allowbreak{}O)\allowbreak{}c4cc(\allowbreak{}O)\allowbreak{}c5c(\allowbreak{}c4)\allowbreak{}C(\allowbreak{}=\allowbreak{}O)\allowbreak{}c4cccc(\allowbreak{}O)\allowbreak{}c4C5=\allowbreak{}O)\allowbreak{}c3C(\allowbreak{}C1)\allowbreak{}C2} & \texttt{CC1=\allowbreak{}CC2Cc3[nH+]\allowbreak{}c4cc(\allowbreak{}Cl)\allowbreak{}ccc4c(\allowbreak{}NCCCCCCCNC(\allowbreak{}=\allowbreak{}O)\allowbreak{}c4cc(\allowbreak{}O)\allowbreak{}c5c(\allowbreak{}c4)\allowbreak{}C(\allowbreak{}=\allowbreak{}O)\allowbreak{}c4cccc(\allowbreak{}O)\allowbreak{}c4C5=\allowbreak{}O)\allowbreak{}c3C(\allowbreak{}C1)\allowbreak{}C2} & 8.50 & +1 & fold 3 & cliff \\
CHEMBL3233829 & \texttt{CC1=\allowbreak{}CC2Cc3nc4cc(\allowbreak{}Cl)\allowbreak{}ccc4c(\allowbreak{}NCCCCCCCCNC(\allowbreak{}=\allowbreak{}O)\allowbreak{}c4cc(\allowbreak{}O)\allowbreak{}c5c(\allowbreak{}c4)\allowbreak{}C(\allowbreak{}=\allowbreak{}O)\allowbreak{}c4cccc(\allowbreak{}O)\allowbreak{}c4C5=\allowbreak{}O)\allowbreak{}c3C(\allowbreak{}C1)\allowbreak{}C2} & \texttt{CC1=\allowbreak{}CC2Cc3[nH+]\allowbreak{}c4cc(\allowbreak{}Cl)\allowbreak{}ccc4c(\allowbreak{}NCCCCCCCCNC(\allowbreak{}=\allowbreak{}O)\allowbreak{}c4cc(\allowbreak{}O)\allowbreak{}c5c(\allowbreak{}c4)\allowbreak{}C(\allowbreak{}=\allowbreak{}O)\allowbreak{}c4cccc(\allowbreak{}O)\allowbreak{}c4C5=\allowbreak{}O)\allowbreak{}c3C(\allowbreak{}C1)\allowbreak{}C2} & 8.36 & +1 & fold 3 & cliff \\
CHEMBL3233831 & \texttt{CC1=\allowbreak{}C[C@H]\allowbreak{}2Cc3nc4cc(\allowbreak{}Cl)\allowbreak{}ccc4c(\allowbreak{}NCCCCCCCCCNC(\allowbreak{}=\allowbreak{}O)\allowbreak{}c4cc(\allowbreak{}O)\allowbreak{}c5c(\allowbreak{}c4)\allowbreak{}C(\allowbreak{}=\allowbreak{}O)\allowbreak{}c4cccc(\allowbreak{}O)\allowbreak{}c4C5=\allowbreak{}O)\allowbreak{}c3[C@@H]\allowbreak{}(\allowbreak{}C1)\allowbreak{}C2} & \texttt{CC1=\allowbreak{}C[C@H]\allowbreak{}2Cc3[nH+]\allowbreak{}c4cc(\allowbreak{}Cl)\allowbreak{}ccc4c(\allowbreak{}NCCCCCCCCCNC(\allowbreak{}=\allowbreak{}O)\allowbreak{}c4cc(\allowbreak{}O)\allowbreak{}c5c(\allowbreak{}c4)\allowbreak{}C(\allowbreak{}=\allowbreak{}O)\allowbreak{}c4cccc(\allowbreak{}O)\allowbreak{}c4C5=\allowbreak{}O)\allowbreak{}c3[C@@H]\allowbreak{}(\allowbreak{}C1)\allowbreak{}C2} & 5.53 & +1 & fold 3 & cliff \\
CHEMBL32453 & \texttt{CCCCCCCNC(\allowbreak{}=\allowbreak{}O)\allowbreak{}Oc1ccc2c(\allowbreak{}c1)\allowbreak{}[C@@H]\allowbreak{}1CCN(\allowbreak{}CC)\allowbreak{}C1C2} & \texttt{CCCCCCCNC(\allowbreak{}=\allowbreak{}O)\allowbreak{}Oc1ccc2c(\allowbreak{}c1)\allowbreak{}[C@@H]\allowbreak{}1CC[NH+]\allowbreak{}(\allowbreak{}CC)\allowbreak{}C1C2} & 6.42 & +1 & test & control \\
CHEMBL32543 & \texttt{CCCCCCNC(\allowbreak{}=\allowbreak{}O)\allowbreak{}Oc1ccc2c(\allowbreak{}c1)\allowbreak{}[C@@H]\allowbreak{}1CCN(\allowbreak{}CCC)\allowbreak{}C1C2} & \texttt{CCCCCCNC(\allowbreak{}=\allowbreak{}O)\allowbreak{}Oc1ccc2c(\allowbreak{}c1)\allowbreak{}[C@@H]\allowbreak{}1CC[NH+]\allowbreak{}(\allowbreak{}CCC)\allowbreak{}C1C2} & 6.28 & +1 & fold 1 & control \\
CHEMBL3262662 & \texttt{CCCN(\allowbreak{}CCC)\allowbreak{}CCCCCCOc1ccc(\allowbreak{}CC2NCCc3cc(\allowbreak{}OC)\allowbreak{}c(\allowbreak{}OC)\allowbreak{}cc32)\allowbreak{}cc1} & \texttt{CCC[NH+]\allowbreak{}(\allowbreak{}CCC)\allowbreak{}CCCCCCOc1ccc(\allowbreak{}CC2[NH2+]\allowbreak{}CCc3cc(\allowbreak{}OC)\allowbreak{}c(\allowbreak{}OC)\allowbreak{}cc32)\allowbreak{}cc1} & 5.46 & +2 & fold 2 & control \\
CHEMBL3262663 & \texttt{CCCCN(\allowbreak{}CCCC)\allowbreak{}CCCCCCOc1ccc(\allowbreak{}CC2NCCc3cc(\allowbreak{}OC)\allowbreak{}c(\allowbreak{}OC)\allowbreak{}cc32)\allowbreak{}cc1} & \texttt{CCCC[NH+]\allowbreak{}(\allowbreak{}CCCC)\allowbreak{}CCCCCCOc1ccc(\allowbreak{}CC2[NH2+]\allowbreak{}CCc3cc(\allowbreak{}OC)\allowbreak{}c(\allowbreak{}OC)\allowbreak{}cc32)\allowbreak{}cc1} & 5.42 & +2 & fold 2 & control \\
CHEMBL32823 & \texttt{c1ccc2c(\allowbreak{}NCCCCCCCNc3c4c(\allowbreak{}nc5ccccc35)\allowbreak{}CCCC4)\allowbreak{}c3c(\allowbreak{}nc2c1)\allowbreak{}CCCC3} & \texttt{c1ccc2c(\allowbreak{}NCCCCCCCNc3c4c(\allowbreak{}[nH+]\allowbreak{}c5ccccc35)\allowbreak{}CCCC4)\allowbreak{}c3c(\allowbreak{}[nH+]\allowbreak{}c2c1)\allowbreak{}CCCC3} & 9.16 & +2 & fold 0 & cliff \\
CHEMBL330883 & \texttt{Nc1c2c(\allowbreak{}nc3ccccc13)\allowbreak{}CCCCC2} & \texttt{Nc1c2c(\allowbreak{}[nH+]\allowbreak{}c3ccccc13)\allowbreak{}CCCCC2} & 7.41 & +1 & fold 2 & cliff \\
CHEMBL33098 & \texttt{CCCCCCCNC(\allowbreak{}=\allowbreak{}O)\allowbreak{}Oc1ccc2c(\allowbreak{}c1)\allowbreak{}[C@@H]\allowbreak{}1CCN(\allowbreak{}CCC)\allowbreak{}C1C2} & \texttt{CCCCCCCNC(\allowbreak{}=\allowbreak{}O)\allowbreak{}Oc1ccc2c(\allowbreak{}c1)\allowbreak{}[C@@H]\allowbreak{}1CC[NH+]\allowbreak{}(\allowbreak{}CCC)\allowbreak{}C1C2} & 6.20 & +1 & fold 1 & control \\
CHEMBL3322158 & \texttt{O=\allowbreak{}C(\allowbreak{}Cc1cc(\allowbreak{}=\allowbreak{}O)\allowbreak{}oc2cc(\allowbreak{}O)\allowbreak{}ccc12)\allowbreak{}NCCNc1c2c(\allowbreak{}nc3ccccc13)\allowbreak{}CCCC2} & \texttt{O=\allowbreak{}C(\allowbreak{}Cc1cc(\allowbreak{}=\allowbreak{}O)\allowbreak{}oc2cc(\allowbreak{}O)\allowbreak{}ccc12)\allowbreak{}NCCNc1c2c(\allowbreak{}[nH+]\allowbreak{}c3ccccc13)\allowbreak{}CCCC2} & 5.54 & +1 & fold 0 & cliff \\
CHEMBL3322160 & \texttt{O=\allowbreak{}C(\allowbreak{}Cc1cc(\allowbreak{}=\allowbreak{}O)\allowbreak{}oc2cc(\allowbreak{}O)\allowbreak{}ccc12)\allowbreak{}NCCCCNc1c2c(\allowbreak{}nc3ccccc13)\allowbreak{}CCCC2} & \texttt{O=\allowbreak{}C(\allowbreak{}Cc1cc(\allowbreak{}=\allowbreak{}O)\allowbreak{}oc2cc(\allowbreak{}O)\allowbreak{}ccc12)\allowbreak{}NCCCCNc1c2c(\allowbreak{}[nH+]\allowbreak{}c3ccccc13)\allowbreak{}CCCC2} & 7.81 & +1 & fold 0 & cliff \\
CHEMBL3322233 & \texttt{O=\allowbreak{}C(\allowbreak{}NCCCCNc1c2c(\allowbreak{}nc3cc(\allowbreak{}Cl)\allowbreak{}ccc13)\allowbreak{}CCCC2)\allowbreak{}c1ccc2nc(\allowbreak{}-c3ccc(\allowbreak{}Cl)\allowbreak{}cc3)\allowbreak{}c3c(\allowbreak{}c2c1)\allowbreak{}NCCC3} & \texttt{O=\allowbreak{}C(\allowbreak{}NCCCCNc1c2c(\allowbreak{}[nH+]\allowbreak{}c3cc(\allowbreak{}Cl)\allowbreak{}ccc13)\allowbreak{}CCCC2)\allowbreak{}c1ccc2[nH+]\allowbreak{}c(\allowbreak{}-c3ccc(\allowbreak{}Cl)\allowbreak{}cc3)\allowbreak{}c3c(\allowbreak{}c2c1)\allowbreak{}NCCC3} & 7.85 & +2 & test & control \\
CHEMBL3322234 & \texttt{O=\allowbreak{}C(\allowbreak{}NCCCCCNc1c2c(\allowbreak{}nc3cc(\allowbreak{}Cl)\allowbreak{}ccc13)\allowbreak{}CCCC2)\allowbreak{}c1ccc2nc(\allowbreak{}-c3ccc(\allowbreak{}Cl)\allowbreak{}cc3)\allowbreak{}c3c(\allowbreak{}c2c1)\allowbreak{}NCCC3} & \texttt{O=\allowbreak{}C(\allowbreak{}NCCCCCNc1c2c(\allowbreak{}[nH+]\allowbreak{}c3cc(\allowbreak{}Cl)\allowbreak{}ccc13)\allowbreak{}CCCC2)\allowbreak{}c1ccc2[nH+]\allowbreak{}c(\allowbreak{}-c3ccc(\allowbreak{}Cl)\allowbreak{}cc3)\allowbreak{}c3c(\allowbreak{}c2c1)\allowbreak{}NCCC3} & 7.85 & +2 & test & control \\
CHEMBL3323053 & \texttt{O=\allowbreak{}c1c2cc(\allowbreak{}OCCCN3CCCCC3)\allowbreak{}ccc2nc2n1CCC2} & \texttt{O=\allowbreak{}c1c2cc(\allowbreak{}OCCC[NH+]\allowbreak{}3CCCCC3)\allowbreak{}ccc2nc2n1CCC2} & 5.96 & +1 & fold 4 & control \\
CHEMBL3323056 & \texttt{CCN(\allowbreak{}CC)\allowbreak{}CCCOc1ccc2nc3n(\allowbreak{}c(\allowbreak{}=\allowbreak{}O)\allowbreak{}c2c1)\allowbreak{}CCC3} & \texttt{CC[NH+]\allowbreak{}(\allowbreak{}CC)\allowbreak{}CCCOc1ccc2nc3n(\allowbreak{}c(\allowbreak{}=\allowbreak{}O)\allowbreak{}c2c1)\allowbreak{}CCC3} & 5.92 & +1 & test & control \\
CHEMBL3323062 & \texttt{CCN(\allowbreak{}CC)\allowbreak{}CCCOc1ccc2nc3n(\allowbreak{}c(\allowbreak{}=\allowbreak{}O)\allowbreak{}c2c1)\allowbreak{}CCCC3} & \texttt{CC[NH+]\allowbreak{}(\allowbreak{}CC)\allowbreak{}CCCOc1ccc2nc3n(\allowbreak{}c(\allowbreak{}=\allowbreak{}O)\allowbreak{}c2c1)\allowbreak{}CCCC3} & 5.82 & +1 & test & control \\
CHEMBL3323065 & \texttt{O=\allowbreak{}c1c2cc(\allowbreak{}OCCCN3CCCCC3)\allowbreak{}ccc2nc2n1CCCCC2} & \texttt{O=\allowbreak{}c1c2cc(\allowbreak{}OCCC[NH+]\allowbreak{}3CCCCC3)\allowbreak{}ccc2nc2n1CCCCC2} & 6.09 & +1 & fold 4 & control \\
CHEMBL3335030 & \texttt{COc1cccc(\allowbreak{}CN2CCN(\allowbreak{}CCOc3ccc4c5c(\allowbreak{}c(\allowbreak{}=\allowbreak{}O)\allowbreak{}oc4c3)\allowbreak{}CCCC5)\allowbreak{}CC2)\allowbreak{}c1} & \texttt{COc1cccc(\allowbreak{}C[NH+]\allowbreak{}2CC[NH+]\allowbreak{}(\allowbreak{}CCOc3ccc4c5c(\allowbreak{}c(\allowbreak{}=\allowbreak{}O)\allowbreak{}oc4c3)\allowbreak{}CCCC5)\allowbreak{}CC2)\allowbreak{}c1} & 5.14 & +2 & fold 4 & control \\
CHEMBL3335049 & \texttt{COc1cccc(\allowbreak{}CN2CCN(\allowbreak{}CCCOc3ccc4c5c(\allowbreak{}c(\allowbreak{}=\allowbreak{}O)\allowbreak{}oc4c3)\allowbreak{}CCCC5)\allowbreak{}CC2)\allowbreak{}c1} & \texttt{COc1cccc(\allowbreak{}C[NH+]\allowbreak{}2CC[NH+]\allowbreak{}(\allowbreak{}CCCOc3ccc4c5c(\allowbreak{}c(\allowbreak{}=\allowbreak{}O)\allowbreak{}oc4c3)\allowbreak{}CCCC5)\allowbreak{}CC2)\allowbreak{}c1} & 5.21 & +2 & fold 4 & control \\
CHEMBL3343931 & \texttt{N[C@@H]\allowbreak{}(\allowbreak{}Cc1c[nH]\allowbreak{}c2ccccc12)\allowbreak{}C(\allowbreak{}=\allowbreak{}O)\allowbreak{}NCCCCCCNc1c2c(\allowbreak{}nc3ccccc13)\allowbreak{}CCCC2} & \texttt{[NH3+]\allowbreak{}[C@@H]\allowbreak{}(\allowbreak{}Cc1c[nH]\allowbreak{}c2ccccc12)\allowbreak{}C(\allowbreak{}=\allowbreak{}O)\allowbreak{}NCCCCCCNc1c2c(\allowbreak{}[nH+]\allowbreak{}c3ccccc13)\allowbreak{}CCCC2} & 6.97 & +2 & fold 0 & control \\
CHEMBL3356538 & \texttt{COc1ccc2nc3c(\allowbreak{}c(\allowbreak{}NCCNC4=\allowbreak{}CC(\allowbreak{}=\allowbreak{}O)\allowbreak{}c5c(\allowbreak{}O)\allowbreak{}cccc5C4=\allowbreak{}O)\allowbreak{}c2c1)\allowbreak{}CCCC3} & \texttt{COc1ccc2[nH+]\allowbreak{}c3c(\allowbreak{}c(\allowbreak{}NCC[NH2+]\allowbreak{}C4=\allowbreak{}CC(\allowbreak{}=\allowbreak{}O)\allowbreak{}c5c(\allowbreak{}O)\allowbreak{}cccc5C4=\allowbreak{}O)\allowbreak{}c2c1)\allowbreak{}CCCC3} & 6.19 & +2 & test & control \\
CHEMBL3356539 & \texttt{COc1ccc2nc3c(\allowbreak{}c(\allowbreak{}NCCCNC4=\allowbreak{}CC(\allowbreak{}=\allowbreak{}O)\allowbreak{}c5c(\allowbreak{}O)\allowbreak{}cccc5C4=\allowbreak{}O)\allowbreak{}c2c1)\allowbreak{}CCCC3} & \texttt{COc1ccc2[nH+]\allowbreak{}c3c(\allowbreak{}c(\allowbreak{}NCCC[NH2+]\allowbreak{}C4=\allowbreak{}CC(\allowbreak{}=\allowbreak{}O)\allowbreak{}c5c(\allowbreak{}O)\allowbreak{}cccc5C4=\allowbreak{}O)\allowbreak{}c2c1)\allowbreak{}CCCC3} & 6.46 & +2 & test & control \\
CHEMBL3402709 & \texttt{CN(\allowbreak{}CCCCCCCN1C(\allowbreak{}=\allowbreak{}O)\allowbreak{}c2ccccc2C1=\allowbreak{}O)\allowbreak{}Cc1ccccc1} & \texttt{C[NH+]\allowbreak{}(\allowbreak{}CCCCCCCN1C(\allowbreak{}=\allowbreak{}O)\allowbreak{}c2ccccc2C1=\allowbreak{}O)\allowbreak{}Cc1ccccc1} & 6.00 & +1 & fold 4 & control \\
CHEMBL3402710 & \texttt{CN(\allowbreak{}CCCCCCCCN1C(\allowbreak{}=\allowbreak{}O)\allowbreak{}c2ccccc2C1=\allowbreak{}O)\allowbreak{}Cc1ccccc1} & \texttt{C[NH+]\allowbreak{}(\allowbreak{}CCCCCCCCN1C(\allowbreak{}=\allowbreak{}O)\allowbreak{}c2ccccc2C1=\allowbreak{}O)\allowbreak{}Cc1ccccc1} & 6.07 & +1 & fold 4 & control \\
CHEMBL340427 & \texttt{CNC(\allowbreak{}=\allowbreak{}O)\allowbreak{}Oc1cccc(\allowbreak{}CN(\allowbreak{}C)\allowbreak{}CCCOc2ccc3c(\allowbreak{}=\allowbreak{}O)\allowbreak{}c4ccccc4oc3c2)\allowbreak{}c1} & \texttt{CNC(\allowbreak{}=\allowbreak{}O)\allowbreak{}Oc1cccc(\allowbreak{}C[NH+]\allowbreak{}(\allowbreak{}C)\allowbreak{}CCCOc2ccc3c(\allowbreak{}=\allowbreak{}O)\allowbreak{}c4ccccc4oc3c2)\allowbreak{}c1} & 9.52 & +1 & fold 3 & cliff \\
CHEMBL3410953 & \texttt{Cc1c(\allowbreak{}C)\allowbreak{}c2c(\allowbreak{}c(\allowbreak{}C)\allowbreak{}c1O)\allowbreak{}CCC(\allowbreak{}C)\allowbreak{}(\allowbreak{}C(\allowbreak{}=\allowbreak{}O)\allowbreak{}NCCCCCNc1c3c(\allowbreak{}nc4ccccc14)\allowbreak{}CCCC3)\allowbreak{}O2} & \texttt{Cc1c(\allowbreak{}C)\allowbreak{}c2c(\allowbreak{}c(\allowbreak{}C)\allowbreak{}c1O)\allowbreak{}CCC(\allowbreak{}C)\allowbreak{}(\allowbreak{}C(\allowbreak{}=\allowbreak{}O)\allowbreak{}NCCCCCNc1c3c(\allowbreak{}[nH+]\allowbreak{}c4ccccc14)\allowbreak{}CCCC3)\allowbreak{}O2} & 6.44 & +1 & fold 0 & control \\
CHEMBL3410955 & \texttt{Cc1c(\allowbreak{}C)\allowbreak{}c2c(\allowbreak{}c(\allowbreak{}C)\allowbreak{}c1O)\allowbreak{}CCC(\allowbreak{}C)\allowbreak{}(\allowbreak{}C(\allowbreak{}=\allowbreak{}O)\allowbreak{}NCCCCCCCNc1c3c(\allowbreak{}nc4ccccc14)\allowbreak{}CCCC3)\allowbreak{}O2} & \texttt{Cc1c(\allowbreak{}C)\allowbreak{}c2c(\allowbreak{}c(\allowbreak{}C)\allowbreak{}c1O)\allowbreak{}CCC(\allowbreak{}C)\allowbreak{}(\allowbreak{}C(\allowbreak{}=\allowbreak{}O)\allowbreak{}NCCCCCCCNc1c3c(\allowbreak{}[nH+]\allowbreak{}c4ccccc14)\allowbreak{}CCCC3)\allowbreak{}O2} & 6.50 & +1 & fold 0 & control \\
CHEMBL3410957 & \texttt{Cc1c(\allowbreak{}C)\allowbreak{}c2c(\allowbreak{}c(\allowbreak{}C)\allowbreak{}c1O)\allowbreak{}CCC(\allowbreak{}C)\allowbreak{}(\allowbreak{}C(\allowbreak{}=\allowbreak{}O)\allowbreak{}NCCCCCCCCCCNc1c3c(\allowbreak{}nc4ccccc14)\allowbreak{}CCCC3)\allowbreak{}O2} & \texttt{Cc1c(\allowbreak{}C)\allowbreak{}c2c(\allowbreak{}c(\allowbreak{}C)\allowbreak{}c1O)\allowbreak{}CCC(\allowbreak{}C)\allowbreak{}(\allowbreak{}C(\allowbreak{}=\allowbreak{}O)\allowbreak{}NCCCCCCCCCCNc1c3c(\allowbreak{}[nH+]\allowbreak{}c4ccccc14)\allowbreak{}CCCC3)\allowbreak{}O2} & 6.34 & +1 & fold 0 & control \\
CHEMBL3415558 & \texttt{O=\allowbreak{}C1c2ccccc2C(\allowbreak{}=\allowbreak{}O)\allowbreak{}N1CCCCCNCc1ccccc1} & \texttt{O=\allowbreak{}C1c2ccccc2C(\allowbreak{}=\allowbreak{}O)\allowbreak{}N1CCCCC[NH2+]\allowbreak{}Cc1ccccc1} & 6.70 & +1 & test & control \\
CHEMBL3415560 & \texttt{O=\allowbreak{}C1c2ccccc2C(\allowbreak{}=\allowbreak{}O)\allowbreak{}N1CCCCCCCNCc1ccccc1} & \texttt{O=\allowbreak{}C1c2ccccc2C(\allowbreak{}=\allowbreak{}O)\allowbreak{}N1CCCCCCC[NH2+]\allowbreak{}Cc1ccccc1} & 6.52 & +1 & test & control \\
CHEMBL3415568 & \texttt{O=\allowbreak{}C1c2ccccc2C(\allowbreak{}=\allowbreak{}O)\allowbreak{}N1CCCCCNCc1ccc(\allowbreak{}F)\allowbreak{}cc1} & \texttt{O=\allowbreak{}C1c2ccccc2C(\allowbreak{}=\allowbreak{}O)\allowbreak{}N1CCCCC[NH2+]\allowbreak{}Cc1ccc(\allowbreak{}F)\allowbreak{}cc1} & 6.57 & +1 & fold 2 & control \\
CHEMBL3415570 & \texttt{O=\allowbreak{}C1c2ccccc2C(\allowbreak{}=\allowbreak{}O)\allowbreak{}N1CCCCCCCNCc1ccc(\allowbreak{}F)\allowbreak{}cc1} & \texttt{O=\allowbreak{}C1c2ccccc2C(\allowbreak{}=\allowbreak{}O)\allowbreak{}N1CCCCCCC[NH2+]\allowbreak{}Cc1ccc(\allowbreak{}F)\allowbreak{}cc1} & 6.36 & +1 & fold 2 & control \\
CHEMBL3415575 & \texttt{O=\allowbreak{}C1c2ccccc2C(\allowbreak{}=\allowbreak{}O)\allowbreak{}N1CCCCCCCNCc1ccc(\allowbreak{}F)\allowbreak{}cc1F} & \texttt{O=\allowbreak{}C1c2ccccc2C(\allowbreak{}=\allowbreak{}O)\allowbreak{}N1CCCCCCC[NH2+]\allowbreak{}Cc1ccc(\allowbreak{}F)\allowbreak{}cc1F} & 6.42 & +1 & fold 1 & control \\
CHEMBL3415576 & \texttt{O=\allowbreak{}C1c2ccccc2C(\allowbreak{}=\allowbreak{}O)\allowbreak{}N1CCCCCCCCNCc1ccc(\allowbreak{}F)\allowbreak{}cc1F} & \texttt{O=\allowbreak{}C1c2ccccc2C(\allowbreak{}=\allowbreak{}O)\allowbreak{}N1CCCCCCCC[NH2+]\allowbreak{}Cc1ccc(\allowbreak{}F)\allowbreak{}cc1F} & 6.24 & +1 & fold 1 & control \\
CHEMBL3415580 & \texttt{O=\allowbreak{}C1c2ccccc2C(\allowbreak{}=\allowbreak{}O)\allowbreak{}N1CCCCCCCNCc1cccc(\allowbreak{}Cl)\allowbreak{}c1} & \texttt{O=\allowbreak{}C1c2ccccc2C(\allowbreak{}=\allowbreak{}O)\allowbreak{}N1CCCCCCC[NH2+]\allowbreak{}Cc1cccc(\allowbreak{}Cl)\allowbreak{}c1} & 6.19 & +1 & fold 3 & control \\
CHEMBL3415581 & \texttt{O=\allowbreak{}C1c2ccccc2C(\allowbreak{}=\allowbreak{}O)\allowbreak{}N1CCCCCCCCNCc1cccc(\allowbreak{}Cl)\allowbreak{}c1} & \texttt{O=\allowbreak{}C1c2ccccc2C(\allowbreak{}=\allowbreak{}O)\allowbreak{}N1CCCCCCCC[NH2+]\allowbreak{}Cc1cccc(\allowbreak{}Cl)\allowbreak{}c1} & 6.16 & +1 & fold 3 & control \\
CHEMBL3416657 & \texttt{COc1cc(\allowbreak{}/\allowbreak{}C=\allowbreak{}C2\textbackslash{}Oc3cc(\allowbreak{}OCCN4CCCC4)\allowbreak{}ccc3C2=\allowbreak{}O)\allowbreak{}cc(\allowbreak{}OC)\allowbreak{}c1OC} & \texttt{COc1cc(\allowbreak{}/\allowbreak{}C=\allowbreak{}C2\textbackslash{}Oc3cc(\allowbreak{}OCC[NH+]\allowbreak{}4CCCC4)\allowbreak{}ccc3C2=\allowbreak{}O)\allowbreak{}cc(\allowbreak{}OC)\allowbreak{}c1OC} & 6.19 & +1 & fold 2 & control \\
CHEMBL3416658 & \texttt{COc1cc(\allowbreak{}/\allowbreak{}C=\allowbreak{}C2\textbackslash{}Oc3cc(\allowbreak{}OCCN4CCCCC4)\allowbreak{}ccc3C2=\allowbreak{}O)\allowbreak{}cc(\allowbreak{}OC)\allowbreak{}c1OC} & \texttt{COc1cc(\allowbreak{}/\allowbreak{}C=\allowbreak{}C2\textbackslash{}Oc3cc(\allowbreak{}OCC[NH+]\allowbreak{}4CCCCC4)\allowbreak{}ccc3C2=\allowbreak{}O)\allowbreak{}cc(\allowbreak{}OC)\allowbreak{}c1OC} & 6.09 & +1 & fold 2 & control \\
CHEMBL3585374 & \texttt{Cc1ccc2[nH]\allowbreak{}c3c(\allowbreak{}c2c1)\allowbreak{}CN(\allowbreak{}CCCCCCN1CCc2[nH]\allowbreak{}c4ccc(\allowbreak{}C)\allowbreak{}cc4c2C1)\allowbreak{}CC3} & \texttt{Cc1ccc2[nH]\allowbreak{}c3c(\allowbreak{}c2c1)\allowbreak{}C[NH+]\allowbreak{}(\allowbreak{}CCCCCC[NH+]\allowbreak{}1CCc2[nH]\allowbreak{}c4ccc(\allowbreak{}C)\allowbreak{}cc4c2C1)\allowbreak{}CC3} & 6.72 & +2 & fold 4 & control \\
CHEMBL3600551 & \texttt{CC1=\allowbreak{}CC2Cc3nc4cc(\allowbreak{}Cl)\allowbreak{}ccc4c(\allowbreak{}NCCCCCCCCCCCCNc4c5c(\allowbreak{}nc6ccccc46)\allowbreak{}CCCC5)\allowbreak{}c3C(\allowbreak{}C1)\allowbreak{}C2} & \texttt{CC1=\allowbreak{}CC2Cc3[nH+]\allowbreak{}c4cc(\allowbreak{}Cl)\allowbreak{}ccc4c(\allowbreak{}NCCCCCCCCCCCCNc4c5c(\allowbreak{}[nH+]\allowbreak{}c6ccccc46)\allowbreak{}CCCC5)\allowbreak{}c3C(\allowbreak{}C1)\allowbreak{}C2} & 8.72 & +2 & fold 4 & control \\
CHEMBL3600556 & \texttt{CC1=\allowbreak{}CC2Cc3nc4ccccc4c(\allowbreak{}NCCCCCCCCCCCCNc4c5c(\allowbreak{}nc6cc(\allowbreak{}Cl)\allowbreak{}ccc46)\allowbreak{}CCCC5)\allowbreak{}c3C(\allowbreak{}C1)\allowbreak{}C2} & \texttt{CC1=\allowbreak{}CC2Cc3[nH+]\allowbreak{}c4ccccc4c(\allowbreak{}NCCCCCCCCCCCCNc4c5c(\allowbreak{}[nH+]\allowbreak{}c6cc(\allowbreak{}Cl)\allowbreak{}ccc46)\allowbreak{}CCCC5)\allowbreak{}c3C(\allowbreak{}C1)\allowbreak{}C2} & 8.44 & +2 & fold 4 & control \\
CHEMBL3632626 & \texttt{COc1ccc2nc3c(\allowbreak{}c(\allowbreak{}NCCCCNC(\allowbreak{}=\allowbreak{}O)\allowbreak{}C4(\allowbreak{}C)\allowbreak{}CCc5c(\allowbreak{}C)\allowbreak{}c(\allowbreak{}O)\allowbreak{}c(\allowbreak{}C)\allowbreak{}c(\allowbreak{}C)\allowbreak{}c5O4)\allowbreak{}c2c1)\allowbreak{}CCCC3} & \texttt{COc1ccc2[nH+]\allowbreak{}c3c(\allowbreak{}c(\allowbreak{}NCCCCNC(\allowbreak{}=\allowbreak{}O)\allowbreak{}C4(\allowbreak{}C)\allowbreak{}CCc5c(\allowbreak{}C)\allowbreak{}c(\allowbreak{}O)\allowbreak{}c(\allowbreak{}C)\allowbreak{}c(\allowbreak{}C)\allowbreak{}c5O4)\allowbreak{}c2c1)\allowbreak{}CCCC3} & 5.14 & +1 & fold 0 & control \\
CHEMBL3632848 & \texttt{COc1ccc2nc3c(\allowbreak{}c(\allowbreak{}NCCCCCNC(\allowbreak{}=\allowbreak{}O)\allowbreak{}C4(\allowbreak{}C)\allowbreak{}CCc5c(\allowbreak{}C)\allowbreak{}c(\allowbreak{}O)\allowbreak{}c(\allowbreak{}C)\allowbreak{}c(\allowbreak{}C)\allowbreak{}c5O4)\allowbreak{}c2c1)\allowbreak{}CCCC3} & \texttt{COc1ccc2[nH+]\allowbreak{}c3c(\allowbreak{}c(\allowbreak{}NCCCCCNC(\allowbreak{}=\allowbreak{}O)\allowbreak{}C4(\allowbreak{}C)\allowbreak{}CCc5c(\allowbreak{}C)\allowbreak{}c(\allowbreak{}O)\allowbreak{}c(\allowbreak{}C)\allowbreak{}c(\allowbreak{}C)\allowbreak{}c5O4)\allowbreak{}c2c1)\allowbreak{}CCCC3} & 4.88 & +1 & fold 2 & control \\
CHEMBL3632849 & \texttt{COc1ccc2nc3c(\allowbreak{}c(\allowbreak{}NCCCCCCNC(\allowbreak{}=\allowbreak{}O)\allowbreak{}C4(\allowbreak{}C)\allowbreak{}CCc5c(\allowbreak{}C)\allowbreak{}c(\allowbreak{}O)\allowbreak{}c(\allowbreak{}C)\allowbreak{}c(\allowbreak{}C)\allowbreak{}c5O4)\allowbreak{}c2c1)\allowbreak{}CCCC3} & \texttt{COc1ccc2[nH+]\allowbreak{}c3c(\allowbreak{}c(\allowbreak{}NCCCCCCNC(\allowbreak{}=\allowbreak{}O)\allowbreak{}C4(\allowbreak{}C)\allowbreak{}CCc5c(\allowbreak{}C)\allowbreak{}c(\allowbreak{}O)\allowbreak{}c(\allowbreak{}C)\allowbreak{}c(\allowbreak{}C)\allowbreak{}c5O4)\allowbreak{}c2c1)\allowbreak{}CCCC3} & 5.18 & +1 & fold 0 & control \\
CHEMBL3632850 & \texttt{COc1ccc2nc3c(\allowbreak{}c(\allowbreak{}NCCCCCCCNC(\allowbreak{}=\allowbreak{}O)\allowbreak{}C4(\allowbreak{}C)\allowbreak{}CCc5c(\allowbreak{}C)\allowbreak{}c(\allowbreak{}O)\allowbreak{}c(\allowbreak{}C)\allowbreak{}c(\allowbreak{}C)\allowbreak{}c5O4)\allowbreak{}c2c1)\allowbreak{}CCCC3} & \texttt{COc1ccc2[nH+]\allowbreak{}c3c(\allowbreak{}c(\allowbreak{}NCCCCCCCNC(\allowbreak{}=\allowbreak{}O)\allowbreak{}C4(\allowbreak{}C)\allowbreak{}CCc5c(\allowbreak{}C)\allowbreak{}c(\allowbreak{}O)\allowbreak{}c(\allowbreak{}C)\allowbreak{}c(\allowbreak{}C)\allowbreak{}c5O4)\allowbreak{}c2c1)\allowbreak{}CCCC3} & 5.17 & +1 & fold 2 & control \\
CHEMBL3632987 & \texttt{Cc1c(\allowbreak{}C)\allowbreak{}c2c(\allowbreak{}c(\allowbreak{}C)\allowbreak{}c1O)\allowbreak{}CCC(\allowbreak{}C)\allowbreak{}(\allowbreak{}C(\allowbreak{}=\allowbreak{}O)\allowbreak{}NCCCNc1c3c(\allowbreak{}nc4ccccc14)\allowbreak{}CCCC3)\allowbreak{}O2} & \texttt{Cc1c(\allowbreak{}C)\allowbreak{}c2c(\allowbreak{}c(\allowbreak{}C)\allowbreak{}c1O)\allowbreak{}CCC(\allowbreak{}C)\allowbreak{}(\allowbreak{}C(\allowbreak{}=\allowbreak{}O)\allowbreak{}NCCCNc1c3c(\allowbreak{}[nH+]\allowbreak{}c4ccccc14)\allowbreak{}CCCC3)\allowbreak{}O2} & 6.58 & +1 & fold 0 & control \\
CHEMBL3632990 & \texttt{Cc1c(\allowbreak{}C)\allowbreak{}c2c(\allowbreak{}c(\allowbreak{}C)\allowbreak{}c1O)\allowbreak{}CCC(\allowbreak{}C)\allowbreak{}(\allowbreak{}C(\allowbreak{}=\allowbreak{}O)\allowbreak{}NCCCCNc1c3c(\allowbreak{}nc4cc(\allowbreak{}Cl)\allowbreak{}ccc14)\allowbreak{}CCCC3)\allowbreak{}O2} & \texttt{Cc1c(\allowbreak{}C)\allowbreak{}c2c(\allowbreak{}c(\allowbreak{}C)\allowbreak{}c1O)\allowbreak{}CCC(\allowbreak{}C)\allowbreak{}(\allowbreak{}C(\allowbreak{}=\allowbreak{}O)\allowbreak{}NCCCCNc1c3c(\allowbreak{}[nH+]\allowbreak{}c4cc(\allowbreak{}Cl)\allowbreak{}ccc14)\allowbreak{}CCCC3)\allowbreak{}O2} & 6.82 & +1 & fold 1 & control \\
CHEMBL3632991 & \texttt{Cc1c(\allowbreak{}C)\allowbreak{}c2c(\allowbreak{}c(\allowbreak{}C)\allowbreak{}c1O)\allowbreak{}CCC(\allowbreak{}C)\allowbreak{}(\allowbreak{}C(\allowbreak{}=\allowbreak{}O)\allowbreak{}NCCCCCNc1c3c(\allowbreak{}nc4cc(\allowbreak{}Cl)\allowbreak{}ccc14)\allowbreak{}CCCC3)\allowbreak{}O2} & \texttt{Cc1c(\allowbreak{}C)\allowbreak{}c2c(\allowbreak{}c(\allowbreak{}C)\allowbreak{}c1O)\allowbreak{}CCC(\allowbreak{}C)\allowbreak{}(\allowbreak{}C(\allowbreak{}=\allowbreak{}O)\allowbreak{}NCCCCCNc1c3c(\allowbreak{}[nH+]\allowbreak{}c4cc(\allowbreak{}Cl)\allowbreak{}ccc14)\allowbreak{}CCCC3)\allowbreak{}O2} & 6.92 & +1 & fold 1 & control \\
CHEMBL3642041 & \texttt{CC(\allowbreak{}=\allowbreak{}O)\allowbreak{}NCCc1c[nH]\allowbreak{}c2ccc(\allowbreak{}OC(\allowbreak{}=\allowbreak{}O)\allowbreak{}NCCNc3c4c(\allowbreak{}nc5ccccc35)\allowbreak{}CCCC4)\allowbreak{}cc12} & \texttt{CC(\allowbreak{}=\allowbreak{}O)\allowbreak{}NCCc1c[nH]\allowbreak{}c2ccc(\allowbreak{}OC(\allowbreak{}=\allowbreak{}O)\allowbreak{}NCCNc3c4c(\allowbreak{}[nH+]\allowbreak{}c5ccccc35)\allowbreak{}CCCC4)\allowbreak{}cc12} & 5.96 & +1 & fold 3 & cliff \\
CHEMBL3642042 & \texttt{CC(\allowbreak{}=\allowbreak{}O)\allowbreak{}NCCc1c[nH]\allowbreak{}c2ccc(\allowbreak{}OC(\allowbreak{}=\allowbreak{}O)\allowbreak{}NCCCCCCNc3c4c(\allowbreak{}nc5ccccc35)\allowbreak{}CCCC4)\allowbreak{}cc12} & \texttt{CC(\allowbreak{}=\allowbreak{}O)\allowbreak{}NCCc1c[nH]\allowbreak{}c2ccc(\allowbreak{}OC(\allowbreak{}=\allowbreak{}O)\allowbreak{}NCCCCCCNc3c4c(\allowbreak{}[nH+]\allowbreak{}c5ccccc35)\allowbreak{}CCCC4)\allowbreak{}cc12} & 6.75 & +1 & fold 3 & cliff \\
CHEMBL3642044 & \texttt{CC(\allowbreak{}=\allowbreak{}O)\allowbreak{}NCCc1c[nH]\allowbreak{}c2ccc(\allowbreak{}OC(\allowbreak{}=\allowbreak{}O)\allowbreak{}NCCNc3c4c(\allowbreak{}nc5cc(\allowbreak{}Cl)\allowbreak{}ccc35)\allowbreak{}CCCC4)\allowbreak{}cc12} & \texttt{CC(\allowbreak{}=\allowbreak{}O)\allowbreak{}NCCc1c[nH]\allowbreak{}c2ccc(\allowbreak{}OC(\allowbreak{}=\allowbreak{}O)\allowbreak{}NCCNc3c4c(\allowbreak{}[nH+]\allowbreak{}c5cc(\allowbreak{}Cl)\allowbreak{}ccc35)\allowbreak{}CCCC4)\allowbreak{}cc12} & 7.70 & +1 & fold 3 & control \\
CHEMBL3642046 & \texttt{CC(\allowbreak{}=\allowbreak{}O)\allowbreak{}NCCc1c[nH]\allowbreak{}c2ccc(\allowbreak{}OC(\allowbreak{}=\allowbreak{}O)\allowbreak{}NCCCNc3c4c(\allowbreak{}nc5ccccc35)\allowbreak{}CCCC4)\allowbreak{}cc12} & \texttt{CC(\allowbreak{}=\allowbreak{}O)\allowbreak{}NCCc1c[nH]\allowbreak{}c2ccc(\allowbreak{}OC(\allowbreak{}=\allowbreak{}O)\allowbreak{}NCCCNc3c4c(\allowbreak{}[nH+]\allowbreak{}c5ccccc35)\allowbreak{}CCCC4)\allowbreak{}cc12} & 6.57 & +1 & fold 3 & cliff \\
CHEMBL3642047 & \texttt{CC(\allowbreak{}=\allowbreak{}O)\allowbreak{}NCCc1c[nH]\allowbreak{}c2ccc(\allowbreak{}OC(\allowbreak{}=\allowbreak{}O)\allowbreak{}NCCCCNc3c4c(\allowbreak{}nc5ccccc35)\allowbreak{}CCCC4)\allowbreak{}cc12} & \texttt{CC(\allowbreak{}=\allowbreak{}O)\allowbreak{}NCCc1c[nH]\allowbreak{}c2ccc(\allowbreak{}OC(\allowbreak{}=\allowbreak{}O)\allowbreak{}NCCCCNc3c4c(\allowbreak{}[nH+]\allowbreak{}c5ccccc35)\allowbreak{}CCCC4)\allowbreak{}cc12} & 6.42 & +1 & fold 3 & cliff \\
CHEMBL3642048 & \texttt{CC(\allowbreak{}=\allowbreak{}O)\allowbreak{}NCCc1c[nH]\allowbreak{}c2ccc(\allowbreak{}OC(\allowbreak{}=\allowbreak{}O)\allowbreak{}NCCCCCNc3c4c(\allowbreak{}nc5ccccc35)\allowbreak{}CCCC4)\allowbreak{}cc12} & \texttt{CC(\allowbreak{}=\allowbreak{}O)\allowbreak{}NCCc1c[nH]\allowbreak{}c2ccc(\allowbreak{}OC(\allowbreak{}=\allowbreak{}O)\allowbreak{}NCCCCCNc3c4c(\allowbreak{}[nH+]\allowbreak{}c5ccccc35)\allowbreak{}CCCC4)\allowbreak{}cc12} & 6.73 & +1 & fold 3 & cliff \\
CHEMBL3642049 & \texttt{CC(\allowbreak{}=\allowbreak{}O)\allowbreak{}NCCc1c[nH]\allowbreak{}c2ccc(\allowbreak{}OC(\allowbreak{}=\allowbreak{}O)\allowbreak{}NCCCCCCCCNc3c4c(\allowbreak{}nc5ccccc35)\allowbreak{}CCCC4)\allowbreak{}cc12} & \texttt{CC(\allowbreak{}=\allowbreak{}O)\allowbreak{}NCCc1c[nH]\allowbreak{}c2ccc(\allowbreak{}OC(\allowbreak{}=\allowbreak{}O)\allowbreak{}NCCCCCCCCNc3c4c(\allowbreak{}[nH+]\allowbreak{}c5ccccc35)\allowbreak{}CCCC4)\allowbreak{}cc12} & 8.34 & +1 & fold 3 & cliff \\
CHEMBL3642051 & \texttt{CC(\allowbreak{}=\allowbreak{}O)\allowbreak{}NCCc1c[nH]\allowbreak{}c2ccc(\allowbreak{}OC(\allowbreak{}=\allowbreak{}O)\allowbreak{}NCCNc3c4c(\allowbreak{}nc5ccc(\allowbreak{}Cl)\allowbreak{}cc35)\allowbreak{}CCCC4)\allowbreak{}cc12} & \texttt{CC(\allowbreak{}=\allowbreak{}O)\allowbreak{}NCCc1c[nH]\allowbreak{}c2ccc(\allowbreak{}OC(\allowbreak{}=\allowbreak{}O)\allowbreak{}NCCNc3c4c(\allowbreak{}[nH+]\allowbreak{}c5ccc(\allowbreak{}Cl)\allowbreak{}cc35)\allowbreak{}CCCC4)\allowbreak{}cc12} & 5.48 & +1 & test & cliff \\
CHEMBL3642052 & \texttt{CC(\allowbreak{}=\allowbreak{}O)\allowbreak{}NCCc1c[nH]\allowbreak{}c2ccc(\allowbreak{}OC(\allowbreak{}=\allowbreak{}O)\allowbreak{}NCCCCCCCCCCNc3c4c(\allowbreak{}nc5ccccc35)\allowbreak{}CCCC4)\allowbreak{}cc12} & \texttt{CC(\allowbreak{}=\allowbreak{}O)\allowbreak{}NCCc1c[nH]\allowbreak{}c2ccc(\allowbreak{}OC(\allowbreak{}=\allowbreak{}O)\allowbreak{}NCCCCCCCCCCNc3c4c(\allowbreak{}[nH+]\allowbreak{}c5ccccc35)\allowbreak{}CCCC4)\allowbreak{}cc12} & 8.93 & +1 & fold 3 & cliff \\
CHEMBL3642053 & \texttt{CC(\allowbreak{}=\allowbreak{}O)\allowbreak{}NCCc1c[nH]\allowbreak{}c2ccc(\allowbreak{}OC(\allowbreak{}=\allowbreak{}O)\allowbreak{}NCCCCCCCCCCCCNc3c4c(\allowbreak{}nc5ccccc35)\allowbreak{}CCCC4)\allowbreak{}cc12} & \texttt{CC(\allowbreak{}=\allowbreak{}O)\allowbreak{}NCCc1c[nH]\allowbreak{}c2ccc(\allowbreak{}OC(\allowbreak{}=\allowbreak{}O)\allowbreak{}NCCCCCCCCCCCCNc3c4c(\allowbreak{}[nH+]\allowbreak{}c5ccccc35)\allowbreak{}CCCC4)\allowbreak{}cc12} & 8.42 & +1 & fold 3 & cliff \\
CHEMBL3642054 & \texttt{CC(\allowbreak{}=\allowbreak{}O)\allowbreak{}NCCc1c[nH]\allowbreak{}c2ccc(\allowbreak{}OC(\allowbreak{}=\allowbreak{}O)\allowbreak{}NCCCCCCCCCCNc3c4c(\allowbreak{}nc5ccc(\allowbreak{}Cl)\allowbreak{}cc35)\allowbreak{}CCCC4)\allowbreak{}cc12} & \texttt{CC(\allowbreak{}=\allowbreak{}O)\allowbreak{}NCCc1c[nH]\allowbreak{}c2ccc(\allowbreak{}OC(\allowbreak{}=\allowbreak{}O)\allowbreak{}NCCCCCCCCCCNc3c4c(\allowbreak{}[nH+]\allowbreak{}c5ccc(\allowbreak{}Cl)\allowbreak{}cc35)\allowbreak{}CCCC4)\allowbreak{}cc12} & 8.85 & +1 & test & cliff \\
CHEMBL3642056 & \texttt{CC(\allowbreak{}=\allowbreak{}O)\allowbreak{}NCCc1c[nH]\allowbreak{}c2ccc(\allowbreak{}OC(\allowbreak{}=\allowbreak{}O)\allowbreak{}NCCCCCCCCCCCCNc3c4c(\allowbreak{}nc5cc(\allowbreak{}Cl)\allowbreak{}ccc35)\allowbreak{}CCCC4)\allowbreak{}cc12} & \texttt{CC(\allowbreak{}=\allowbreak{}O)\allowbreak{}NCCc1c[nH]\allowbreak{}c2ccc(\allowbreak{}OC(\allowbreak{}=\allowbreak{}O)\allowbreak{}NCCCCCCCCCCCCNc3c4c(\allowbreak{}[nH+]\allowbreak{}c5cc(\allowbreak{}Cl)\allowbreak{}ccc35)\allowbreak{}CCCC4)\allowbreak{}cc12} & 7.59 & +1 & fold 3 & control \\
CHEMBL3642057 & \texttt{CC(\allowbreak{}=\allowbreak{}O)\allowbreak{}NCCc1c[nH]\allowbreak{}c2ccc(\allowbreak{}OC(\allowbreak{}=\allowbreak{}O)\allowbreak{}NCCCCCCCCCCCCNc3c4c(\allowbreak{}nc5ccc(\allowbreak{}Cl)\allowbreak{}cc35)\allowbreak{}CCCC4)\allowbreak{}cc12} & \texttt{CC(\allowbreak{}=\allowbreak{}O)\allowbreak{}NCCc1c[nH]\allowbreak{}c2ccc(\allowbreak{}OC(\allowbreak{}=\allowbreak{}O)\allowbreak{}NCCCCCCCCCCCCNc3c4c(\allowbreak{}[nH+]\allowbreak{}c5ccc(\allowbreak{}Cl)\allowbreak{}cc35)\allowbreak{}CCCC4)\allowbreak{}cc12} & 8.15 & +1 & test & cliff \\
CHEMBL3752119 & \texttt{COc1cccc2cc(\allowbreak{}C(\allowbreak{}=\allowbreak{}O)\allowbreak{}NCCCCCCNc3c4c(\allowbreak{}nc5ccccc35)\allowbreak{}CCCCC4)\allowbreak{}oc12} & \texttt{COc1cccc2cc(\allowbreak{}C(\allowbreak{}=\allowbreak{}O)\allowbreak{}NCCCCCCNc3c4c(\allowbreak{}[nH+]\allowbreak{}c5ccccc35)\allowbreak{}CCCCC4)\allowbreak{}oc12} & 7.49 & +1 & fold 3 & control \\
CHEMBL3752555 & \texttt{O=\allowbreak{}C(\allowbreak{}NCCCCCCNc1c2c(\allowbreak{}nc3ccccc13)\allowbreak{}CCC2)\allowbreak{}c1cc2ccccc2o1} & \texttt{O=\allowbreak{}C(\allowbreak{}NCCCCCCNc1c2c(\allowbreak{}[nH+]\allowbreak{}c3ccccc13)\allowbreak{}CCC2)\allowbreak{}c1cc2ccccc2o1} & 6.26 & +1 & test & control \\
CHEMBL3752926 & \texttt{COc1cccc2cc(\allowbreak{}C(\allowbreak{}=\allowbreak{}O)\allowbreak{}NCCCCCCCNc3c4c(\allowbreak{}nc5ccccc35)\allowbreak{}CCCC4)\allowbreak{}oc12} & \texttt{COc1cccc2cc(\allowbreak{}C(\allowbreak{}=\allowbreak{}O)\allowbreak{}NCCCCCCCNc3c4c(\allowbreak{}[nH+]\allowbreak{}c5ccccc35)\allowbreak{}CCCC4)\allowbreak{}oc12} & 7.36 & +1 & fold 3 & control \\
CHEMBL3754480 & \texttt{O=\allowbreak{}C(\allowbreak{}NCCCCCCCNc1c2c(\allowbreak{}nc3ccccc13)\allowbreak{}CCCCC2)\allowbreak{}c1cc2ccccc2o1} & \texttt{O=\allowbreak{}C(\allowbreak{}NCCCCCCCNc1c2c(\allowbreak{}[nH+]\allowbreak{}c3ccccc13)\allowbreak{}CCCCC2)\allowbreak{}c1cc2ccccc2o1} & 6.49 & +1 & test & control \\
CHEMBL3827011 & \texttt{CCCCCCCCC1OCc2c[n+]\allowbreak{}(\allowbreak{}C)\allowbreak{}c(\allowbreak{}C)\allowbreak{}c(\allowbreak{}OC(\allowbreak{}=\allowbreak{}O)\allowbreak{}N(\allowbreak{}C)\allowbreak{}C)\allowbreak{}c2CO1} & \texttt{CCCCCCCCC1OCc2c[n+]\allowbreak{}(\allowbreak{}C)\allowbreak{}c(\allowbreak{}C)\allowbreak{}c(\allowbreak{}OC(\allowbreak{}=\allowbreak{}O)\allowbreak{}N(\allowbreak{}C)\allowbreak{}C)\allowbreak{}c2CO1} & 5.89 & +1 & fold 2 & control \\
CHEMBL3828028 & \texttt{CCCCCCCC1OCc2c[n+]\allowbreak{}(\allowbreak{}C)\allowbreak{}c(\allowbreak{}C)\allowbreak{}c(\allowbreak{}OC(\allowbreak{}=\allowbreak{}O)\allowbreak{}N(\allowbreak{}C)\allowbreak{}C)\allowbreak{}c2CO1} & \texttt{CCCCCCCC1OCc2c[n+]\allowbreak{}(\allowbreak{}C)\allowbreak{}c(\allowbreak{}C)\allowbreak{}c(\allowbreak{}OC(\allowbreak{}=\allowbreak{}O)\allowbreak{}N(\allowbreak{}C)\allowbreak{}C)\allowbreak{}c2CO1} & 5.68 & +1 & fold 2 & control \\
CHEMBL3828293 & \texttt{CN(\allowbreak{}CCCCCOc1ccc2ccc(\allowbreak{}=\allowbreak{}O)\allowbreak{}oc2c1)\allowbreak{}Cc1cccc(\allowbreak{}OC(\allowbreak{}=\allowbreak{}O)\allowbreak{}NCCCCCc2ccccc2)\allowbreak{}c1} & \texttt{C[NH+]\allowbreak{}(\allowbreak{}CCCCCOc1ccc2ccc(\allowbreak{}=\allowbreak{}O)\allowbreak{}oc2c1)\allowbreak{}Cc1cccc(\allowbreak{}OC(\allowbreak{}=\allowbreak{}O)\allowbreak{}NCCCCCc2ccccc2)\allowbreak{}c1} & 7.05 & +1 & fold 4 & control \\
CHEMBL384886 & \texttt{COc1cc(\allowbreak{}CCC(\allowbreak{}=\allowbreak{}O)\allowbreak{}NCCCC(\allowbreak{}=\allowbreak{}O)\allowbreak{}NNc2c3c(\allowbreak{}nc4ccccc24)\allowbreak{}CCCC3)\allowbreak{}cc(\allowbreak{}OC)\allowbreak{}c1OC} & \texttt{COc1cc(\allowbreak{}CCC(\allowbreak{}=\allowbreak{}O)\allowbreak{}NCCCC(\allowbreak{}=\allowbreak{}O)\allowbreak{}NNc2c3c(\allowbreak{}[nH+]\allowbreak{}c4ccccc24)\allowbreak{}CCCC3)\allowbreak{}cc(\allowbreak{}OC)\allowbreak{}c1OC} & 7.14 & +1 & fold 2 & control \\
CHEMBL3884618 & \texttt{O=\allowbreak{}S(\allowbreak{}=\allowbreak{}O)\allowbreak{}(\allowbreak{}c1ccccc1)\allowbreak{}n1ccc2c(\allowbreak{}N3CCN(\allowbreak{}CCCNc4c5c(\allowbreak{}nc6ccccc46)\allowbreak{}CCCC5)\allowbreak{}CC3)\allowbreak{}cccc21} & \texttt{O=\allowbreak{}S(\allowbreak{}=\allowbreak{}O)\allowbreak{}(\allowbreak{}c1ccccc1)\allowbreak{}n1ccc2c(\allowbreak{}N3CC[NH+]\allowbreak{}(\allowbreak{}CCCNc4c5c(\allowbreak{}[nH+]\allowbreak{}c6ccccc46)\allowbreak{}CCCC5)\allowbreak{}CC3)\allowbreak{}cccc21} & 8.01 & +2 & fold 2 & control \\
CHEMBL3884690 & \texttt{O=\allowbreak{}S(\allowbreak{}=\allowbreak{}O)\allowbreak{}(\allowbreak{}c1ccccc1)\allowbreak{}n1ccc2c(\allowbreak{}N3CCN(\allowbreak{}CCCCCCCNc4c5c(\allowbreak{}nc6ccccc46)\allowbreak{}CCCC5)\allowbreak{}CC3)\allowbreak{}cccc21} & \texttt{O=\allowbreak{}S(\allowbreak{}=\allowbreak{}O)\allowbreak{}(\allowbreak{}c1ccccc1)\allowbreak{}n1ccc2c(\allowbreak{}N3CC[NH+]\allowbreak{}(\allowbreak{}CCCCCCCNc4c5c(\allowbreak{}[nH+]\allowbreak{}c6ccccc46)\allowbreak{}CCCC5)\allowbreak{}CC3)\allowbreak{}cccc21} & 8.21 & +2 & fold 2 & control \\
CHEMBL389283 & \texttt{c1cc[n+]\allowbreak{}(\allowbreak{}CCCCCC[n+]\allowbreak{}2ccccc2)\allowbreak{}cc1} & \texttt{c1cc[n+]\allowbreak{}(\allowbreak{}CCCCCC[n+]\allowbreak{}2ccccc2)\allowbreak{}cc1} & 4.20 & +2 & fold 1 & cliff \\
CHEMBL3901794 & \texttt{Cc1ccc(\allowbreak{}/\allowbreak{}C=\allowbreak{}C2/\allowbreak{}CNCC3=\allowbreak{}C2N=\allowbreak{}C2SC=\allowbreak{}C(\allowbreak{}c4ccc(\allowbreak{}C)\allowbreak{}cc4)\allowbreak{}N2C3c2ccc(\allowbreak{}C)\allowbreak{}cc2)\allowbreak{}cc1} & \texttt{Cc1ccc(\allowbreak{}/\allowbreak{}C=\allowbreak{}C2/\allowbreak{}C[NH2+]\allowbreak{}CC3=\allowbreak{}C2N=\allowbreak{}C2SC=\allowbreak{}C(\allowbreak{}c4ccc(\allowbreak{}C)\allowbreak{}cc4)\allowbreak{}N2C3c2ccc(\allowbreak{}C)\allowbreak{}cc2)\allowbreak{}cc1} & 5.75 & +1 & fold 1 & control \\
CHEMBL390737 & \texttt{CN(\allowbreak{}CCCCCCOc1ccc2c(\allowbreak{}=\allowbreak{}O)\allowbreak{}c3ccccc3oc2c1)\allowbreak{}Cc1ccccc1} & \texttt{C[NH+]\allowbreak{}(\allowbreak{}CCCCCCOc1ccc2c(\allowbreak{}=\allowbreak{}O)\allowbreak{}c3ccccc3oc2c1)\allowbreak{}Cc1ccccc1} & 5.54 & +1 & fold 1 & control \\
CHEMBL3945423 & \texttt{Cc1ccc(\allowbreak{}C2=\allowbreak{}CSC3=\allowbreak{}NC4=\allowbreak{}C(\allowbreak{}CNC/\allowbreak{}C4=\allowbreak{}C/\allowbreak{}c4ccc(\allowbreak{}Cl)\allowbreak{}cc4)\allowbreak{}C(\allowbreak{}c4ccc(\allowbreak{}Cl)\allowbreak{}cc4)\allowbreak{}N23)\allowbreak{}cc1} & \texttt{Cc1ccc(\allowbreak{}C2=\allowbreak{}CSC3=\allowbreak{}NC4=\allowbreak{}C(\allowbreak{}C[NH2+]\allowbreak{}C/\allowbreak{}C4=\allowbreak{}C/\allowbreak{}c4ccc(\allowbreak{}Cl)\allowbreak{}cc4)\allowbreak{}C(\allowbreak{}c4ccc(\allowbreak{}Cl)\allowbreak{}cc4)\allowbreak{}N23)\allowbreak{}cc1} & 5.96 & +1 & fold 1 & control \\
CHEMBL4060617 & \texttt{O=\allowbreak{}C(\allowbreak{}CCCC1CCSS1)\allowbreak{}NCCCNc1c2c(\allowbreak{}nc3cc(\allowbreak{}Cl)\allowbreak{}ccc13)\allowbreak{}CCCC2} & \texttt{O=\allowbreak{}C(\allowbreak{}CCCC1CCSS1)\allowbreak{}NCCCNc1c2c(\allowbreak{}[nH+]\allowbreak{}c3cc(\allowbreak{}Cl)\allowbreak{}ccc13)\allowbreak{}CCCC2} & 6.60 & +1 & fold 3 & cliff \\
CHEMBL4063045 & \texttt{COc1cc(\allowbreak{}/\allowbreak{}C=\allowbreak{}C/\allowbreak{}C(\allowbreak{}=\allowbreak{}O)\allowbreak{}Nc2ccccc2)\allowbreak{}ccc1OCCCCCCN(\allowbreak{}C)\allowbreak{}C} & \texttt{COc1cc(\allowbreak{}/\allowbreak{}C=\allowbreak{}C/\allowbreak{}C(\allowbreak{}=\allowbreak{}O)\allowbreak{}Nc2ccccc2)\allowbreak{}ccc1OCCCCCC[NH+]\allowbreak{}(\allowbreak{}C)\allowbreak{}C} & 5.08 & +1 & fold 3 & control \\
CHEMBL4084714 & \texttt{COc1cc(\allowbreak{}/\allowbreak{}C=\allowbreak{}C/\allowbreak{}C(\allowbreak{}=\allowbreak{}O)\allowbreak{}Nc2ccccc2)\allowbreak{}ccc1OCCCCCCCCN(\allowbreak{}C)\allowbreak{}C} & \texttt{COc1cc(\allowbreak{}/\allowbreak{}C=\allowbreak{}C/\allowbreak{}C(\allowbreak{}=\allowbreak{}O)\allowbreak{}Nc2ccccc2)\allowbreak{}ccc1OCCCCCCCC[NH+]\allowbreak{}(\allowbreak{}C)\allowbreak{}C} & 5.05 & +1 & fold 3 & control \\
CHEMBL4090413 & \texttt{COc1cc(\allowbreak{}/\allowbreak{}C=\allowbreak{}C/\allowbreak{}C(\allowbreak{}=\allowbreak{}O)\allowbreak{}Nc2ccccc2)\allowbreak{}ccc1OCCCCCN1CCCCC1} & \texttt{COc1cc(\allowbreak{}/\allowbreak{}C=\allowbreak{}C/\allowbreak{}C(\allowbreak{}=\allowbreak{}O)\allowbreak{}Nc2ccccc2)\allowbreak{}ccc1OCCCCC[NH+]\allowbreak{}1CCCCC1} & 5.32 & +1 & test & control \\
CHEMBL4094981 & \texttt{CCN(\allowbreak{}CC)\allowbreak{}CCCCCCOc1ccc(\allowbreak{}/\allowbreak{}C=\allowbreak{}C/\allowbreak{}C(\allowbreak{}=\allowbreak{}O)\allowbreak{}Nc2ccccc2)\allowbreak{}cc1OC} & \texttt{CC[NH+]\allowbreak{}(\allowbreak{}CC)\allowbreak{}CCCCCCOc1ccc(\allowbreak{}/\allowbreak{}C=\allowbreak{}C/\allowbreak{}C(\allowbreak{}=\allowbreak{}O)\allowbreak{}Nc2ccccc2)\allowbreak{}cc1OC} & 5.74 & +1 & fold 4 & control \\
CHEMBL4098132 & \texttt{COc1cc(\allowbreak{}/\allowbreak{}C=\allowbreak{}C/\allowbreak{}C(\allowbreak{}=\allowbreak{}O)\allowbreak{}Nc2ccccc2)\allowbreak{}ccc1OCCCCCCCCN1CCCCC1} & \texttt{COc1cc(\allowbreak{}/\allowbreak{}C=\allowbreak{}C/\allowbreak{}C(\allowbreak{}=\allowbreak{}O)\allowbreak{}Nc2ccccc2)\allowbreak{}ccc1OCCCCCCCC[NH+]\allowbreak{}1CCCCC1} & 5.38 & +1 & test & control \\
CHEMBL4105158 & \texttt{CCN(\allowbreak{}CC)\allowbreak{}CCCCCCCCOc1ccc(\allowbreak{}/\allowbreak{}C=\allowbreak{}C/\allowbreak{}C(\allowbreak{}=\allowbreak{}O)\allowbreak{}Nc2ccccc2)\allowbreak{}cc1OC} & \texttt{CC[NH+]\allowbreak{}(\allowbreak{}CC)\allowbreak{}CCCCCCCCOc1ccc(\allowbreak{}/\allowbreak{}C=\allowbreak{}C/\allowbreak{}C(\allowbreak{}=\allowbreak{}O)\allowbreak{}Nc2ccccc2)\allowbreak{}cc1OC} & 5.87 & +1 & fold 4 & control \\
CHEMBL418712 & \texttt{CN(\allowbreak{}C)\allowbreak{}Cc1ccc(\allowbreak{}CSCCCCCCCCCSCc2ccc(\allowbreak{}CN(\allowbreak{}C)\allowbreak{}C)\allowbreak{}o2)\allowbreak{}o1} & \texttt{C[NH+]\allowbreak{}(\allowbreak{}C)\allowbreak{}Cc1ccc(\allowbreak{}CSCCCCCCCCCSCc2ccc(\allowbreak{}C[NH+]\allowbreak{}(\allowbreak{}C)\allowbreak{}C)\allowbreak{}o2)\allowbreak{}o1} & 6.34 & +2 & fold 0 & control \\
CHEMBL4202784 & \texttt{COc1ccc(\allowbreak{}Cn2cc(\allowbreak{}C(\allowbreak{}=\allowbreak{}O)\allowbreak{}NCCCCCNc3c4c(\allowbreak{}nc5ccccc35)\allowbreak{}CCCC4)\allowbreak{}c(\allowbreak{}=\allowbreak{}O)\allowbreak{}c3ccccc32)\allowbreak{}cc1} & \texttt{COc1ccc(\allowbreak{}Cn2cc(\allowbreak{}C(\allowbreak{}=\allowbreak{}O)\allowbreak{}NCCCCCNc3c4c(\allowbreak{}[nH+]\allowbreak{}c5ccccc35)\allowbreak{}CCCC4)\allowbreak{}c(\allowbreak{}=\allowbreak{}O)\allowbreak{}c3ccccc32)\allowbreak{}cc1} & 6.53 & +1 & fold 2 & cliff \\
CHEMBL4204913 & \texttt{COc1ccc(\allowbreak{}Cn2cc(\allowbreak{}C(\allowbreak{}=\allowbreak{}O)\allowbreak{}NCCCCCNc3c4c(\allowbreak{}nc5ccc(\allowbreak{}OC)\allowbreak{}cc35)\allowbreak{}CCCC4)\allowbreak{}c(\allowbreak{}=\allowbreak{}O)\allowbreak{}c3ccccc32)\allowbreak{}cc1} & \texttt{COc1ccc(\allowbreak{}Cn2cc(\allowbreak{}C(\allowbreak{}=\allowbreak{}O)\allowbreak{}NCCCCCNc3c4c(\allowbreak{}[nH+]\allowbreak{}c5ccc(\allowbreak{}OC)\allowbreak{}cc35)\allowbreak{}CCCC4)\allowbreak{}c(\allowbreak{}=\allowbreak{}O)\allowbreak{}c3ccccc32)\allowbreak{}cc1} & 5.10 & +1 & fold 0 & control \\
CHEMBL4205910 & \texttt{COc1ccc(\allowbreak{}Cn2cc(\allowbreak{}C(\allowbreak{}=\allowbreak{}O)\allowbreak{}NCCCCCCNc3c4c(\allowbreak{}nc5ccc(\allowbreak{}OC)\allowbreak{}cc35)\allowbreak{}CCCC4)\allowbreak{}c(\allowbreak{}=\allowbreak{}O)\allowbreak{}c3ccccc32)\allowbreak{}cc1} & \texttt{COc1ccc(\allowbreak{}Cn2cc(\allowbreak{}C(\allowbreak{}=\allowbreak{}O)\allowbreak{}NCCCCCCNc3c4c(\allowbreak{}[nH+]\allowbreak{}c5ccc(\allowbreak{}OC)\allowbreak{}cc35)\allowbreak{}CCCC4)\allowbreak{}c(\allowbreak{}=\allowbreak{}O)\allowbreak{}c3ccccc32)\allowbreak{}cc1} & 5.00 & +1 & fold 0 & control \\
CHEMBL4206484 & \texttt{COc1ccc(\allowbreak{}Cn2cc(\allowbreak{}C(\allowbreak{}=\allowbreak{}O)\allowbreak{}NCCCCCCNc3c4c(\allowbreak{}nc5ccccc35)\allowbreak{}CCCC4)\allowbreak{}c(\allowbreak{}=\allowbreak{}O)\allowbreak{}c3ccccc32)\allowbreak{}cc1} & \texttt{COc1ccc(\allowbreak{}Cn2cc(\allowbreak{}C(\allowbreak{}=\allowbreak{}O)\allowbreak{}NCCCCCCNc3c4c(\allowbreak{}[nH+]\allowbreak{}c5ccccc35)\allowbreak{}CCCC4)\allowbreak{}c(\allowbreak{}=\allowbreak{}O)\allowbreak{}c3ccccc32)\allowbreak{}cc1} & 6.71 & +1 & fold 2 & cliff \\
CHEMBL4208322 & \texttt{CCCCNCC(\allowbreak{}=\allowbreak{}O)\allowbreak{}Nc1ccc2nc3n(\allowbreak{}c(\allowbreak{}=\allowbreak{}O)\allowbreak{}c2c1)\allowbreak{}CCC3} & \texttt{CCCC[NH2+]\allowbreak{}CC(\allowbreak{}=\allowbreak{}O)\allowbreak{}Nc1ccc2nc3n(\allowbreak{}c(\allowbreak{}=\allowbreak{}O)\allowbreak{}c2c1)\allowbreak{}CCC3} & 7.67 & +1 & fold 0 & control \\
CHEMBL4209755 & \texttt{COc1ccc(\allowbreak{}Cn2cc(\allowbreak{}C(\allowbreak{}=\allowbreak{}O)\allowbreak{}NCCCCNc3c4c(\allowbreak{}nc5ccccc35)\allowbreak{}CCCC4)\allowbreak{}c(\allowbreak{}=\allowbreak{}O)\allowbreak{}c3ccccc32)\allowbreak{}cc1} & \texttt{COc1ccc(\allowbreak{}Cn2cc(\allowbreak{}C(\allowbreak{}=\allowbreak{}O)\allowbreak{}NCCCCNc3c4c(\allowbreak{}[nH+]\allowbreak{}c5ccccc35)\allowbreak{}CCCC4)\allowbreak{}c(\allowbreak{}=\allowbreak{}O)\allowbreak{}c3ccccc32)\allowbreak{}cc1} & 6.31 & +1 & fold 2 & cliff \\
CHEMBL4210274 & \texttt{COc1ccc(\allowbreak{}Cn2cc(\allowbreak{}C(\allowbreak{}=\allowbreak{}O)\allowbreak{}NCCCCCCCCNc3c4c(\allowbreak{}nc5ccc(\allowbreak{}OC)\allowbreak{}cc35)\allowbreak{}CCCC4)\allowbreak{}c(\allowbreak{}=\allowbreak{}O)\allowbreak{}c3ccccc32)\allowbreak{}cc1} & \texttt{COc1ccc(\allowbreak{}Cn2cc(\allowbreak{}C(\allowbreak{}=\allowbreak{}O)\allowbreak{}NCCCCCCCCNc3c4c(\allowbreak{}[nH+]\allowbreak{}c5ccc(\allowbreak{}OC)\allowbreak{}cc35)\allowbreak{}CCCC4)\allowbreak{}c(\allowbreak{}=\allowbreak{}O)\allowbreak{}c3ccccc32)\allowbreak{}cc1} & 4.13 & +1 & fold 2 & cliff \\
CHEMBL4212151 & \texttt{COc1ccc(\allowbreak{}Cn2cc(\allowbreak{}C(\allowbreak{}=\allowbreak{}O)\allowbreak{}NCCNc3c4c(\allowbreak{}nc5ccc(\allowbreak{}OC)\allowbreak{}cc35)\allowbreak{}CCCC4)\allowbreak{}c(\allowbreak{}=\allowbreak{}O)\allowbreak{}c3ccccc32)\allowbreak{}cc1} & \texttt{COc1ccc(\allowbreak{}Cn2cc(\allowbreak{}C(\allowbreak{}=\allowbreak{}O)\allowbreak{}NCCNc3c4c(\allowbreak{}[nH+]\allowbreak{}c5ccc(\allowbreak{}OC)\allowbreak{}cc35)\allowbreak{}CCCC4)\allowbreak{}c(\allowbreak{}=\allowbreak{}O)\allowbreak{}c3ccccc32)\allowbreak{}cc1} & 5.58 & +1 & test & control \\
CHEMBL4214900 & \texttt{COc1ccc(\allowbreak{}Cn2cc(\allowbreak{}C(\allowbreak{}=\allowbreak{}O)\allowbreak{}NCCCCCCCNc3c4c(\allowbreak{}nc5ccccc35)\allowbreak{}CCCC4)\allowbreak{}c(\allowbreak{}=\allowbreak{}O)\allowbreak{}c3ccccc32)\allowbreak{}cc1} & \texttt{COc1ccc(\allowbreak{}Cn2cc(\allowbreak{}C(\allowbreak{}=\allowbreak{}O)\allowbreak{}NCCCCCCCNc3c4c(\allowbreak{}[nH+]\allowbreak{}c5ccccc35)\allowbreak{}CCCC4)\allowbreak{}c(\allowbreak{}=\allowbreak{}O)\allowbreak{}c3ccccc32)\allowbreak{}cc1} & 6.71 & +1 & fold 2 & cliff \\
CHEMBL421501 & \texttt{c1ccc2c(\allowbreak{}NCCCCCCNc3c4c(\allowbreak{}nc5ccccc35)\allowbreak{}CCC4)\allowbreak{}c3c(\allowbreak{}nc2c1)\allowbreak{}CCC3} & \texttt{c1ccc2c(\allowbreak{}NCCCCCCNc3c4c(\allowbreak{}[nH+]\allowbreak{}c5ccccc35)\allowbreak{}CCC4)\allowbreak{}c3c(\allowbreak{}[nH+]\allowbreak{}c2c1)\allowbreak{}CCC3} & 6.94 & +2 & fold 0 & cliff \\
CHEMBL4215154 & \texttt{COc1ccc(\allowbreak{}Cn2cc(\allowbreak{}C(\allowbreak{}=\allowbreak{}O)\allowbreak{}NCCCCNc3c4c(\allowbreak{}nc5cc(\allowbreak{}Cl)\allowbreak{}ccc35)\allowbreak{}CCCC4)\allowbreak{}c(\allowbreak{}=\allowbreak{}O)\allowbreak{}c3ccccc32)\allowbreak{}cc1} & \texttt{COc1ccc(\allowbreak{}Cn2cc(\allowbreak{}C(\allowbreak{}=\allowbreak{}O)\allowbreak{}NCCCCNc3c4c(\allowbreak{}[nH+]\allowbreak{}c5cc(\allowbreak{}Cl)\allowbreak{}ccc35)\allowbreak{}CCCC4)\allowbreak{}c(\allowbreak{}=\allowbreak{}O)\allowbreak{}c3ccccc32)\allowbreak{}cc1} & 7.05 & +1 & fold 4 & control \\
CHEMBL4217126 & \texttt{COc1ccc(\allowbreak{}Cn2cc(\allowbreak{}C(\allowbreak{}=\allowbreak{}O)\allowbreak{}NCCCNc3c4c(\allowbreak{}nc5ccc(\allowbreak{}OC)\allowbreak{}cc35)\allowbreak{}CCCC4)\allowbreak{}c(\allowbreak{}=\allowbreak{}O)\allowbreak{}c3ccccc32)\allowbreak{}cc1} & \texttt{COc1ccc(\allowbreak{}Cn2cc(\allowbreak{}C(\allowbreak{}=\allowbreak{}O)\allowbreak{}NCCCNc3c4c(\allowbreak{}[nH+]\allowbreak{}c5ccc(\allowbreak{}OC)\allowbreak{}cc35)\allowbreak{}CCCC4)\allowbreak{}c(\allowbreak{}=\allowbreak{}O)\allowbreak{}c3ccccc32)\allowbreak{}cc1} & 5.65 & +1 & test & control \\
CHEMBL4217139 & \texttt{CCCCCCCCCCCCNCC(\allowbreak{}=\allowbreak{}O)\allowbreak{}Nc1ccc2nc3n(\allowbreak{}c(\allowbreak{}=\allowbreak{}O)\allowbreak{}c2c1)\allowbreak{}CCC3} & \texttt{CCCCCCCCCCCC[NH2+]\allowbreak{}CC(\allowbreak{}=\allowbreak{}O)\allowbreak{}Nc1ccc2nc3n(\allowbreak{}c(\allowbreak{}=\allowbreak{}O)\allowbreak{}c2c1)\allowbreak{}CCC3} & 7.43 & +1 & fold 0 & control \\
CHEMBL4217346 & \texttt{COc1ccc(\allowbreak{}Cn2cc(\allowbreak{}C(\allowbreak{}=\allowbreak{}O)\allowbreak{}NCCCCCCNc3c4c(\allowbreak{}nc5cc(\allowbreak{}Cl)\allowbreak{}ccc35)\allowbreak{}CCCC4)\allowbreak{}c(\allowbreak{}=\allowbreak{}O)\allowbreak{}c3ccccc32)\allowbreak{}cc1} & \texttt{COc1ccc(\allowbreak{}Cn2cc(\allowbreak{}C(\allowbreak{}=\allowbreak{}O)\allowbreak{}NCCCCCCNc3c4c(\allowbreak{}[nH+]\allowbreak{}c5cc(\allowbreak{}Cl)\allowbreak{}ccc35)\allowbreak{}CCCC4)\allowbreak{}c(\allowbreak{}=\allowbreak{}O)\allowbreak{}c3ccccc32)\allowbreak{}cc1} & 7.03 & +1 & fold 4 & control \\
CHEMBL426441 & \texttt{CN1C(\allowbreak{}=\allowbreak{}O)\allowbreak{}Cc2cc3onc(\allowbreak{}CCC4CCN(\allowbreak{}Cc5ccccc5)\allowbreak{}CC4)\allowbreak{}c3cc21} & \texttt{CN1C(\allowbreak{}=\allowbreak{}O)\allowbreak{}Cc2cc3onc(\allowbreak{}CCC4CC[NH+]\allowbreak{}(\allowbreak{}Cc5ccccc5)\allowbreak{}CC4)\allowbreak{}c3cc21} & 9.32 & +1 & fold 1 & control \\
CHEMBL4278287 & \texttt{COc1ccc2nc3c(\allowbreak{}c(\allowbreak{}NCCCNC(\allowbreak{}=\allowbreak{}O)\allowbreak{}c4sc(\allowbreak{}-c5ccnc(\allowbreak{}NC(\allowbreak{}=\allowbreak{}O)\allowbreak{}C6CC6)\allowbreak{}c5)\allowbreak{}nc4OC)\allowbreak{}c2c1)\allowbreak{}CCCC3} & \texttt{COc1ccc2[nH+]\allowbreak{}c3c(\allowbreak{}c(\allowbreak{}NCCCNC(\allowbreak{}=\allowbreak{}O)\allowbreak{}c4sc(\allowbreak{}-c5ccnc(\allowbreak{}NC(\allowbreak{}=\allowbreak{}O)\allowbreak{}C6CC6)\allowbreak{}c5)\allowbreak{}nc4OC)\allowbreak{}c2c1)\allowbreak{}CCCC3} & 7.42 & +1 & test & control \\
CHEMBL4284047 & \texttt{Cn1c2nc3ccccc3c-2c(\allowbreak{}NCCCCCCNc2c3c(\allowbreak{}nc4ccccc24)\allowbreak{}CCCC3)\allowbreak{}c2ccccc21} & \texttt{Cn1c2nc3ccccc3c-2c(\allowbreak{}NCCCCCCNc2c3c(\allowbreak{}[nH+]\allowbreak{}c4ccccc24)\allowbreak{}CCCC3)\allowbreak{}c2ccccc21} & 8.90 & +1 & fold 0 & control \\
CHEMBL4285581 & \texttt{COc1ccc2nc3c(\allowbreak{}c(\allowbreak{}NCCNC(\allowbreak{}=\allowbreak{}O)\allowbreak{}c4sc(\allowbreak{}-c5ccnc(\allowbreak{}NC(\allowbreak{}=\allowbreak{}O)\allowbreak{}C6CC6)\allowbreak{}c5)\allowbreak{}nc4OC)\allowbreak{}c2c1)\allowbreak{}CCCC3} & \texttt{COc1ccc2[nH+]\allowbreak{}c3c(\allowbreak{}c(\allowbreak{}NCCNC(\allowbreak{}=\allowbreak{}O)\allowbreak{}c4sc(\allowbreak{}-c5ccnc(\allowbreak{}NC(\allowbreak{}=\allowbreak{}O)\allowbreak{}C6CC6)\allowbreak{}c5)\allowbreak{}nc4OC)\allowbreak{}c2c1)\allowbreak{}CCCC3} & 7.64 & +1 & test & control \\
CHEMBL4294640 & \texttt{Cn1c2nc3ccccc3c-2c(\allowbreak{}NCCCCCCCCNc2c3c(\allowbreak{}nc4ccccc24)\allowbreak{}CCCC3)\allowbreak{}c2ccccc21} & \texttt{Cn1c2nc3ccccc3c-2c(\allowbreak{}NCCCCCCCCNc2c3c(\allowbreak{}[nH+]\allowbreak{}c4ccccc24)\allowbreak{}CCCC3)\allowbreak{}c2ccccc21} & 9.02 & +1 & fold 0 & control \\
CHEMBL431519 & \texttt{CN1CC[C@@]\allowbreak{}2(\allowbreak{}C)\allowbreak{}c3cc(\allowbreak{}OC(\allowbreak{}=\allowbreak{}O)\allowbreak{}NCCCCCCCCCCN4CCOCC4)\allowbreak{}ccc3N(\allowbreak{}C)\allowbreak{}[C@@H]\allowbreak{}12} & \texttt{CN1c2ccc(\allowbreak{}OC(\allowbreak{}=\allowbreak{}O)\allowbreak{}NCCCCCCCCCC[NH+]\allowbreak{}3CCOCC3)\allowbreak{}cc2[C@]\allowbreak{}2(\allowbreak{}C)\allowbreak{}CC[NH+]\allowbreak{}(\allowbreak{}C)\allowbreak{}[C@H]\allowbreak{}12} & 7.40 & +2 & fold 2 & cliff \\
CHEMBL434378 & \texttt{COc1cc(\allowbreak{}C(\allowbreak{}=\allowbreak{}O)\allowbreak{}NCCCCCC(\allowbreak{}=\allowbreak{}O)\allowbreak{}NNc2c3c(\allowbreak{}nc4ccccc24)\allowbreak{}CCCC3)\allowbreak{}cc(\allowbreak{}OC)\allowbreak{}c1OC} & \texttt{COc1cc(\allowbreak{}C(\allowbreak{}=\allowbreak{}O)\allowbreak{}NCCCCCC(\allowbreak{}=\allowbreak{}O)\allowbreak{}NNc2c3c(\allowbreak{}[nH+]\allowbreak{}c4ccccc24)\allowbreak{}CCCC3)\allowbreak{}cc(\allowbreak{}OC)\allowbreak{}c1OC} & 7.74 & +1 & test & control \\
CHEMBL4438626 & \texttt{O=\allowbreak{}C(\allowbreak{}CCCNc1c2c(\allowbreak{}nc3cc(\allowbreak{}Cl)\allowbreak{}ccc13)\allowbreak{}CCCC2)\allowbreak{}NC12CC3CC(\allowbreak{}F)\allowbreak{}(\allowbreak{}CC(\allowbreak{}C1)\allowbreak{}c1ccccc13)\allowbreak{}C2} & \texttt{O=\allowbreak{}C(\allowbreak{}CCCNc1c2c(\allowbreak{}[nH+]\allowbreak{}c3cc(\allowbreak{}Cl)\allowbreak{}ccc13)\allowbreak{}CCCC2)\allowbreak{}NC12CC3CC(\allowbreak{}F)\allowbreak{}(\allowbreak{}CC(\allowbreak{}C1)\allowbreak{}c1ccccc13)\allowbreak{}C2} & 8.85 & +1 & test & control \\
CHEMBL4455690 & \texttt{Cc1ccc2oc(\allowbreak{}-c3ccc(\allowbreak{}OCCCCCCN4CCCCCC4)\allowbreak{}cc3)\allowbreak{}cc(\allowbreak{}=\allowbreak{}O)\allowbreak{}c2c1} & \texttt{Cc1ccc2oc(\allowbreak{}-c3ccc(\allowbreak{}OCCCCCC[NH+]\allowbreak{}4CCCCCC4)\allowbreak{}cc3)\allowbreak{}cc(\allowbreak{}=\allowbreak{}O)\allowbreak{}c2c1} & 6.32 & +1 & fold 2 & control \\
CHEMBL4462369 & \texttt{O=\allowbreak{}C(\allowbreak{}CCCCNc1c2c(\allowbreak{}nc3cc(\allowbreak{}Cl)\allowbreak{}ccc13)\allowbreak{}CCCC2)\allowbreak{}NC12CC3CC(\allowbreak{}F)\allowbreak{}(\allowbreak{}CC(\allowbreak{}C1)\allowbreak{}c1ccccc13)\allowbreak{}C2} & \texttt{O=\allowbreak{}C(\allowbreak{}CCCCNc1c2c(\allowbreak{}[nH+]\allowbreak{}c3cc(\allowbreak{}Cl)\allowbreak{}ccc13)\allowbreak{}CCCC2)\allowbreak{}NC12CC3CC(\allowbreak{}F)\allowbreak{}(\allowbreak{}CC(\allowbreak{}C1)\allowbreak{}c1ccccc13)\allowbreak{}C2} & 8.89 & +1 & test & control \\
CHEMBL4467355 & \texttt{COc1ccc2nc3c(\allowbreak{}c(\allowbreak{}NCCCCCCNC(\allowbreak{}=\allowbreak{}O)\allowbreak{}[C@@H]\allowbreak{}(\allowbreak{}N)\allowbreak{}Cc4c[nH]\allowbreak{}c5ccccc45)\allowbreak{}c2c1)\allowbreak{}CCCC3} & \texttt{COc1ccc2[nH+]\allowbreak{}c3c(\allowbreak{}c(\allowbreak{}NCCCCCCNC(\allowbreak{}=\allowbreak{}O)\allowbreak{}[C@@H]\allowbreak{}(\allowbreak{}[NH3+]\allowbreak{})\allowbreak{}Cc4c[nH]\allowbreak{}c5ccccc45)\allowbreak{}c2c1)\allowbreak{}CCCC3} & 6.03 & +2 & fold 2 & control \\
CHEMBL4467520 & \texttt{Cc1ccc2[nH]\allowbreak{}c3c(\allowbreak{}c2c1)\allowbreak{}CN(\allowbreak{}CCCCCN1CCc2[nH]\allowbreak{}c4ccc(\allowbreak{}C)\allowbreak{}cc4c2C1)\allowbreak{}CC3} & \texttt{Cc1ccc2[nH]\allowbreak{}c3c(\allowbreak{}c2c1)\allowbreak{}C[NH+]\allowbreak{}(\allowbreak{}CCCCC[NH+]\allowbreak{}1CCc2[nH]\allowbreak{}c4ccc(\allowbreak{}C)\allowbreak{}cc4c2C1)\allowbreak{}CC3} & 6.71 & +2 & fold 4 & control \\
CHEMBL4474399 & \texttt{Cc1ccc2oc(\allowbreak{}-c3ccc(\allowbreak{}OCCCCCCN4CCCCC4)\allowbreak{}cc3)\allowbreak{}cc(\allowbreak{}=\allowbreak{}O)\allowbreak{}c2c1} & \texttt{Cc1ccc2oc(\allowbreak{}-c3ccc(\allowbreak{}OCCCCCC[NH+]\allowbreak{}4CCCCC4)\allowbreak{}cc3)\allowbreak{}cc(\allowbreak{}=\allowbreak{}O)\allowbreak{}c2c1} & 6.44 & +1 & fold 2 & control \\
CHEMBL4536838 & \texttt{COc1ccc2[nH]\allowbreak{}cc(\allowbreak{}CCNC(\allowbreak{}=\allowbreak{}O)\allowbreak{}CN(\allowbreak{}CCC3CCN(\allowbreak{}Cc4ccccc4)\allowbreak{}CC3)\allowbreak{}C(\allowbreak{}=\allowbreak{}O)\allowbreak{}/\allowbreak{}C=\allowbreak{}C/\allowbreak{}c3ccc(\allowbreak{}O)\allowbreak{}c(\allowbreak{}OC)\allowbreak{}c3)\allowbreak{}c2c1} & \texttt{COc1ccc2[nH]\allowbreak{}cc(\allowbreak{}CCNC(\allowbreak{}=\allowbreak{}O)\allowbreak{}CN(\allowbreak{}CCC3CC[NH+]\allowbreak{}(\allowbreak{}Cc4ccccc4)\allowbreak{}CC3)\allowbreak{}C(\allowbreak{}=\allowbreak{}O)\allowbreak{}/\allowbreak{}C=\allowbreak{}C/\allowbreak{}c3ccc(\allowbreak{}O)\allowbreak{}c(\allowbreak{}OC)\allowbreak{}c3)\allowbreak{}c2c1} & 7.51 & +1 & fold 3 & control \\
CHEMBL4541077 & \texttt{COc1ccc2[nH]\allowbreak{}cc(\allowbreak{}CCNC(\allowbreak{}=\allowbreak{}O)\allowbreak{}CN(\allowbreak{}CCCC3CCN(\allowbreak{}Cc4ccccc4)\allowbreak{}CC3)\allowbreak{}C(\allowbreak{}=\allowbreak{}O)\allowbreak{}/\allowbreak{}C=\allowbreak{}C/\allowbreak{}c3ccc(\allowbreak{}O)\allowbreak{}c(\allowbreak{}OC)\allowbreak{}c3)\allowbreak{}c2c1} & \texttt{COc1ccc2[nH]\allowbreak{}cc(\allowbreak{}CCNC(\allowbreak{}=\allowbreak{}O)\allowbreak{}CN(\allowbreak{}CCCC3CC[NH+]\allowbreak{}(\allowbreak{}Cc4ccccc4)\allowbreak{}CC3)\allowbreak{}C(\allowbreak{}=\allowbreak{}O)\allowbreak{}/\allowbreak{}C=\allowbreak{}C/\allowbreak{}c3ccc(\allowbreak{}O)\allowbreak{}c(\allowbreak{}OC)\allowbreak{}c3)\allowbreak{}c2c1} & 7.58 & +1 & fold 3 & control \\
CHEMBL4551316 & \texttt{COc1ccc(\allowbreak{}CN2CCC(\allowbreak{}Nc3nc(\allowbreak{}-c4ccc(\allowbreak{}Cl)\allowbreak{}cc4)\allowbreak{}c(\allowbreak{}-c4ccc(\allowbreak{}Cl)\allowbreak{}cc4)\allowbreak{}s3)\allowbreak{}CC2)\allowbreak{}cc1} & \texttt{COc1ccc(\allowbreak{}C[NH+]\allowbreak{}2CCC(\allowbreak{}Nc3[nH+]\allowbreak{}c(\allowbreak{}-c4ccc(\allowbreak{}Cl)\allowbreak{}cc4)\allowbreak{}c(\allowbreak{}-c4ccc(\allowbreak{}Cl)\allowbreak{}cc4)\allowbreak{}s3)\allowbreak{}CC2)\allowbreak{}cc1} & 5.95 & +2 & test & control \\
CHEMBL4556372 & \texttt{COc1ccc2nc3c(\allowbreak{}c(\allowbreak{}NCCCNC(\allowbreak{}=\allowbreak{}O)\allowbreak{}[C@@H]\allowbreak{}(\allowbreak{}N)\allowbreak{}Cc4c[nH]\allowbreak{}c5ccccc45)\allowbreak{}c2c1)\allowbreak{}CCCC3} & \texttt{COc1ccc2[nH+]\allowbreak{}c3c(\allowbreak{}c(\allowbreak{}NCCCNC(\allowbreak{}=\allowbreak{}O)\allowbreak{}[C@@H]\allowbreak{}(\allowbreak{}[NH3+]\allowbreak{})\allowbreak{}Cc4c[nH]\allowbreak{}c5ccccc45)\allowbreak{}c2c1)\allowbreak{}CCCC3} & 5.89 & +2 & fold 2 & control \\
CHEMBL4559477 & \texttt{N[C@@H]\allowbreak{}(\allowbreak{}Cc1c[nH]\allowbreak{}c2ccccc12)\allowbreak{}C(\allowbreak{}=\allowbreak{}O)\allowbreak{}NCCCCCCCNc1c2c(\allowbreak{}nc3ccccc13)\allowbreak{}CCCC2} & \texttt{[NH3+]\allowbreak{}[C@@H]\allowbreak{}(\allowbreak{}Cc1c[nH]\allowbreak{}c2ccccc12)\allowbreak{}C(\allowbreak{}=\allowbreak{}O)\allowbreak{}NCCCCCCCNc1c2c(\allowbreak{}[nH+]\allowbreak{}c3ccccc13)\allowbreak{}CCCC2} & 6.92 & +2 & fold 0 & control \\
CHEMBL4562938 & \texttt{COc1ccc(\allowbreak{}CN2CCC(\allowbreak{}Nc3nc(\allowbreak{}-c4ccc(\allowbreak{}OC)\allowbreak{}cc4)\allowbreak{}c(\allowbreak{}-c4ccc(\allowbreak{}OC)\allowbreak{}cc4)\allowbreak{}s3)\allowbreak{}CC2)\allowbreak{}cc1} & \texttt{COc1ccc(\allowbreak{}C[NH+]\allowbreak{}2CCC(\allowbreak{}Nc3[nH+]\allowbreak{}c(\allowbreak{}-c4ccc(\allowbreak{}OC)\allowbreak{}cc4)\allowbreak{}c(\allowbreak{}-c4ccc(\allowbreak{}OC)\allowbreak{}cc4)\allowbreak{}s3)\allowbreak{}CC2)\allowbreak{}cc1} & 6.11 & +2 & test & control \\
CHEMBL473866 & \texttt{Clc1ccc2c(\allowbreak{}NCCCNCCCCCOc3cccc4[nH]\allowbreak{}c5ccccc5c34)\allowbreak{}c3c(\allowbreak{}nc2c1)\allowbreak{}CCCC3} & \texttt{Clc1ccc2c(\allowbreak{}NCCC[NH2+]\allowbreak{}CCCCCOc3cccc4[nH]\allowbreak{}c5ccccc5c34)\allowbreak{}c3c(\allowbreak{}[nH+]\allowbreak{}c2c1)\allowbreak{}CCCC3} & 8.81 & +2 & fold 0 & control \\
CHEMBL4747041 & \texttt{COc1cc2c(\allowbreak{}cc1OCc1ccccc1)\allowbreak{}CCC(\allowbreak{}C)\allowbreak{}(\allowbreak{}CC(\allowbreak{}=\allowbreak{}O)\allowbreak{}NCCCCCCCNc1c3c(\allowbreak{}nc4cc(\allowbreak{}Cl)\allowbreak{}ccc14)\allowbreak{}CCCC3)\allowbreak{}O2} & \texttt{COc1cc2c(\allowbreak{}cc1OCc1ccccc1)\allowbreak{}CCC(\allowbreak{}C)\allowbreak{}(\allowbreak{}CC(\allowbreak{}=\allowbreak{}O)\allowbreak{}NCCCCCCCNc1c3c(\allowbreak{}[nH+]\allowbreak{}c4cc(\allowbreak{}Cl)\allowbreak{}ccc14)\allowbreak{}CCCC3)\allowbreak{}O2} & 8.16 & +1 & fold 3 & control \\
CHEMBL4749678 & \texttt{CNC(\allowbreak{}=\allowbreak{}O)\allowbreak{}c1sc(\allowbreak{}-c2ccnc(\allowbreak{}NC(\allowbreak{}=\allowbreak{}O)\allowbreak{}CNCCNc3c4c(\allowbreak{}nc5ccccc35)\allowbreak{}CCCC4)\allowbreak{}c2)\allowbreak{}nc1OCC1CCC1} & \texttt{CNC(\allowbreak{}=\allowbreak{}O)\allowbreak{}c1sc(\allowbreak{}-c2ccnc(\allowbreak{}NC(\allowbreak{}=\allowbreak{}O)\allowbreak{}C[NH2+]\allowbreak{}CCNc3c4c(\allowbreak{}[nH+]\allowbreak{}c5ccccc35)\allowbreak{}CCCC4)\allowbreak{}c2)\allowbreak{}nc1OCC1CCC1} & 6.81 & +2 & fold 2 & cliff \\
CHEMBL4757922 & \texttt{COc1cc2c(\allowbreak{}cc1O)\allowbreak{}CCC(\allowbreak{}C)\allowbreak{}(\allowbreak{}CC(\allowbreak{}=\allowbreak{}O)\allowbreak{}NCCCCCCCCCNc1c3c(\allowbreak{}nc4ccccc14)\allowbreak{}CCCC3)\allowbreak{}O2} & \texttt{COc1cc2c(\allowbreak{}cc1O)\allowbreak{}CCC(\allowbreak{}C)\allowbreak{}(\allowbreak{}CC(\allowbreak{}=\allowbreak{}O)\allowbreak{}NCCCCCCCCCNc1c3c(\allowbreak{}[nH+]\allowbreak{}c4ccccc14)\allowbreak{}CCCC3)\allowbreak{}O2} & 7.24 & +1 & fold 3 & control \\
CHEMBL4776999 & \texttt{CCCC[N+]\allowbreak{}1(\allowbreak{}C)\allowbreak{}CCN(\allowbreak{}c2ccc(\allowbreak{}/\allowbreak{}C=\allowbreak{}C3\textbackslash{}Cc4cc(\allowbreak{}OC)\allowbreak{}c(\allowbreak{}OC)\allowbreak{}cc4C3=\allowbreak{}O)\allowbreak{}cc2)\allowbreak{}CC1} & \texttt{CCCC[N+]\allowbreak{}1(\allowbreak{}C)\allowbreak{}CCN(\allowbreak{}c2ccc(\allowbreak{}/\allowbreak{}C=\allowbreak{}C3\textbackslash{}Cc4cc(\allowbreak{}OC)\allowbreak{}c(\allowbreak{}OC)\allowbreak{}cc4C3=\allowbreak{}O)\allowbreak{}cc2)\allowbreak{}CC1} & 6.27 & +1 & fold 2 & control \\
CHEMBL478666 & \texttt{Clc1ccc2c(\allowbreak{}NCCCNCCCOc3cccc4[nH]\allowbreak{}c5ccccc5c34)\allowbreak{}c3c(\allowbreak{}nc2c1)\allowbreak{}CCCC3} & \texttt{Clc1ccc2c(\allowbreak{}NCCC[NH2+]\allowbreak{}CCCOc3cccc4[nH]\allowbreak{}c5ccccc5c34)\allowbreak{}c3c(\allowbreak{}[nH+]\allowbreak{}c2c1)\allowbreak{}CCCC3} & 8.67 & +2 & fold 0 & control \\
CHEMBL4789145 & \texttt{CNC(\allowbreak{}=\allowbreak{}O)\allowbreak{}c1sc(\allowbreak{}-c2ccnc(\allowbreak{}NC(\allowbreak{}=\allowbreak{}O)\allowbreak{}CNCCNc3c4c(\allowbreak{}nc5ccccc35)\allowbreak{}CCCC4)\allowbreak{}c2)\allowbreak{}nc1OCC1CC1} & \texttt{CNC(\allowbreak{}=\allowbreak{}O)\allowbreak{}c1sc(\allowbreak{}-c2ccnc(\allowbreak{}NC(\allowbreak{}=\allowbreak{}O)\allowbreak{}C[NH2+]\allowbreak{}CCNc3c4c(\allowbreak{}[nH+]\allowbreak{}c5ccccc35)\allowbreak{}CCCC4)\allowbreak{}c2)\allowbreak{}nc1OCC1CC1} & 8.92 & +2 & fold 2 & cliff \\
CHEMBL4789172 & \texttt{COc1cc2c(\allowbreak{}cc1OCc1ccccc1)\allowbreak{}CCC(\allowbreak{}C)\allowbreak{}(\allowbreak{}CC(\allowbreak{}=\allowbreak{}O)\allowbreak{}NCCCCCCNc1c3c(\allowbreak{}nc4cc(\allowbreak{}Cl)\allowbreak{}ccc14)\allowbreak{}CCCC3)\allowbreak{}O2} & \texttt{COc1cc2c(\allowbreak{}cc1OCc1ccccc1)\allowbreak{}CCC(\allowbreak{}C)\allowbreak{}(\allowbreak{}CC(\allowbreak{}=\allowbreak{}O)\allowbreak{}NCCCCCCNc1c3c(\allowbreak{}[nH+]\allowbreak{}c4cc(\allowbreak{}Cl)\allowbreak{}ccc14)\allowbreak{}CCCC3)\allowbreak{}O2} & 7.96 & +1 & fold 3 & control \\
CHEMBL4792460 & \texttt{COc1cc2c(\allowbreak{}cc1O)\allowbreak{}CCC(\allowbreak{}C)\allowbreak{}(\allowbreak{}CC(\allowbreak{}=\allowbreak{}O)\allowbreak{}NCCCCCCCCNc1c3c(\allowbreak{}nc4ccccc14)\allowbreak{}CCCC3)\allowbreak{}O2} & \texttt{COc1cc2c(\allowbreak{}cc1O)\allowbreak{}CCC(\allowbreak{}C)\allowbreak{}(\allowbreak{}CC(\allowbreak{}=\allowbreak{}O)\allowbreak{}NCCCCCCCCNc1c3c(\allowbreak{}[nH+]\allowbreak{}c4ccccc14)\allowbreak{}CCCC3)\allowbreak{}O2} & 7.16 & +1 & fold 3 & control \\
CHEMBL4797677 & \texttt{CCCCC[N+]\allowbreak{}1(\allowbreak{}C)\allowbreak{}CCN(\allowbreak{}c2ccc(\allowbreak{}/\allowbreak{}C=\allowbreak{}C3\textbackslash{}Cc4cc(\allowbreak{}OC)\allowbreak{}c(\allowbreak{}OC)\allowbreak{}cc4C3=\allowbreak{}O)\allowbreak{}cc2)\allowbreak{}CC1} & \texttt{CCCCC[N+]\allowbreak{}1(\allowbreak{}C)\allowbreak{}CCN(\allowbreak{}c2ccc(\allowbreak{}/\allowbreak{}C=\allowbreak{}C3\textbackslash{}Cc4cc(\allowbreak{}OC)\allowbreak{}c(\allowbreak{}OC)\allowbreak{}cc4C3=\allowbreak{}O)\allowbreak{}cc2)\allowbreak{}CC1} & 6.27 & +1 & fold 2 & control \\
CHEMBL4798419 & \texttt{CNC(\allowbreak{}=\allowbreak{}O)\allowbreak{}c1sc(\allowbreak{}-c2ccnc(\allowbreak{}NC(\allowbreak{}=\allowbreak{}O)\allowbreak{}CNCCCCCCCCNc3c4c(\allowbreak{}nc5ccccc35)\allowbreak{}CCCC4)\allowbreak{}c2)\allowbreak{}nc1OCC1CC1} & \texttt{CNC(\allowbreak{}=\allowbreak{}O)\allowbreak{}c1sc(\allowbreak{}-c2ccnc(\allowbreak{}NC(\allowbreak{}=\allowbreak{}O)\allowbreak{}C[NH2+]\allowbreak{}CCCCCCCCNc3c4c(\allowbreak{}[nH+]\allowbreak{}c5ccccc35)\allowbreak{}CCCC4)\allowbreak{}c2)\allowbreak{}nc1OCC1CC1} & 7.82 & +2 & fold 4 & control \\
CHEMBL4800403 & \texttt{CNC(\allowbreak{}=\allowbreak{}O)\allowbreak{}c1sc(\allowbreak{}-c2ccnc(\allowbreak{}NC(\allowbreak{}=\allowbreak{}O)\allowbreak{}CNCCCCCCNc3c4c(\allowbreak{}nc5ccccc35)\allowbreak{}CCCC4)\allowbreak{}c2)\allowbreak{}nc1OCC1CC1} & \texttt{CNC(\allowbreak{}=\allowbreak{}O)\allowbreak{}c1sc(\allowbreak{}-c2ccnc(\allowbreak{}NC(\allowbreak{}=\allowbreak{}O)\allowbreak{}C[NH2+]\allowbreak{}CCCCCCNc3c4c(\allowbreak{}[nH+]\allowbreak{}c5ccccc35)\allowbreak{}CCCC4)\allowbreak{}c2)\allowbreak{}nc1OCC1CC1} & 7.66 & +2 & fold 4 & control \\
CHEMBL4856125 & \texttt{Cc1cccc2nc3c(\allowbreak{}c(\allowbreak{}N)\allowbreak{}c12)\allowbreak{}CCCCC3} & \texttt{Cc1cccc2[nH+]\allowbreak{}c3c(\allowbreak{}c(\allowbreak{}N)\allowbreak{}c12)\allowbreak{}CCCCC3} & 6.59 & +1 & test & control \\
CHEMBL4859128 & \texttt{Cc1cccc2nc3c(\allowbreak{}c(\allowbreak{}N)\allowbreak{}c12)\allowbreak{}CCC3} & \texttt{Cc1cccc2[nH+]\allowbreak{}c3c(\allowbreak{}c(\allowbreak{}N)\allowbreak{}c12)\allowbreak{}CCC3} & 6.38 & +1 & test & control \\
CHEMBL5178885 & \texttt{O=\allowbreak{}C(\allowbreak{}Nc1ccc(\allowbreak{}OC(\allowbreak{}F)\allowbreak{}(\allowbreak{}F)\allowbreak{}F)\allowbreak{}cc1)\allowbreak{}NC1CCN(\allowbreak{}C(\allowbreak{}=\allowbreak{}O)\allowbreak{}CCCCNc2c3c(\allowbreak{}nc4cc(\allowbreak{}Cl)\allowbreak{}ccc24)\allowbreak{}CCCC3)\allowbreak{}CC1} & \texttt{O=\allowbreak{}C(\allowbreak{}Nc1ccc(\allowbreak{}OC(\allowbreak{}F)\allowbreak{}(\allowbreak{}F)\allowbreak{}F)\allowbreak{}cc1)\allowbreak{}NC1CCN(\allowbreak{}C(\allowbreak{}=\allowbreak{}O)\allowbreak{}CCCCNc2c3c(\allowbreak{}[nH+]\allowbreak{}c4cc(\allowbreak{}Cl)\allowbreak{}ccc24)\allowbreak{}CCCC3)\allowbreak{}CC1} & 7.89 & +1 & fold 4 & control \\
CHEMBL5204900 & \texttt{O=\allowbreak{}C(\allowbreak{}Nc1ccc(\allowbreak{}OC(\allowbreak{}F)\allowbreak{}(\allowbreak{}F)\allowbreak{}F)\allowbreak{}cc1)\allowbreak{}NC1CCN(\allowbreak{}C(\allowbreak{}=\allowbreak{}O)\allowbreak{}CCNc2c3c(\allowbreak{}nc4cc(\allowbreak{}Cl)\allowbreak{}ccc24)\allowbreak{}CCCC3)\allowbreak{}CC1} & \texttt{O=\allowbreak{}C(\allowbreak{}Nc1ccc(\allowbreak{}OC(\allowbreak{}F)\allowbreak{}(\allowbreak{}F)\allowbreak{}F)\allowbreak{}cc1)\allowbreak{}NC1CCN(\allowbreak{}C(\allowbreak{}=\allowbreak{}O)\allowbreak{}CCNc2c3c(\allowbreak{}[nH+]\allowbreak{}c4cc(\allowbreak{}Cl)\allowbreak{}ccc24)\allowbreak{}CCCC3)\allowbreak{}CC1} & 7.84 & +1 & fold 4 & control \\
CHEMBL5220283 & \texttt{CNC(\allowbreak{}=\allowbreak{}O)\allowbreak{}Oc1cccc(\allowbreak{}CN(\allowbreak{}C)\allowbreak{}CCCCCCOc2ccc3c(\allowbreak{}c2)\allowbreak{}O/\allowbreak{}C(\allowbreak{}=\allowbreak{}C\textbackslash{}c2cc(\allowbreak{}OC)\allowbreak{}c(\allowbreak{}OC)\allowbreak{}c(\allowbreak{}OC)\allowbreak{}c2)\allowbreak{}C3=\allowbreak{}O)\allowbreak{}c1} & \texttt{CNC(\allowbreak{}=\allowbreak{}O)\allowbreak{}Oc1cccc(\allowbreak{}C[NH+]\allowbreak{}(\allowbreak{}C)\allowbreak{}CCCCCCOc2ccc3c(\allowbreak{}c2)\allowbreak{}O/\allowbreak{}C(\allowbreak{}=\allowbreak{}C\textbackslash{}c2cc(\allowbreak{}OC)\allowbreak{}c(\allowbreak{}OC)\allowbreak{}c(\allowbreak{}OC)\allowbreak{}c2)\allowbreak{}C3=\allowbreak{}O)\allowbreak{}c1} & 9.28 & +1 & fold 4 & cliff \\
CHEMBL5272675 & \texttt{COc1ccc2c(\allowbreak{}c1)\allowbreak{}c(\allowbreak{}C=\allowbreak{}O)\allowbreak{}cn2CCCCCCCN(\allowbreak{}C)\allowbreak{}Cc1cccc(\allowbreak{}O)\allowbreak{}c1} & \texttt{COc1ccc2c(\allowbreak{}c1)\allowbreak{}c(\allowbreak{}C=\allowbreak{}O)\allowbreak{}cn2CCCCCCC[NH+]\allowbreak{}(\allowbreak{}C)\allowbreak{}Cc1cccc(\allowbreak{}O)\allowbreak{}c1} & 5.26 & +1 & fold 3 & control \\
CHEMBL5278778 & \texttt{COc1ccc2c(\allowbreak{}c1)\allowbreak{}c(\allowbreak{}C=\allowbreak{}O)\allowbreak{}cn2CCCCCCCCCCN(\allowbreak{}C)\allowbreak{}Cc1ccc(\allowbreak{}O)\allowbreak{}cc1} & \texttt{COc1ccc2c(\allowbreak{}c1)\allowbreak{}c(\allowbreak{}C=\allowbreak{}O)\allowbreak{}cn2CCCCCCCCCC[NH+]\allowbreak{}(\allowbreak{}C)\allowbreak{}Cc1ccc(\allowbreak{}O)\allowbreak{}cc1} & 6.22 & +1 & test & cliff \\
CHEMBL5283238 & \texttt{COc1ccc2c(\allowbreak{}c1)\allowbreak{}c(\allowbreak{}C=\allowbreak{}O)\allowbreak{}cn2CCCCCCCCCCCCN(\allowbreak{}C)\allowbreak{}Cc1ccc(\allowbreak{}O)\allowbreak{}cc1} & \texttt{COc1ccc2c(\allowbreak{}c1)\allowbreak{}c(\allowbreak{}C=\allowbreak{}O)\allowbreak{}cn2CCCCCCCCCCCC[NH+]\allowbreak{}(\allowbreak{}C)\allowbreak{}Cc1ccc(\allowbreak{}O)\allowbreak{}cc1} & 4.22 & +1 & test & cliff \\
CHEMBL5283851 & \texttt{COc1ccc2c(\allowbreak{}c1)\allowbreak{}c(\allowbreak{}C=\allowbreak{}O)\allowbreak{}cn2CCCCCCCCCCN(\allowbreak{}C)\allowbreak{}Cc1cccc(\allowbreak{}O)\allowbreak{}c1} & \texttt{COc1ccc2c(\allowbreak{}c1)\allowbreak{}c(\allowbreak{}C=\allowbreak{}O)\allowbreak{}cn2CCCCCCCCCC[NH+]\allowbreak{}(\allowbreak{}C)\allowbreak{}Cc1cccc(\allowbreak{}O)\allowbreak{}c1} & 5.34 & +1 & fold 3 & control \\
CHEMBL5285645 & \texttt{O=\allowbreak{}C1CC2Cc3ccsc3C2=\allowbreak{}NN1CCCCCCN1CCN(\allowbreak{}Cc2cccc3c2ccn3S(\allowbreak{}=\allowbreak{}O)\allowbreak{}(\allowbreak{}=\allowbreak{}O)\allowbreak{}c2ccccc2)\allowbreak{}CC1} & \texttt{O=\allowbreak{}C1CC2Cc3ccsc3C2=\allowbreak{}NN1CCCCCC[NH+]\allowbreak{}1CC[NH+]\allowbreak{}(\allowbreak{}Cc2cccc3c2ccn3S(\allowbreak{}=\allowbreak{}O)\allowbreak{}(\allowbreak{}=\allowbreak{}O)\allowbreak{}c2ccccc2)\allowbreak{}CC1} & 5.04 & +2 & test & control \\
CHEMBL5287809 & \texttt{O=\allowbreak{}C1CC2Cc3ccsc3C2=\allowbreak{}NN1CCCCCN1CCN(\allowbreak{}Cc2cccc3c2ccn3S(\allowbreak{}=\allowbreak{}O)\allowbreak{}(\allowbreak{}=\allowbreak{}O)\allowbreak{}c2ccccc2)\allowbreak{}CC1} & \texttt{O=\allowbreak{}C1CC2Cc3ccsc3C2=\allowbreak{}NN1CCCCC[NH+]\allowbreak{}1CC[NH+]\allowbreak{}(\allowbreak{}Cc2cccc3c2ccn3S(\allowbreak{}=\allowbreak{}O)\allowbreak{}(\allowbreak{}=\allowbreak{}O)\allowbreak{}c2ccccc2)\allowbreak{}CC1} & 5.16 & +2 & test & control \\
CHEMBL542467 & \texttt{CNC(\allowbreak{}=\allowbreak{}O)\allowbreak{}Oc1cccc(\allowbreak{}CN(\allowbreak{}C)\allowbreak{}CCCCCCCCCOc2ccc3c(\allowbreak{}=\allowbreak{}O)\allowbreak{}c4ccccc4oc3c2)\allowbreak{}c1} & \texttt{CNC(\allowbreak{}=\allowbreak{}O)\allowbreak{}Oc1cccc(\allowbreak{}C[NH+]\allowbreak{}(\allowbreak{}C)\allowbreak{}CCCCCCCCCOc2ccc3c(\allowbreak{}=\allowbreak{}O)\allowbreak{}c4ccccc4oc3c2)\allowbreak{}c1} & 7.93 & +1 & fold 2 & control \\
CHEMBL544823 & \texttt{CNC(\allowbreak{}=\allowbreak{}O)\allowbreak{}Oc1cccc(\allowbreak{}CN(\allowbreak{}C)\allowbreak{}CCCCCCCCCCOc2ccc3c(\allowbreak{}=\allowbreak{}O)\allowbreak{}c4ccccc4oc3c2)\allowbreak{}c1} & \texttt{CNC(\allowbreak{}=\allowbreak{}O)\allowbreak{}Oc1cccc(\allowbreak{}C[NH+]\allowbreak{}(\allowbreak{}C)\allowbreak{}CCCCCCCCCCOc2ccc3c(\allowbreak{}=\allowbreak{}O)\allowbreak{}c4ccccc4oc3c2)\allowbreak{}c1} & 7.66 & +1 & fold 2 & control \\
CHEMBL5564643 & \texttt{O=\allowbreak{}C(\allowbreak{}NCCCNc1c2c(\allowbreak{}nc3ccccc13)\allowbreak{}CCCC2)\allowbreak{}Oc1ccc2cc[nH]\allowbreak{}c2c1} & \texttt{O=\allowbreak{}C(\allowbreak{}NCCCNc1c2c(\allowbreak{}[nH+]\allowbreak{}c3ccccc13)\allowbreak{}CCCC2)\allowbreak{}Oc1ccc2cc[nH]\allowbreak{}c2c1} & 6.26 & +1 & fold 2 & control \\
CHEMBL5565153 & \texttt{CC(\allowbreak{}C)\allowbreak{}(\allowbreak{}C)\allowbreak{}c1ccc(\allowbreak{}Oc2ccc3nc4c(\allowbreak{}c(\allowbreak{}N)\allowbreak{}c3c2)\allowbreak{}CCC4)\allowbreak{}cc1} & \texttt{CC(\allowbreak{}C)\allowbreak{}(\allowbreak{}C)\allowbreak{}c1ccc(\allowbreak{}Oc2ccc3[nH+]\allowbreak{}c4c(\allowbreak{}c(\allowbreak{}N)\allowbreak{}c3c2)\allowbreak{}CCC4)\allowbreak{}cc1} & 5.56 & +1 & fold 0 & control \\
CHEMBL5565570 & \texttt{Nc1c2c(\allowbreak{}nc3ccc(\allowbreak{}Oc4ccc(\allowbreak{}Cl)\allowbreak{}cc4)\allowbreak{}cc13)\allowbreak{}CCC2} & \texttt{Nc1c2c(\allowbreak{}[nH+]\allowbreak{}c3ccc(\allowbreak{}Oc4ccc(\allowbreak{}Cl)\allowbreak{}cc4)\allowbreak{}cc13)\allowbreak{}CCC2} & 5.29 & +1 & fold 1 & control \\
CHEMBL5571191 & \texttt{Nc1c2c(\allowbreak{}nc3ccc(\allowbreak{}Oc4ccccc4)\allowbreak{}cc13)\allowbreak{}CCCCC2} & \texttt{Nc1c2c(\allowbreak{}[nH+]\allowbreak{}c3ccc(\allowbreak{}Oc4ccccc4)\allowbreak{}cc13)\allowbreak{}CCCCC2} & 5.27 & +1 & test & control \\
CHEMBL5572528 & \texttt{CC(\allowbreak{}C)\allowbreak{}(\allowbreak{}C)\allowbreak{}c1ccc(\allowbreak{}Oc2ccc3nc4c(\allowbreak{}c(\allowbreak{}N)\allowbreak{}c3c2)\allowbreak{}CCCC4)\allowbreak{}cc1} & \texttt{CC(\allowbreak{}C)\allowbreak{}(\allowbreak{}C)\allowbreak{}c1ccc(\allowbreak{}Oc2ccc3[nH+]\allowbreak{}c4c(\allowbreak{}c(\allowbreak{}N)\allowbreak{}c3c2)\allowbreak{}CCCC4)\allowbreak{}cc1} & 5.71 & +1 & fold 0 & control \\
CHEMBL5573162 & \texttt{Nc1c2c(\allowbreak{}nc3ccc(\allowbreak{}Oc4ccc(\allowbreak{}Cl)\allowbreak{}cc4)\allowbreak{}cc13)\allowbreak{}CCCCC2} & \texttt{Nc1c2c(\allowbreak{}[nH+]\allowbreak{}c3ccc(\allowbreak{}Oc4ccc(\allowbreak{}Cl)\allowbreak{}cc4)\allowbreak{}cc13)\allowbreak{}CCCCC2} & 5.38 & +1 & fold 1 & control \\
CHEMBL5573526 & \texttt{COc1cccc(\allowbreak{}Oc2ccc3nc4c(\allowbreak{}c(\allowbreak{}N)\allowbreak{}c3c2)\allowbreak{}CCC4)\allowbreak{}c1} & \texttt{COc1cccc(\allowbreak{}Oc2ccc3[nH+]\allowbreak{}c4c(\allowbreak{}c(\allowbreak{}N)\allowbreak{}c3c2)\allowbreak{}CCC4)\allowbreak{}c1} & 5.28 & +1 & fold 3 & control \\
CHEMBL5574879 & \texttt{O=\allowbreak{}C(\allowbreak{}NCCCCCNc1c2c(\allowbreak{}nc3ccccc13)\allowbreak{}CCCC2)\allowbreak{}Oc1ccc2cc[nH]\allowbreak{}c2c1} & \texttt{O=\allowbreak{}C(\allowbreak{}NCCCCCNc1c2c(\allowbreak{}[nH+]\allowbreak{}c3ccccc13)\allowbreak{}CCCC2)\allowbreak{}Oc1ccc2cc[nH]\allowbreak{}c2c1} & 6.38 & +1 & fold 2 & control \\
CHEMBL5575070 & \texttt{COc1ccc(\allowbreak{}Oc2ccc3nc4c(\allowbreak{}c(\allowbreak{}N)\allowbreak{}c3c2)\allowbreak{}CCC4)\allowbreak{}cc1} & \texttt{COc1ccc(\allowbreak{}Oc2ccc3[nH+]\allowbreak{}c4c(\allowbreak{}c(\allowbreak{}N)\allowbreak{}c3c2)\allowbreak{}CCC4)\allowbreak{}cc1} & 5.37 & +1 & fold 4 & control \\
CHEMBL5575441 & \texttt{O=\allowbreak{}C(\allowbreak{}NCCCCCCNc1c2c(\allowbreak{}nc3cc(\allowbreak{}F)\allowbreak{}ccc13)\allowbreak{}CCCC2)\allowbreak{}Oc1ccc2cc[nH]\allowbreak{}c2c1} & \texttt{O=\allowbreak{}C(\allowbreak{}NCCCCCCNc1c2c(\allowbreak{}[nH+]\allowbreak{}c3cc(\allowbreak{}F)\allowbreak{}ccc13)\allowbreak{}CCCC2)\allowbreak{}Oc1ccc2cc[nH]\allowbreak{}c2c1} & 6.53 & +1 & fold 3 & control \\
CHEMBL5583827 & \texttt{O=\allowbreak{}C(\allowbreak{}NCCCCCNc1c2c(\allowbreak{}nc3cc(\allowbreak{}F)\allowbreak{}ccc13)\allowbreak{}CCCC2)\allowbreak{}Oc1ccc2cc[nH]\allowbreak{}c2c1} & \texttt{O=\allowbreak{}C(\allowbreak{}NCCCCCNc1c2c(\allowbreak{}[nH+]\allowbreak{}c3cc(\allowbreak{}F)\allowbreak{}ccc13)\allowbreak{}CCCC2)\allowbreak{}Oc1ccc2cc[nH]\allowbreak{}c2c1} & 6.56 & +1 & fold 3 & control \\
CHEMBL5591282 & \texttt{Nc1c2c(\allowbreak{}nc3ccc(\allowbreak{}Oc4ccccc4)\allowbreak{}cc13)\allowbreak{}CCC2} & \texttt{Nc1c2c(\allowbreak{}[nH+]\allowbreak{}c3ccc(\allowbreak{}Oc4ccccc4)\allowbreak{}cc13)\allowbreak{}CCC2} & 5.19 & +1 & test & control \\
CHEMBL5592792 & \texttt{COc1ccc(\allowbreak{}Oc2ccc3nc4c(\allowbreak{}c(\allowbreak{}N)\allowbreak{}c3c2)\allowbreak{}CCCC4)\allowbreak{}cc1} & \texttt{COc1ccc(\allowbreak{}Oc2ccc3[nH+]\allowbreak{}c4c(\allowbreak{}c(\allowbreak{}N)\allowbreak{}c3c2)\allowbreak{}CCCC4)\allowbreak{}cc1} & 5.54 & +1 & fold 4 & control \\
CHEMBL5593367 & \texttt{COc1cccc(\allowbreak{}Oc2ccc3nc4c(\allowbreak{}c(\allowbreak{}N)\allowbreak{}c3c2)\allowbreak{}CCCC4)\allowbreak{}c1} & \texttt{COc1cccc(\allowbreak{}Oc2ccc3[nH+]\allowbreak{}c4c(\allowbreak{}c(\allowbreak{}N)\allowbreak{}c3c2)\allowbreak{}CCCC4)\allowbreak{}c1} & 5.58 & +1 & fold 3 & control \\
CHEMBL594629 & \texttt{CC(\allowbreak{}C)\allowbreak{}(\allowbreak{}C)\allowbreak{}c1cc[n+]\allowbreak{}(\allowbreak{}CCCCCCCCCCCC[n+]\allowbreak{}2ccc(\allowbreak{}C(\allowbreak{}C)\allowbreak{}(\allowbreak{}C)\allowbreak{}C)\allowbreak{}cc2)\allowbreak{}cc1} & \texttt{CC(\allowbreak{}C)\allowbreak{}(\allowbreak{}C)\allowbreak{}c1cc[n+]\allowbreak{}(\allowbreak{}CCCCCCCCCCCC[n+]\allowbreak{}2ccc(\allowbreak{}C(\allowbreak{}C)\allowbreak{}(\allowbreak{}C)\allowbreak{}C)\allowbreak{}cc2)\allowbreak{}cc1} & 8.15 & +2 & fold 0 & cliff \\
CHEMBL598824 & \texttt{CCN(\allowbreak{}CC)\allowbreak{}CCCCCCCOc1ccc2c(\allowbreak{}c1)\allowbreak{}C(\allowbreak{}=\allowbreak{}O)\allowbreak{}/\allowbreak{}C(\allowbreak{}=\allowbreak{}C/\allowbreak{}c1ccc(\allowbreak{}CN(\allowbreak{}C)\allowbreak{}Cc3ccccc3)\allowbreak{}cc1)\allowbreak{}C2} & \texttt{CC[NH+]\allowbreak{}(\allowbreak{}CC)\allowbreak{}CCCCCCCOc1ccc2c(\allowbreak{}c1)\allowbreak{}C(\allowbreak{}=\allowbreak{}O)\allowbreak{}/\allowbreak{}C(\allowbreak{}=\allowbreak{}C/\allowbreak{}c1ccc(\allowbreak{}C[NH+]\allowbreak{}(\allowbreak{}C)\allowbreak{}Cc3ccccc3)\allowbreak{}cc1)\allowbreak{}C2} & 6.29 & +2 & test & control \\
CHEMBL599058 & \texttt{CCN(\allowbreak{}CC)\allowbreak{}CCCCCCOc1ccc2c(\allowbreak{}c1)\allowbreak{}C(\allowbreak{}=\allowbreak{}O)\allowbreak{}/\allowbreak{}C(\allowbreak{}=\allowbreak{}C/\allowbreak{}c1ccc(\allowbreak{}CN(\allowbreak{}C)\allowbreak{}Cc3ccccc3)\allowbreak{}cc1)\allowbreak{}C2} & \texttt{CC[NH+]\allowbreak{}(\allowbreak{}CC)\allowbreak{}CCCCCCOc1ccc2c(\allowbreak{}c1)\allowbreak{}C(\allowbreak{}=\allowbreak{}O)\allowbreak{}/\allowbreak{}C(\allowbreak{}=\allowbreak{}C/\allowbreak{}c1ccc(\allowbreak{}C[NH+]\allowbreak{}(\allowbreak{}C)\allowbreak{}Cc3ccccc3)\allowbreak{}cc1)\allowbreak{}C2} & 6.42 & +2 & test & control \\
CHEMBL601410 & \texttt{c1ccc2c(\allowbreak{}c1)\allowbreak{}ccc[n+]\allowbreak{}2CCCCCC[n+]\allowbreak{}1cccc2ccccc21} & \texttt{c1ccc2c(\allowbreak{}c1)\allowbreak{}ccc[n+]\allowbreak{}2CCCCCC[n+]\allowbreak{}1cccc2ccccc21} & 4.60 & +2 & test & cliff \\
CHEMBL62664 & \texttt{CN1CC[C@@]\allowbreak{}2(\allowbreak{}C)\allowbreak{}c3cc(\allowbreak{}OC(\allowbreak{}=\allowbreak{}O)\allowbreak{}NCCCCCCCCCCCN4CCOCC4)\allowbreak{}ccc3N(\allowbreak{}C)\allowbreak{}[C@@H]\allowbreak{}12} & \texttt{CN1c2ccc(\allowbreak{}OC(\allowbreak{}=\allowbreak{}O)\allowbreak{}NCCCCCCCCCCC[NH+]\allowbreak{}3CCOCC3)\allowbreak{}cc2[C@]\allowbreak{}2(\allowbreak{}C)\allowbreak{}CC[NH+]\allowbreak{}(\allowbreak{}C)\allowbreak{}[C@H]\allowbreak{}12} & 7.40 & +2 & fold 2 & cliff \\
CHEMBL63173 & \texttt{CN1CC[C@@]\allowbreak{}2(\allowbreak{}C)\allowbreak{}c3cc(\allowbreak{}OC(\allowbreak{}=\allowbreak{}O)\allowbreak{}NCCCCN4CCOCC4)\allowbreak{}ccc3N(\allowbreak{}C)\allowbreak{}[C@@H]\allowbreak{}12} & \texttt{CN1c2ccc(\allowbreak{}OC(\allowbreak{}=\allowbreak{}O)\allowbreak{}NCCCC[NH+]\allowbreak{}3CCOCC3)\allowbreak{}cc2[C@]\allowbreak{}2(\allowbreak{}C)\allowbreak{}CC[NH+]\allowbreak{}(\allowbreak{}C)\allowbreak{}[C@H]\allowbreak{}12} & 4.68 & +2 & fold 2 & cliff \\
CHEMBL64011 & \texttt{CN1CC[C@@]\allowbreak{}2(\allowbreak{}C)\allowbreak{}c3cc(\allowbreak{}OC(\allowbreak{}=\allowbreak{}O)\allowbreak{}NCCCCCCCCCN4CCOCC4)\allowbreak{}ccc3N(\allowbreak{}C)\allowbreak{}[C@@H]\allowbreak{}12} & \texttt{CN1c2ccc(\allowbreak{}OC(\allowbreak{}=\allowbreak{}O)\allowbreak{}NCCCCCCCCC[NH+]\allowbreak{}3CCOCC3)\allowbreak{}cc2[C@]\allowbreak{}2(\allowbreak{}C)\allowbreak{}CC[NH+]\allowbreak{}(\allowbreak{}C)\allowbreak{}[C@H]\allowbreak{}12} & 7.22 & +2 & fold 2 & cliff \\
CHEMBL65018 & \texttt{CN1CC[C@@]\allowbreak{}2(\allowbreak{}C)\allowbreak{}c3cc(\allowbreak{}OC(\allowbreak{}=\allowbreak{}O)\allowbreak{}NCCCCCCCCN4CCOCC4)\allowbreak{}ccc3N(\allowbreak{}C)\allowbreak{}[C@@H]\allowbreak{}12} & \texttt{CN1c2ccc(\allowbreak{}OC(\allowbreak{}=\allowbreak{}O)\allowbreak{}NCCCCCCCC[NH+]\allowbreak{}3CCOCC3)\allowbreak{}cc2[C@]\allowbreak{}2(\allowbreak{}C)\allowbreak{}CC[NH+]\allowbreak{}(\allowbreak{}C)\allowbreak{}[C@H]\allowbreak{}12} & 7.10 & +2 & fold 2 & cliff \\
CHEMBL65585 & \texttt{CN1CC[C@@]\allowbreak{}2(\allowbreak{}C)\allowbreak{}c3cc(\allowbreak{}OC(\allowbreak{}=\allowbreak{}O)\allowbreak{}NCCCCCCCCCCCCN4CCOCC4)\allowbreak{}ccc3N(\allowbreak{}C)\allowbreak{}[C@@H]\allowbreak{}12} & \texttt{CN1c2ccc(\allowbreak{}OC(\allowbreak{}=\allowbreak{}O)\allowbreak{}NCCCCCCCCCCCC[NH+]\allowbreak{}3CCOCC3)\allowbreak{}cc2[C@]\allowbreak{}2(\allowbreak{}C)\allowbreak{}CC[NH+]\allowbreak{}(\allowbreak{}C)\allowbreak{}[C@H]\allowbreak{}12} & 7.16 & +2 & fold 2 & cliff \\
CHEMBL73593 & \texttt{O=\allowbreak{}C(\allowbreak{}CCC1CCN(\allowbreak{}Cc2ccccc2)\allowbreak{}CC1)\allowbreak{}c1ccc2c(\allowbreak{}c1)\allowbreak{}NCCC2} & \texttt{O=\allowbreak{}C(\allowbreak{}CCC1CC[NH+]\allowbreak{}(\allowbreak{}Cc2ccccc2)\allowbreak{}CC1)\allowbreak{}c1ccc2c(\allowbreak{}c1)\allowbreak{}NCCC2} & 6.79 & +1 & fold 1 & control \\
CHEMBL75013 & \texttt{O=\allowbreak{}C(\allowbreak{}CCC1CCN(\allowbreak{}Cc2ccccc2)\allowbreak{}CC1)\allowbreak{}c1ccc2c(\allowbreak{}c1)\allowbreak{}NCCCC2} & \texttt{O=\allowbreak{}C(\allowbreak{}CCC1CC[NH+]\allowbreak{}(\allowbreak{}Cc2ccccc2)\allowbreak{}CC1)\allowbreak{}c1ccc2c(\allowbreak{}c1)\allowbreak{}NCCCC2} & 6.86 & +1 & fold 1 & control \\
CHEMBL76014 & \texttt{c1ccc2c(\allowbreak{}NCCCCCCCNc3c4c(\allowbreak{}nc5ccccc35)\allowbreak{}CCC4)\allowbreak{}c3c(\allowbreak{}nc2c1)\allowbreak{}CCC3} & \texttt{c1ccc2c(\allowbreak{}NCCCCCCCNc3c4c(\allowbreak{}[nH+]\allowbreak{}c5ccccc35)\allowbreak{}CCC4)\allowbreak{}c3c(\allowbreak{}[nH+]\allowbreak{}c2c1)\allowbreak{}CCC3} & 7.12 & +2 & fold 0 & cliff \\
CHEMBL94546 & \texttt{CN1C(\allowbreak{}=\allowbreak{}O)\allowbreak{}Cc2cc3c(\allowbreak{}CCC4CCN(\allowbreak{}Cc5ccccc5)\allowbreak{}CC4)\allowbreak{}noc3cc21} & \texttt{CN1C(\allowbreak{}=\allowbreak{}O)\allowbreak{}Cc2cc3c(\allowbreak{}CCC4CC[NH+]\allowbreak{}(\allowbreak{}Cc5ccccc5)\allowbreak{}CC4)\allowbreak{}noc3cc21} & 9.32 & +1 & fold 1 & control \\
\end{longtable}

\normalsize

\clearpage

Each descriptor is named here as it appears in the deposited descriptor tables. A
quantity discussed in the main text can therefore be traced to the column it was
computed in. Rungs~2, 2e, 3 and~4 share one panel of fourteen gorge-lining residues and
compute a per-residue quantity for each of them. All three of these tables therefore
carry a block of fourteen entries at the end, followed in the rung-4 table by the
heavy-atom count alone. What that quantity is differs between them. It is
a contact count on one docked pose, the same count taken as a median over eight
receptor conformations, a Coulomb and Lennard-Jones sum over the atoms of the residue,
and the fraction of simulated frames in which the contact was present. The residue
numbering is that of the mature human enzyme as deposited in 4EY7. We read the names
from the receptor used in the calculations rather than assigning them by hand. The
\SI{4.5}{\angstrom} cutoff that defines a contact is the same at rungs~2, 2e and~4, so
a per-residue contact entry means the same geometric thing wherever it appears.

Five of the rung-2 descriptors measure the ligand's position along the axis between
the two anionic sites. They are defined against the single 4EY7 geometry, so the
ensemble carries the remaining 29 and not those five. The rung-4 table carries two
further columns. The first is the spread of a descriptor between ligands divided by its
spread between replicate simulations of one ligand. The second is the verdict that
ratio produced at a threshold of two. A descriptor whose replicate spread was zero has
no ratio, and that case splits in two. A quantity that never varies at all carries no
information, while one that is almost always zero is reproducible only by its own
sparsity. The heavy-atom count appears in the same table because it is carried beside
either feature set as a size covariate. It is not a dynamic quantity and was not put to
the filter.

\begin{table}[htbp]
  \centering
  \footnotesize
  \caption{\textbf{Rung~2 and rung~2e descriptors.} The 33 geometric quantities
    measured on the docked pose and the heavy-atom covariate carried beside them, with
    the 29 of the 34 that the receptor ensemble carries over to rung~2e.}
  \label{tab:sirung2}
  \begin{tabular}{l>{\raggedright\arraybackslash}p{0.54\linewidth}ll}
\toprule
Descriptor & Definition & Unit & Rung~2e \\
\midrule
\addlinespace
\multicolumn{4}{l}{\itshape Engagement of the two anionic sites} \\
\addlinespace[2pt]
\texttt{d\_cas} & Closest approach of a ligand heavy atom to a Trp86 side-chain heavy atom & \si{\angstrom} & yes \\
\texttt{d\_pas} & The same distance to Trp286 & \si{\angstrom} & yes \\
\texttt{d\_gate} & The same distance to Tyr337 & \si{\angstrom} & yes \\
\texttt{bridges} & Both $d_{\mathrm{cas}}$ and $d_{\mathrm{pas}}$ within \SI{4.5}{\angstrom} & 0/1 & yes \\
\texttt{bridge\_gap} & The larger of $d_{\mathrm{cas}}$ and $d_{\mathrm{pas}}$ & \si{\angstrom} & yes \\
\texttt{catpi\_cas} & Closest approach of a formally cationic ligand atom to the centroid of the Trp86 six-membered ring & \si{\angstrom} & yes \\
\texttt{catpi\_pas} & The same distance to the Trp286 ring & \si{\angstrom} & yes \\
\texttt{catpi\_both} & Both cation--$\pi$ distances within \SI{6}{\angstrom} & 0/1 & yes \\
\addlinespace
\multicolumn{4}{l}{\itshape Position along the gorge axis} \\
\addlinespace[2pt]
\texttt{axis\_span\_frac} & Extent of the ligand's projection onto the Trp86$\rightarrow$Trp286 axis, as a fraction of its \SI{15.2}{\angstrom} length & --- & no \\
\texttt{axis\_reach\_pas} & Projection of the ligand atom reaching furthest towards the gorge mouth & --- & no \\
\texttt{axis\_foot\_cas} & Projection of the ligand atom sitting lowest towards the gorge floor & --- & no \\
\texttt{d\_cas\_centroid} & Closest approach of a ligand heavy atom to the Trp86 side-chain centroid & \si{\angstrom} & no \\
\texttt{d\_pas\_centroid} & The same distance to the Trp286 side-chain centroid & \si{\angstrom} & no \\
\addlinespace
\multicolumn{4}{l}{\itshape Burial and contact density} \\
\addlinespace[2pt]
\texttt{n\_contacts\_4A} & Ligand heavy atoms having a protein heavy atom within \SI{4}{\angstrom} & count & yes \\
\texttt{n\_prot\_5A} & Protein heavy atoms within \SI{5}{\angstrom} of any ligand heavy atom & count & yes \\
\texttt{n\_prot\_8A} & Protein heavy atoms within \SI{8}{\angstrom} of any ligand heavy atom & count & yes \\
\texttt{frac\_lig\_buried} & Fraction of ligand heavy atoms having a protein heavy atom within \SI{4.5}{\angstrom} & --- & yes \\
\addlinespace
\multicolumn{4}{l}{\itshape Shape of the bound conformation} \\
\addlinespace[2pt]
\texttt{rgyr} & Radius of gyration of the ligand heavy atoms, unweighted & \si{\angstrom} & yes \\
\texttt{end\_to\_end} & Largest distance between any two ligand heavy atoms & \si{\angstrom} & yes \\
\texttt{n\_heavy} & Ligand heavy atoms, carried as a size covariate & count & yes \\
\addlinespace
\multicolumn{4}{l}{\itshape Per-residue contacts} \\
\addlinespace[2pt]
\texttt{c\_72} & Ligand and residue heavy-atom pairs within \SI{4.5}{\angstrom} of Tyr72 & count & yes \\
\texttt{c\_74} & Ligand and residue heavy-atom pairs within \SI{4.5}{\angstrom} of Asp74 & count & yes \\
\texttt{c\_86} & Ligand and residue heavy-atom pairs within \SI{4.5}{\angstrom} of Trp86 & count & yes \\
\texttt{c\_121} & Ligand and residue heavy-atom pairs within \SI{4.5}{\angstrom} of Gly121 & count & yes \\
\texttt{c\_124} & Ligand and residue heavy-atom pairs within \SI{4.5}{\angstrom} of Tyr124 & count & yes \\
\texttt{c\_203} & Ligand and residue heavy-atom pairs within \SI{4.5}{\angstrom} of Ser203 & count & yes \\
\texttt{c\_286} & Ligand and residue heavy-atom pairs within \SI{4.5}{\angstrom} of Trp286 & count & yes \\
\texttt{c\_295} & Ligand and residue heavy-atom pairs within \SI{4.5}{\angstrom} of Phe295 & count & yes \\
\texttt{c\_297} & Ligand and residue heavy-atom pairs within \SI{4.5}{\angstrom} of Phe297 & count & yes \\
\texttt{c\_334} & Ligand and residue heavy-atom pairs within \SI{4.5}{\angstrom} of Glu334 & count & yes \\
\texttt{c\_337} & Ligand and residue heavy-atom pairs within \SI{4.5}{\angstrom} of Tyr337 & count & yes \\
\texttt{c\_338} & Ligand and residue heavy-atom pairs within \SI{4.5}{\angstrom} of Phe338 & count & yes \\
\texttt{c\_341} & Ligand and residue heavy-atom pairs within \SI{4.5}{\angstrom} of Tyr341 & count & yes \\
\texttt{c\_447} & Ligand and residue heavy-atom pairs within \SI{4.5}{\angstrom} of His447 & count & yes \\
\bottomrule
\end{tabular}

\end{table}

\begin{table}[htbp]
  \centering
  \footnotesize
  \caption{\textbf{Rung~3 descriptors.} The 22 MM/GBSA quantities, computed on the
    rung~2 pose after minimisation of the ligand hydrogens. Differences within a matched
    pair are what the models see, because the absolute values are not interpretable.}
  \label{tab:sirung3}
  \begin{tabular}{l>{\raggedright\arraybackslash}p{0.60\linewidth}l}
\toprule
Descriptor & Definition & Unit \\
\midrule
\addlinespace
\multicolumn{3}{l}{\itshape Interaction energy, receptor to ligand} \\
\addlinespace[2pt]
\texttt{dE\_elec} & Coulomb term between receptor and ligand atoms, on the minimised complex & kcal/mol \\
\texttt{dE\_vdw} & Lennard-Jones term between the same atoms & kcal/mol \\
\texttt{dE\_mm} & Sum of the Coulomb and Lennard-Jones terms & kcal/mol \\
\addlinespace
\multicolumn{3}{l}{\itshape Desolvation and the assembled estimate} \\
\addlinespace[2pt]
\texttt{dG\_mm\_gb} & Energy of the complex less that of the receptor and the ligand apart, each evaluated with the generalised-Born term & kcal/mol \\
\texttt{dG\_polar\_gb} & Polar desolvation, i.e. $\Delta G_{\mathrm{mm,gb}}$ less $\Delta E_{\mathrm{mm}}$ & kcal/mol \\
\texttt{dSASA} & Change in solvent-accessible surface area on binding, over interface atoms & \si{\angstrom\squared} \\
\texttt{dG\_nonpolar} & Nonpolar term, \num{0.0072} times the buried surface area & kcal/mol \\
\texttt{dG\_total} & Sum of $\Delta G_{\mathrm{mm,gb}}$ and the nonpolar term & kcal/mol \\
\addlinespace
\multicolumn{3}{l}{\itshape Per-residue decomposition} \\
\addlinespace[2pt]
\texttt{res\_72} & Coulomb and Lennard-Jones terms summed over the atoms of Tyr72 & kcal/mol \\
\texttt{res\_74} & Coulomb and Lennard-Jones terms summed over the atoms of Asp74 & kcal/mol \\
\texttt{res\_86} & Coulomb and Lennard-Jones terms summed over the atoms of Trp86 & kcal/mol \\
\texttt{res\_121} & Coulomb and Lennard-Jones terms summed over the atoms of Gly121 & kcal/mol \\
\texttt{res\_124} & Coulomb and Lennard-Jones terms summed over the atoms of Tyr124 & kcal/mol \\
\texttt{res\_203} & Coulomb and Lennard-Jones terms summed over the atoms of Ser203 & kcal/mol \\
\texttt{res\_286} & Coulomb and Lennard-Jones terms summed over the atoms of Trp286 & kcal/mol \\
\texttt{res\_295} & Coulomb and Lennard-Jones terms summed over the atoms of Phe295 & kcal/mol \\
\texttt{res\_297} & Coulomb and Lennard-Jones terms summed over the atoms of Phe297 & kcal/mol \\
\texttt{res\_334} & Coulomb and Lennard-Jones terms summed over the atoms of Glu334 & kcal/mol \\
\texttt{res\_337} & Coulomb and Lennard-Jones terms summed over the atoms of Tyr337 & kcal/mol \\
\texttt{res\_338} & Coulomb and Lennard-Jones terms summed over the atoms of Phe338 & kcal/mol \\
\texttt{res\_341} & Coulomb and Lennard-Jones terms summed over the atoms of Tyr341 & kcal/mol \\
\texttt{res\_447} & Coulomb and Lennard-Jones terms summed over the atoms of His447 & kcal/mol \\
\bottomrule
\end{tabular}

\end{table}

\begin{table}[htbp]
  \centering
  \footnotesize
  \caption{\textbf{Rung~4 descriptors and the reliability filter.} The 32 dynamic
    descriptors, the ratio of between-ligand to between-replicate spread for each, and
    the verdict that ratio produced.}
  \label{tab:sirung4}
  \begin{tabular}{l>{\raggedright\arraybackslash}p{0.40\linewidth}rl}
\toprule
Descriptor & Definition & $s_{\mathrm{lig}}/s_{\mathrm{rep}}$ & Verdict \\
\midrule
\addlinespace
\multicolumn{4}{l}{\itshape Persistence of the bridging contact} \\
\addlinespace[2pt]
\texttt{occ\_cas} & Fraction of frames with a ligand heavy atom within \SI{4.5}{\angstrom} of Trp86 & 103.10 & usable \\
\texttt{occ\_pas} & The same fraction for Trp286 & 51.54 & usable \\
\texttt{occ\_bridge} & Fraction of frames contacting Trp86 and Trp286 at once & 55.65 & usable \\
\texttt{bridge\_runs\_per\_ns} & Separate bridging episodes begun per nanosecond & 4.22 & usable \\
\texttt{bridge\_longest\_ns} & Longest uninterrupted bridging episode & 1.31 & seed-dominated \\
\addlinespace
\multicolumn{4}{l}{\itshape Whether the ligand stayed where it was docked} \\
\addlinespace[2pt]
\texttt{rmsd\_mean} & Ligand heavy-atom RMSD from the equilibrated pose, averaged over frames, measured after superposing the receptor on its reference by C$\alpha$ & 0.88 & seed-dominated \\
\texttt{rmsd\_max} & Largest value that RMSD reached & 1.32 & seed-dominated \\
\texttt{escaped} & Whether the ligand centre of mass ever left a \SI{12}{\angstrom} sphere about its starting position & --- & constant \\
\texttt{depth\_mean} & Position of the ligand centroid along the Trp86$\rightarrow$Trp286 axis, which is recomputed each frame & 3.08 & usable \\
\texttt{depth\_sd} & Standard deviation of that position & 1.29 & seed-dominated \\
\addlinespace
\multicolumn{4}{l}{\itshape Flexibility, which a single pose has none of} \\
\addlinespace[2pt]
\texttt{rmsf\_lig} & Per-atom fluctuation of the ligand, averaged over its heavy atoms & 1.38 & seed-dominated \\
\texttt{rmsf\_gorge} & The same fluctuation for the 14 gorge-lining side chains & 1.09 & seed-dominated \\
\texttt{gate\_chi1\_sd} & Circular standard deviation of the Tyr337 $\chi_1$ dihedral & 0.66 & seed-dominated \\
\texttt{gate\_flips\_per\_ns} & Crossings of Tyr337 $\chi_1$ between rotamer wells per nanosecond & 0.41 & seed-dominated \\
\texttt{lig\_dihedral\_sd} & Circular standard deviation of the ligand's own rotatable dihedrals, averaged over them & 1.78 & seed-dominated \\
\addlinespace
\multicolumn{4}{l}{\itshape Burial, water and quality control} \\
\addlinespace[2pt]
\texttt{frac\_lig\_buried\_mean} & Fraction of ligand heavy atoms within \SI{4.5}{\angstrom} of the protein, averaged over frames & 1.50 & seed-dominated \\
\texttt{n\_bridge\_water} & Waters within \SI{3.5}{\angstrom} of both a ligand and a protein polar heavy atom, averaged over frames & 1.68 & seed-dominated \\
\texttt{ca\_rmsd\_mean} & Receptor C$\alpha$ RMSD after superposition, kept as a check that the receptor was reassembled whole across the periodic boundary & 1.00 & seed-dominated \\
\addlinespace
\multicolumn{4}{l}{\itshape Per-residue contact persistence} \\
\addlinespace[2pt]
\texttt{occ\_72} & Fraction of frames within \SI{4.5}{\angstrom} of Tyr72 & 4.12 & usable \\
\texttt{occ\_74} & Fraction of frames within \SI{4.5}{\angstrom} of Asp74 & 1.07 & seed-dominated \\
\texttt{occ\_86} & Fraction of frames within \SI{4.5}{\angstrom} of Trp86 & 1241.70 & usable \\
\texttt{occ\_121} & Fraction of frames within \SI{4.5}{\angstrom} of Gly121 & 3.48 & usable \\
\texttt{occ\_124} & Fraction of frames within \SI{4.5}{\angstrom} of Tyr124 & 3.69 & usable \\
\texttt{occ\_203} & Fraction of frames within \SI{4.5}{\angstrom} of Ser203 & 0.99 & seed-dominated \\
\texttt{occ\_286} & Fraction of frames within \SI{4.5}{\angstrom} of Trp286 & 51.51 & usable \\
\texttt{occ\_295} & Fraction of frames within \SI{4.5}{\angstrom} of Phe295 & 1.66 & seed-dominated \\
\texttt{occ\_297} & Fraction of frames within \SI{4.5}{\angstrom} of Phe297 & 0.83 & seed-dominated \\
\texttt{occ\_334} & Fraction of frames within \SI{4.5}{\angstrom} of Glu334 & --- & sparse \\
\texttt{occ\_337} & Fraction of frames within \SI{4.5}{\angstrom} of Tyr337 & 3.11 & usable \\
\texttt{occ\_338} & Fraction of frames within \SI{4.5}{\angstrom} of Phe338 & 0.77 & seed-dominated \\
\texttt{occ\_341} & Fraction of frames within \SI{4.5}{\angstrom} of Tyr341 & 0.82 & seed-dominated \\
\texttt{occ\_447} & Fraction of frames within \SI{4.5}{\angstrom} of His447 & 0.98 & seed-dominated \\
\addlinespace
\multicolumn{4}{l}{\itshape Size covariate} \\
\addlinespace[2pt]
\texttt{n\_heavy} & Ligand heavy atoms, which is not a dynamic quantity and is not put to the filter & --- & not applicable \\
\bottomrule
\end{tabular}

\end{table}

\clearpage

Tables~\ref{tab:benchreg} and~\ref{tab:benchclf} hold the full grid of scores that the
ladder summarises, with every rung crossed against every descriptor arm and every
algorithm family. The headline comparison takes the best value a rung reached over its
arms and its families, which is one number drawn from a block of them, and these two
tables are that block. They are split because the regression and the classification framings
do not share a model suite. Partial least squares has no classification counterpart
and is absent from the second rather than having failed in it, while the linear model of
the regression suite appears there as a logistic regression. Reading down a column shows how
much of a rung's score is the descriptors and how much is the choice of algorithm,
which the single best value conceals. The arms are the alternative feature sets of the
main text. These are the two fingerprint variants of rung~1 and the filtered and
unfiltered sets of rung~4, while rungs~2, 2e and~3 carry one arm each.

\begin{table}[tb]
  \centering
  \caption{\textbf{Every regression family on \framingA{} and \framingB{}.} Each rung
    with every descriptor arm it was scored with, so a quoted best value can be
    located in its grid.}
  \label{tab:benchreg}
  \begin{tabular}{lllrr}
\toprule
Rung & Descriptor arm & Algorithm & \framingA{} $R^2$ & \framingB{} $R^2$ \\
\midrule
Rung 1 2D & binary ECFP4 & Ridge (linear) & +0.175 & +0.106 \\
 &  & PLS & +0.386 & +0.100 \\
 &  & k-NN (k=5) & +0.104 & +0.099 \\
 &  & SVM (RBF) & +0.383 & +0.124 \\
 &  & Random Forest & +0.329 & +0.184 \\
 &  & Gradient Boosting & +0.375 & +0.200 \\
\addlinespace
Rung 1 2D & count ECFP4 + RDKit2D & Ridge (linear) & +0.077 & -513.621 \\
 &  & PLS & +0.357 & -687.740 \\
 &  & k-NN (k=5) & +0.081 & -0.637 \\
 &  & SVM (RBF) & +0.380 & -0.519 \\
 &  & Random Forest & +0.359 & -0.328 \\
 &  & Gradient Boosting & +0.336 & -0.007 \\
\addlinespace
Rung 2 pose & 4EY7 + bridging & Ridge (linear) & +0.045 & -0.186 \\
 &  & PLS & +0.046 & -0.120 \\
 &  & k-NN (k=5) & -0.074 & -0.075 \\
 &  & SVM (RBF) & -0.213 & +0.167 \\
 &  & Random Forest & +0.031 & +0.063 \\
 &  & Gradient Boosting & -0.058 & +0.024 \\
\addlinespace
Rung 2 ensemble & ensemble median & Ridge (linear) & +0.183 & -0.029 \\
 &  & PLS & +0.131 & -0.018 \\
 &  & k-NN (k=5) & +0.155 & +0.134 \\
 &  & SVM (RBF) & +0.201 & +0.274 \\
 &  & Random Forest & +0.215 & +0.043 \\
 &  & Gradient Boosting & +0.182 & +0.190 \\
\addlinespace
Rung 3 MM/GBSA & MM/GBSA & Ridge (linear) & +0.034 & -0.183 \\
 &  & PLS & +0.026 & -0.259 \\
 &  & k-NN (k=5) & -0.045 & -0.239 \\
 &  & SVM (RBF) & -0.045 & -0.228 \\
 &  & Random Forest & +0.118 & -0.056 \\
 &  & Gradient Boosting & +0.041 & -0.110 \\
\addlinespace
Rung 4 MD & replicate-filtered (12) & Ridge (linear) & +0.013 & +0.084 \\
 &  & PLS & -0.028 & +0.072 \\
 &  & k-NN (k=5) & -0.153 & +0.012 \\
 &  & SVM (RBF) & -0.067 & +0.124 \\
 &  & Random Forest & +0.056 & -0.084 \\
 &  & Gradient Boosting & -0.028 & -0.104 \\
\addlinespace
Rung 4 MD & all descriptors (33) & Ridge (linear) & -0.186 & -0.148 \\
 &  & PLS & -0.279 & -0.054 \\
 &  & k-NN (k=5) & +0.014 & -0.060 \\
 &  & SVM (RBF) & -0.307 & -0.108 \\
 &  & Random Forest & +0.068 & -0.073 \\
 &  & Gradient Boosting & -0.110 & -0.166 \\
\addlinespace
\bottomrule
\end{tabular}

\end{table}

\begin{table}[tb]
  \centering
  \caption{\textbf{Every classification family on \framingC{}.} The families differ
    from Table~\ref{tab:benchreg} because \framingC{} is a classification; partial least
    squares has no counterpart here and is absent rather than failed.}
  \label{tab:benchclf}
  \begin{tabular}{lllr}
\toprule
Rung & Descriptor arm & Algorithm & \framingC{} ROC AUC \\
\midrule
Rung 1 2D & binary ECFP4 & Logistic regression & 0.597 \\
 &  & k-NN (k=5) & 0.590 \\
 &  & SVM (RBF) & 0.657 \\
 &  & Random Forest & 0.648 \\
 &  & Gradient Boosting & 0.520 \\
\addlinespace
Rung 1 2D & count ECFP4 + RDKit2D & Logistic regression & 0.609 \\
 &  & k-NN (k=5) & 0.683 \\
 &  & SVM (RBF) & 0.606 \\
 &  & Random Forest & 0.805 \\
 &  & Gradient Boosting & 0.575 \\
\addlinespace
Rung 2 pose & 4EY7 + bridging & Logistic regression & 0.574 \\
 &  & k-NN (k=5) & 0.614 \\
 &  & SVM (RBF) & 0.555 \\
 &  & Random Forest & 0.602 \\
 &  & Gradient Boosting & 0.684 \\
\addlinespace
Rung 2 ensemble & ensemble median & Logistic regression & 0.639 \\
 &  & k-NN (k=5) & 0.681 \\
 &  & SVM (RBF) & 0.635 \\
 &  & Random Forest & 0.705 \\
 &  & Gradient Boosting & 0.673 \\
\addlinespace
Rung 3 MM/GBSA & MM/GBSA & Logistic regression & 0.689 \\
 &  & k-NN (k=5) & 0.698 \\
 &  & SVM (RBF) & 0.630 \\
 &  & Random Forest & 0.753 \\
 &  & Gradient Boosting & 0.655 \\
\addlinespace
Rung 4 MD & replicate-filtered (12) & Logistic regression & 0.749 \\
 &  & k-NN (k=5) & 0.737 \\
 &  & SVM (RBF) & 0.741 \\
 &  & Random Forest & 0.793 \\
 &  & Gradient Boosting & 0.785 \\
\addlinespace
Rung 4 MD & all descriptors (33) & Logistic regression & 0.627 \\
 &  & k-NN (k=5) & 0.587 \\
 &  & SVM (RBF) & 0.535 \\
 &  & Random Forest & 0.721 \\
 &  & Gradient Boosting & 0.700 \\
\addlinespace
\bottomrule
\end{tabular}

\end{table}

\end{document}
